\documentclass[10pt,a4paper]{article}

% ----------------- Paquetes -------------------%
\usepackage[T1]{fontenc}
\usepackage[utf8]{inputenc}
\usepackage[a4paper,margin=2.5cm]{geometry}
\usepackage{mathptmx}             
\usepackage{changepage} 
\usepackage{setspace}
\usepackage[spanish]{babel}
\usepackage[labelfont=bf,labelsep=space]{caption}
\usepackage{amssymb}
\usepackage{amsmath}
\usepackage{ebgaramond}
\usepackage{placeins}
\usepackage{etoolbox}
\usepackage{setspace}
\usepackage{parskip} 
\usepackage{tcolorbox}
\usepackage{needspace}
\usepackage{xparse}
\usepackage{ocgx2}
\usepackage{hyperref}
\usepackage{hypcap}


%%%%%%%%%%%%%%%%%%%%%%%%%%%%%%%%%%%%%%%%%%%%%%%%%%%%%%%%%%%%%%%%
% --- Fija primero tus valores globales "buenos" ---
\setstretch{1.2}                 % Interlineado global
\setlength{\parskip}{0.7\baselineskip}
\setlength{\parindent}{0pt}
\newlength{\GlobalParskip}
\setlength{\GlobalParskip}{\parskip}

\newcounter{eqnum}[subsection]
\renewcommand{\theeqnum}{\arabic{eqnum}}
\newcounter{alphaitem}
\newcounter{numitem}

%% ------------------- Recuadros----------------------%%
\newcommand{\puntosclave}[1]{%
  \begin{tcolorbox}[
    colframe=black,
    title={Puntos clave},
    before upper={\setlength{\parindent}{0.5cm}\setlength{\parskip}{0.5\baselineskip}}
  ]
  #1
  \end{tcolorbox}
}

\newcommand{\modificaciones}[1]{
  \begin{tcolorbox}[
    colback=white,
    colframe=red,
    title={Modificaciones a realizar},
    before upper={\setlength{\parindent}{0.5cm}\setlength{\parskip}{0.5\baselineskip}}
  ]
  #1
  \end{tcolorbox}
}

%% ------------------- Listas----------------------%%
    % --- Lista alfabética ---
    \newenvironment{la}
      {\begin{list}{\alph{alphaitem})}{%
          \usecounter{alphaitem}%
          \setlength{\leftmargin}{1cm}%
          \setlength{\parskip}{3pt}% 
          \setlength{\itemsep}{0pt}% ← espacio entre ítems
          \setlength{\parsep}{0pt}%  ← párrafos dentro de un ítem
      }}
      {\end{list}}
      
    % --- Lista numérica ---
    \newenvironment{lnum}
      {\begin{list}{\arabic{numitem}.}{%
          \usecounter{numitem}%
          \setlength{\leftmargin}{1cm}%
          \setlength{\parskip}{3pt}% 
          \setlength{\itemsep}{0pt}% ← espacio entre ítems
          \setlength{\parsep}{0pt}%  ← párrafos dentro de un ítem
      }}
      {\end{list}}
      
% ------------------- Imágenes----------------------%
    \graphicspath{{Img/}}% Rutas por defecto 
    % --- Imagen adaptativa ---
    % Registros de dimensión definidos una sola vez
    \newdimen\imgcap
    \newdimen\imgmin
    \newdimen\imgmax
    \newdimen\imgremain
    \newdimen\imgh
    \newdimen\imgneed
    
    \newcommand{\imagenfit}[3]{%
      \addvspace{10pt}% separación antes
      \par\begingroup
        % Parámetros de composición (asignaciones, no \newdimen)
        \imgcap=1\baselineskip            % alto estimado de título+fuente
        \imgmin=0.15\textheight
        \imgmax=0.15\textheight
        \imgremain=\pagegoal
        \ifdim\pagegoal=\maxdimen \imgremain=0pt\fi
        \advance\imgremain by -\pagetotal
        \advance\imgremain by -\imgcap
        % Decisión de altura destino
        \ifdim\imgremain<\imgmin
          \newpage
          \imgh=0.20\textheight
        \else
          \imgh=\imgremain
          \ifdim\imgh>\imgmax \imgh=\imgmax \fi
        \fi
        % Reservar espacio para evitar cortes del bloque completo
        \imgneed=\imgh \advance\imgneed by \imgcap
        \Needspace{\imgneed}
        % Cuerpo
        \noindent\begin{minipage}{\textwidth}
          \centering
          \captionof{figure}{\textemdash\ #2}\par\vspace{0.3em}% título arriba
          \includegraphics[width=\linewidth,height=\imgh,keepaspectratio]{\detokenize{#1}}\par
          \vspace{0.3em}\small\textit{Fuente: }#3% fuente abajo
        \end{minipage}%
      \endgroup
      \par\addvspace{10pt}%
    }

%------------ Bloque de RECUADRO ESPECIAL -------------%
    \newenvironment{specialblock}[1]{%
      \begin{quote} % crea sangría a izquierda y derecha
      \small % texto más pequeño
      \noindent\textbf{#1}\\[-0.3em] % título en negrita
      \rule{\linewidth}{0.4pt}\\[0.6em]}{
      \end{quote}}
      

%-------------------------- Author Font -----------------------%
    \newcommand{\authorfont}[1]{{\fontfamily{EBGaramond-TLF}\selectfont #1}}

%-------------------------- Anexos -----------------------%
    \makeatletter
    \newcounter{anexo}
    \renewcommand{\theanexo}{\Roman{anexo}}
    \newcommand{\anexo}[1]{%
      \clearpage
      \phantomsection
      \refstepcounter{anexo}%
      \noindent\textbf{Anexo \theanexo.- }#1\par\medskip
      \addcontentsline{stc}{section}{Anexo \theanexo.- #1}%
    }
    \makeatother

    %-------------------------- Lista de figuras (personal) -----------------------%
\makeatletter
\newcommand{\MyListOfFigures}{%
  \clearpage
  \phantomsection
  {\centering\bfseries\Large Índice de figuras\par}%
  \medskip
  % Entrada en nuestro índice de contenido (stc)
  \addcontentsline{stc}{section}{Índice de figuras}%
  % Recuperación del listado (sin cabecera propia)
  \begingroup
    \parindent=0pt\parskip=2pt
    \@starttoc{lof}%
  \endgroup
}
\makeatother

%-------------------------- Bibliografía -----------------------%
\makeatletter
% Contador para las referencias
\newcounter{myref}
% Cabecera de la bibliografía, entra en el índice 'stc'
\newcommand{\MyBibliography}{%
  \clearpage
  \phantomsection
  {\centering\bfseries\Large Bibliografía y referencias\par}%
  \medskip
  \addcontentsline{stc}{section}{Bibliografía y referencias}%
  \setcounter{myref}{0}%
}
\newcommand{\MyBibItem}[4]{%
  \refstepcounter{myref}%
  \noindent[\themyref]\ #1.\ \emph{#2}. #3. #4\par
}
\makeatother

%--------- Establecer cuatro niveles de índice --------%
    \setcounter{tocdepth}{4}   
    \setcounter{secnumdepth}{4}% numera hasta \paragraph

%%%%%%%%%%%%%%%%%%%%%%%%%%%%%%%%

\makeatletter

    %--------- Interlineado Global --------%
    \edef\GlobalBaseLineStretch{\baselinestretch}
    
    %--------- Espaciado de seccinones --------%
    \renewcommand\section{\@startsection{section}
      {1}{\z@}
      {2.5ex \@plus 1ex \@minus .2ex} % espacio antes
      {1.5ex \@plus .2ex}             % espacio después
      {\normalfont\Large\bfseries}    % formato del título
    }
    \renewcommand\subsection{\@startsection{subsection}{2}{\z@}
      {3ex \@plus 0.8ex \@minus .2ex}
      {1.5ex \@plus .2ex}
      {\normalfont\large\bfseries}
    }
    \renewcommand\subsubsection{\@startsection{subsubsection}{3}{\z@}
      {3ex \@plus 0.5ex \@minus .2ex}
      {1ex \@plus .2ex}
      {\normalfont\normalsize\bfseries}
    }
    \renewcommand\paragraph{\@startsection{paragraph}{4}{\z@}%
    {-2ex \@plus -0.5ex \@minus -.2ex}% espacio antes
    {0.8ex \@plus .2ex}% espacio después
    {\normalfont\normalsize\itshape}}% ← cursiva en lugar de negrita

    %--------- Ecuaciones --------%   
    % Variables globales para almacenar títulos y notas
    \gdef\leq@current@section{}%
    \gdef\leq@current@subsection{}%
    \gdef\leq@last@printed@section{}%
    \gdef\leq@last@printed@subsection{}%
    
    % Comando para añadir nota (invisible en el cuerpo)
    \newcommand{\eqnote}[1]{%
      \addcontentsline{leq}{leqnote}{#1}%
    }
    
    % Comando para ecuaciones
    \newcommand{\eqblock}[2]{%
      % Verificar si necesitamos imprimir el título de sección
      \ifx\leq@current@section\leq@last@printed@section\else
        \addcontentsline{leq}{leqsection}{\leq@current@section}%
        \global\let\leq@last@printed@section\leq@current@section
        \gdef\leq@last@printed@subsection{}% Reset subsección al cambiar sección
      \fi
      % Verificar si necesitamos imprimir el título de subsección
      \ifx\leq@current@subsection\leq@last@printed@subsection\else
        \ifx\leq@current@subsection\@empty\else
          \addcontentsline{leq}{leqsubsection}{\leq@current@subsection}%
          \global\let\leq@last@printed@subsection\leq@current@subsection
        \fi
      \fi
      % Ecuación
      \refstepcounter{eqnum}%
      \begin{equation*}
        #1\tag{E.\theeqnum}
      \end{equation*}%
      \addcontentsline{leq}{leqt}{[E.\theeqnum] -- #2}%
      \addcontentsline{leq}{leqm}{#1}%
    }
    
    %%% ----- Índice de ecuaciones ----------%%%
    % Tamaño para []
    \newcommand{\leqIndexFont}{\normalsize}
    \newcommand{\leqTitleFont}{\tiny}
    \newcommand{\leqNoteFont}{\tiny}
    % Cabecera de sección 
    \newcommand*\l@leqsection[2]{%
      \addvspace{25pt}% Mayor separación antes de nueva sección
      \noindent\leqIndexFont
      \makebox[\linewidth]{\rule{.38\linewidth}{0.4pt}\ \ \textbf{  \MakeUppercase{#1}}\ \ \rule{.38\linewidth}{0.4pt}}%
      \par\vspace{-10pt}
    }
    % Cabecera de subsección 
    \newcommand*\l@leqsubsection[2]{%
      \par\addvspace{15pt}%
      \noindent\leqIndexFont #1
      \par\vspace{-7pt}
    }
    % Título de ecuación 
    \newcommand*\l@leqt[2]{%
    \par\addvspace{10pt}%
      \par\noindent\leqTitleFont #1\par
    }
    % Miniatura de ecuación
    \newcommand*\l@leqm[2]{%
      \vspace{0.3\baselineskip}%
      \par\scriptsize\noindent\hspace*{1.5em}\(#1\)\par%
    }
    % Nota de ecuación 
    \newcommand*\l@leqnote[2]{%
      \vspace{0.2\baselineskip}%
      \par\noindent\leqNoteFont\hspace*{1.5em}\textbf{Nota.--} #1\par\vspace{0.5\baselineskip}%
    }
    % Generación del índice
    \newcommand{\mylistofequations}{%
      \clearpage
      \phantomsection
      % Título visual de la lista (ajústalo a tu gusto):
      {\centering\bfseries\Large Índice de ecuaciones\par}
      % Entrada en nuestro índice 'stc':
      \addcontentsline{stc}{section}{Índice de ecuaciones}%
      % Llamada a tu lista real (usa el nombre exacto de tu macro):
      \@starttoc{leq}%
    }
    
    %--------- Captura de títulos de sección --------%
    \let\leq@original@section\section
    \let\leq@original@subsection\subsection
    % Flag para saber si estamos en una sección asterisco
    \newif\ifLeqStarSection
    \LeqStarSectionfalse
    
    % Redefinir \section CON CONTROL DE ASTERISCO
    \renewcommand{\section}{%
      \@ifstar{\leq@section@star}{\@dblarg\leq@section@nostar}%
    }
    \def\leq@section@star#1{%
      \global\LeqStarSectiontrue
      \gdef\leq@current@section{#1}%
      \gdef\leq@current@subsection{}%
      \leq@original@section*{#1}%
      \global\LeqStarSectionfalse
    }
    \def\leq@section@nostar[#1]#2{%
      \global\LeqStarSectionfalse
      \gdef\leq@current@section{#2}%
      \gdef\leq@current@subsection{}%
      \leq@original@section[#1]{#2}%
    }
    % Redefinir \subsection
    \renewcommand{\subsection}{%
      \@ifstar{\leq@subsection@star}{\@dblarg\leq@subsection@nostar}%
    }
    \def\leq@subsection@star#1{%
      \global\LeqStarSectiontrue
      \gdef\leq@current@subsection{#1}%
      \leq@original@subsection*{#1}%
      \global\LeqStarSectionfalse
    }
    \def\leq@subsection@nostar[#1]#2{%
      \global\LeqStarSectionfalse
      \gdef\leq@current@subsection{#2}%
      \leq@original@subsection[#1]{#2}%
    }

    % ----- Restauración de valores Globales de estilo ----------%
    % Flags
    \newif\ifInFootnote
    \InFootnotefalse
    \pretocmd{\@footnotetext}{\InFootnotetrue}{}{}
    \apptocmd{\@footnotetext}{\InFootnotefalse}{}{}
    % Profundidad de listas (soporta anidadas)
    \newcount\MyListDepth \MyListDepth=0
    \AtBeginEnvironment{la}{\advance\MyListDepth by 1}
    \AfterEndEnvironment{la}{\advance\MyListDepth by -1}
    \AtBeginEnvironment{lnum}{\advance\MyListDepth by 1}
    \AfterEndEnvironment{lnum}{\advance\MyListDepth by -1}
    \AtBeginEnvironment{itemize}{\advance\MyListDepth by 1}
    \AfterEndEnvironment{itemize}{\advance\MyListDepth by -1}
    \AtBeginEnvironment{enumerate}{\advance\MyListDepth by 1}
    \AfterEndEnvironment{enumerate}{\advance\MyListDepth by -1}
    % Restauración diferida: hazlo al INICIO del próximo párrafo real
    \newif\if@pendingRestore
    \@pendingRestorefalse
    \newcommand{\RequestRestore}{\global\@pendingRestoretrue}
    \everypar{%
      \if@pendingRestore
        \global\@pendingRestorefalse
        \ifInFootnote\else
          \ifnum\MyListDepth=0
            \setlength{\parskip}{\GlobalParskip}%
            \setstretch{\GlobalBaseLineStretch}%
          \fi
        \fi
      \fi
    }
    % Al final de bloques:
    \AfterEndEnvironment{equation}{\RequestRestore}
    \AfterEndEnvironment{equation*}{\RequestRestore}
    \AfterEndEnvironment{la}{\RequestRestore}
    \AfterEndEnvironment{lnum}{\RequestRestore}
    % Al final de macros:
    \apptocmd{\eqblock}{\RequestRestore}{}{}
    \apptocmd{\imagenfit}{\RequestRestore}{}{}
    
\makeatother

% ===================== ToC propio con 4 niveles (único bloque) =====================
\makeatletter

% --- Escritores: añadimos una línea al .stc en cada cabecera numerada ---
% Guardamos las versiones actuales para llamar al comportamiento normal
\let\MyTOC@orig@section\section
\let\MyTOC@orig@subsection\subsection
\let\MyTOC@orig@subsubsection\subsubsection
\let\MyTOC@orig@paragraph\paragraph

% SECTION
\renewcommand{\section}{%
  \@ifstar{\MyTOC@section@star}{\@dblarg\MyTOC@section@nostar}%
}
\def\MyTOC@section@star#1{%
  \MyTOC@orig@section*{#1}% sin entrada al .stc
}
\def\MyTOC@section@nostar[#1]#2{%
  \MyTOC@orig@section[{#1}]{#2}% comportamiento normal (escribe al .toc)
  % Escribimos también nuestra línea al .stc (texto con \numberline)
  \addcontentsline{stc}{section}{\protect\numberline{\thesection}#2}%
}

% SUBSECTION
\renewcommand{\subsection}{%
  \@ifstar{\MyTOC@subsection@star}{\@dblarg\MyTOC@subsection@nostar}%
}
\def\MyTOC@subsection@star#1{%
  \MyTOC@orig@subsection*{#1}%
}
\def\MyTOC@subsection@nostar[#1]#2{%
  \MyTOC@orig@subsection[{#1}]{#2}%
  \addcontentsline{stc}{subsection}{\protect\numberline{\thesubsection}#2}%
}

% SUBSUBSECTION
\renewcommand{\subsubsection}{%
  \@ifstar{\MyTOC@subsubsection@star}{\@dblarg\MyTOC@subsubsection@nostar}%
}
\def\MyTOC@subsubsection@star#1{%
  \MyTOC@orig@subsubsection*{#1}%
}
\def\MyTOC@subsubsection@nostar[#1]#2{%
  \MyTOC@orig@subsubsection[{#1}]{#2}%
  \addcontentsline{stc}{subsubsection}{\protect\numberline{\thesubsubsection}#2}%
}

% PARAGRAPH
\renewcommand{\paragraph}{%
  \@ifstar{\MyTOC@paragraph@star}{\@dblarg\MyTOC@paragraph@nostar}%
}
\def\MyTOC@paragraph@star#1{%
  \MyTOC@orig@paragraph*{#1}%
}
\def\MyTOC@paragraph@nostar[#1]#2{%
  \MyTOC@orig@paragraph[{#1}]{#2}%
  % Si no numeras los \paragraph, cambia \theparagraph por texto plano sin \numberline
  \addcontentsline{stc}{paragraph}{\protect\numberline{\theparagraph}#2}%
}

% --- Impresor del índice propio (lee el .stc) ---
\newcommand{\MyTOC}{%
  \par\bigskip
  {\centering\bfseries\Large Índice de contenido\par}%
  \medskip
  \begingroup
    \parindent=0pt\parskip=2pt
    \@starttoc{stc}%
  \endgroup
}

\makeatother
% ===================== fin ToC propio =====================


%%%%%%%%%%%%%%


% --- Datos del tema ---
\title{Empresa Pública\\
\large Razones de su existencia. Política de provatización. Comparaciones internacionales}
\author{Marcelo Ortega}
\date{Septiembre 2026}

\usepackage{soul}\sethlcolor{yellow}% resaltado amarillo: \hl{texto} (Panel TCEE, ⌃H)
\begin{document}
\maketitle
\newpage

% --- Página de resumen y modificaciones ---
\puntosclave{ %% Escribir puntos clave Apuntes CECO
     
    }
\vspace{1cm}

\modificaciones{
    
    }

\newpage
\MyTOC
\newpage

\mylistofequations
\clearpage

\MyListOfFigures
\clearpage


\pagenumbering{arabic}
\setcounter{page}{1}


\section{Introducción}



\subsection{Contextualización}


\subsubsection{Relevancia}




\subsubsection{Problemática}



\subsection{Estructura de la exposición}










\clearpage
\section{Empresa Pública: Marco conceptual}
\puntosclave{
    La empresa pública se puede definir como una entidad jurídica sujeta a control público que opera en el mercado. 
    
    A lo largo de la historia han aparecido distintos argumentos para justificar la existencia de la empresa pública, siendo el principal la existencia de fallos de mercado que impiden la competencia, como el monopolio natural.
    
    Además, la existencia de la empresa pública puede justificarse también por motivos de equidad, para garantizar la seguridad de los suministros de ciertos bienes y servicios o aumentar la capacidad del estado de suavizar el ciclo económico.

    Existen una serie de argumentos teóricos que entienden que la empresa pública va ser relativamente más ineficiente que la privada. Estos argumentos se ligarían, a grandes rasgos, a la doble relación de agencia que existe en su seno: dirección-control y control-propiedad. Además, habitualmente se enfrentan a mayores restricciones jurídicas y burocráticas, por lo que esta asimetría podría incluso verse agravada.
}

Este tema se integra en el bloque de Economía Pública de forma relativamente aislada, en el sentido de que es el único en el que se realiza un análisis teórico de la empresa pública de forma individualizada.

En primera parte del- tema, referida al concepto de empresa puhlic:a y a la justificación de su exist,mcia, es importante realizar una buena definición teórica de lo que es la empresa pública (para delimitar el objeto de estudio del tema) y, te-niendo en cuenta que muchos de los fallos de mercado que justifican la existencia de la empresa pública ya se analizan en el temario del tercer ejercicio. el opositor debeni centrarse en relacionarlos con la empresa pública. poniendo ejemplos prácticos v evitando dedicar excesivo tiempo a explicaciones teóricas que sean objeto de otros temas.

En el segundo apartado de política de privatización, se analizan las limitaciones de la empresa pública desde un punto de vista teórico (en este caso también, se mencionan algunos fallos del sector público que son objeto del temario del tercer ejercicio, de modo que deberán explicarse de forma sintética y evitando dedicar excesivo tiempo a conceptos que fom1en parte de otros temas) para posteriorme111e pasar a explicar cómo diseí'íar un proceso de privatización. Puesto que el proceso de privatización se analiza exclusivamente en este tema. puede realizarse una explicación más exhaustiva, dejando claro cuáles son los trade-o.f.f.s· a los que se enfrenta el ��ector público a la hora de privatizar, así como las alternativas que existen a lo largo del proceso.

En tercer lugar, el apartado de comparaciones i111ernacio11ales pretende ofrecer una visión global del rol que las empresas públicas juegan en el mundo y su evolución más reciente. Va a ser importante formarse una visión de conjunto y poder transmitir los principales retos que existen a nivel regulatorio. asociados, por ejemplo, a las distorsiones a la competencia a nivel internacional que pueden generar las empresas públicas.

Finalmente, el apartado de evidencia empírica pennite dotar de mayor fo1malidad al conjunto del tema, de modo que conviene dedicarle tiempo para explicar con rigor los esmdios seleccionados.

El opositor debería ser capaz de entender y plantear los principales retos metodológicos que plantea el análisis de la eficiencia de la empresa públíca, las soluciones que se han encontrado a los mismos en la evidencia empírica seleccionada y las principales conclusiones alcanzadas.

De cara a actualizar v profundizar en el estudio del tema. los dos primeros apartados, de justificación de la empresa pública y limitaciones que ésta presenta pueden abordarse con cualquier manual de economía pública. a fin de entender los principales argumentos planteados.

En cambio, el apartado III.lll. de diseí'io de la estrategia de privatización, así como el de comparaciones internacionales, se han elaborado en base a publicaciones de distintas instituciones internacionales como la OCDE o el FM 1, que pueden revisarse y actualizar. Finalmente, en el apartado de evidencia empírica se han seleccionado dos artículos en los que se puede profundizar. sin perjuicio de intToducir otros si el opositor lo considerara conveniente.

\textcolor{red}{\textbf{INTRODUCCIÓN}}
La economía pública se encarga de analizar la forma en que la intervención del Estado en la economía afecta a la asignación final en términos en términos de bienestar social. Existen tres principales formas de intervención: la primera es a través del presupuesto público, en la que el gobierno utiliza sus recursos para financiar programas y proyectos que buscan mejorar la calidad de vida de los ciudadanos y el desarrollo económico del país. La segunda forma de intervención es mediante la regulación, en la que se establecen normas y leyes que regulan el comportamiento de las empresas y los mercados. Y, por último, la provisión directa que implica la propiedad y control gubernamental de empresas que prestan servicios o proveen bienes al mercado.

Desde un punto de vista histórico la de empresa pública se ha venido justificando únicamente por razones de eficiencia, normalmente debido a la presencia de fallos de mercado, como los monopolios naturales o las industrias red. No obstante, también han sido relevantes algunos argumentas que versan sobre la equidad. 

Sin embargo, en los años 80 cobran mayor relevancia las críticas a la empresa pública ligadas a su potencial ineficiencia relativa, que sirvieron para justificar sendas políticas de privatización a finales del siglo XX.

En cualquier caso. y a pesar de que el sector público empresarial ha perdido peso y relevancia en muchos países desarrollados, sigue persistiendo en múltiples sectores estratégicos y servicios públicos con presencia de fuertes economías de escala, como la provisión de electricidad o las telecomunicaciones. y las empresas públicas siguen jugando un rol esencial en los países en desarrollo. En este sentido, la existencia de argumentos teóricos tanto a favor como en contra de la eficiencia de la empresa pública han aumentado el interés en la realización de análisis empíricos que puedan dilucidar esta cuestión.

En la actualidad, tanto la teoría económica como la evidencia empírica concluyen que, en sectores en los que la competencia puede ser efectiva e intensa, las empresas privadas tienden a ser más productivas y, por tanto, no estarían justificadas por razones de eficiencia (sin perjuicio de que persigan otros objetivos, como la equidad). Sin embargo, en los sectores que, por sus características tecnológicas, tienden a la fuerte concentración en el mercado, la evidencia empírica no es concluyente. ¿Por qué? Aunque no está claro, se considera que, en esos mercados, las empresas privadas sujetas a regulación no tienen por qué ser más eficientes que las empresas públicas; lo que refuerza la necesidad de poner el énfasis sobre la calidad regulatoria, de forma que persiga la nivelación del terreno de juego entre agentes privados y no la necesidad de la provisión directa de estos bienes y servicios.

A lo largo del tema se va a analizar las justificaciones tradicionales y, posteriormente, las limitaciones y potenciales ineficiencias. Asimismo, se van a exponer los pasos a seguir para asegurar que las privatizaciones se llavan a cabo de forma adecuada y transparente, así como también comparar las diferentes situaciones entre jurisdicciones, en torno a la relevancia del sector público empresarial.

Finalmente, va a citarse literatura empírica relevante para responder a la siguiente pregunta: ¿es la empresa pública necesariamente más ineficiente que la privada?

\subsection{Objeto: ¿Qué es una empresa pública?}
Desde un punto de vista económico, la empresa pública se podría definir como una entidad que combina elementos del Sector Público y del Sector Privado.\footnote{%
    Los elementos propios del sector público se manifiestan en tres aspectos: (i) las decisiones en el seno de la empresa son tomadas son tomadas por agentes vinculados al sector público y su criterio no está estrictamente ligado a la rentabilidad financiera ofrecida; (ii) los beneficios de la empresa pertenecen, al menos en parte, al sector público; y, (iii) se espera que la empresa sea responsable por sus acciones ante el conjunto de la sociedad. Por su parte, los elementos propios de la empresa serían: (i) se espera que la entidad sea viable desde un punto de vista financiero en el largo plazo; y, (ii) sus precios deberían estar relacionados, al menos en parte, con los costes de producción.
}

Estos dos últimos criterios nos servirían para distinguir a las entidades que debemos categorizar siguiendo un enfoque económico, como empresa pública, en contraposición a otras entidades públicas productoras de bienes y servicios. Evidentemente, los criterios expuestos pueden darse de forma parcial y, en consecuencia, no ser perfectamente aplicable; por ello, es interesante conocer la definición que se hace de la empresa pública desde un punto de vista jurídico en el marco del Sistema Europeo de Cuentas (2010).

Para vincular una entidad jurídica al Sector Público es necesario verificar que éste mantiene el control efectivo de ésta, teniendo en cuenta para ello los siguientes criterios:
\begin{la}
    \item titularidad pública de la mayoría de los derechos de voto; (ii) control público del consejo de administración o del órgano ejecutivo;
    \item control público del nombramiento o de la revocación del personal directivo; (iv) control público de los principales comités de la entidad;
    \item posesión pública de una acción de oro. La acción de oro es una herramienta que otorga a su poseedor el derecho de veto en determinadas operaciones, aun teniendo una participación minoritaria en el capital social de la empresa. A pesar de este privilegio, la acción de oro no di fiere en cuanto a los derechos económicos que le corresponden, por lo que su titular recibirá los mismos dividendos que el resto de accionistas en proporción a su participación en la empresa;\footnote{%
        Este tipo de acción fue creada en la década de los 80. durante el proceso de privatizaciones de las empresas públicas europeas. Los gobiernos, al ceder su control accionarial a los inversores privados. querían asegurarse de que seguirían teniendo influencia en la toma de decisiones clave. De esta forma, los gobiernos pudieron conservar una posición privilegiada en la gestión ele las compañías privatizadas·. Sin embargo, en 2003 el Tribunal de Justicia de la Unión Europea estableció que las acciones de oro iban en contra de las libertades comunitarias, declarándolas ilegales. En cualquier caso, siguen existiendo en muchos países del mundo.
    }
    \item normativas especiales. Determinadas empresas privadas operan en sectores muy regulados. En la medida en que esta regulación afecte de forma sustancial en las decisiones empresariales. se puede considerar u n factor decisivo de control a efectos de atribuir la titularidad de la entidad al sector público;
    \item la administración representa el grueso de la demanda. Si una empresa destina todas o la graA mayoría de sus ventas al sector público, existe una influencia dominante de la Administración sobre la misma. pudiéndose considerar control. La asignación o no de la titularidad de la empresa al sector público por este canal de control deberá analizarse caso por caso. En este sentido, si la Administración exige límites explícitos con respecto al trato con otros agentes privados, será dificil de argumentar que no exista control público; y/o, 
    \item La administración presta recursos. En ocasiones, la administración exige a sus prestatarios condiciones o controles que van más allá de los que suelen establecer otras entidades financieras. Cuando estas condiciones influyan de forma relevante las actuaciones de la empresa, podrán considerarse también una fuente de control público.
\end{la} 

Se deberá valorar la concurrencia de estos elementos en su conjunto y caso por caso, no siendo necesario el cumplimiento de todos los criterios expuestos, pero tampoco siendo necesariamente suficiente el cumplimiento de tan solo uno de ellos.


\subsection{Argumentos a favor: ¿Por qué existen las empresas públicas?}
A lo largo de la historia se han utilizado distintos argumentos para justificar la existencia de la empresa pública, cuya importancia ha variado con el tiempo, el país o el sector. 

En este apartado se van a presentar, brevemente, algunos de los argumentos aportados en función de si son relativos a: (i) eficiencia; (ii) equidad y (iii) estabilización. Finalmente, va a analizarse también el rol que las empresas públicas pueden jugar en la lucha contra el cambio climático.

\subsubsection{Razones de eficiencia}
La falta de competencia ha justificado tradicionalmente la producción pública en servicios como las telecomunicaciones, la provisión de agua o electricidad y los servicios postales. Existen muchas razones por las que determinados mercados no pueden llegar a tener una estructura plenamente competitiva, y en ese caso puede llegar a justificarse la existencia de la empresa pública si ésta genera menos costes que la regulación o que la ausencia de intervención.

Los clásicos fallos de mercado que más encontramos en la práctica de este tipo de provisión son la presencia de fuertes economías de escala y la presencia de externalidades.

\textcolor{red}{\textbf{OJO!} Aquí también cabría incluir las tésis de MAZZUCATO: infrainversión, incertidumbre radical, etc.}

\paragraph{Monopolio natural}
En presencia de un monopolio natural, el monopolista puede tomar ventaja de su posición de dominio cargar un ma)or precio. Una solución a este problema es que el gobierno se encargue de la provisión del bien o servicio mediante la creación de una empresa rública. Otra alternativa es permitir la provisión privada rero interviniendo de otras formas: mediante la regulación de precios o con el establecim iento de subsidios (o impuestos) para asegurar l a provisión del bien en todos los segmentos del mercado. y no sólo los rentables ( ror ejemplo, <;e podría subsidiar la provisión de servicios postales en ;,onas rurales cuando no les resulte rentable hacerlo a los agentes privados). La existencia de monopolios naturales ha j ustificado, por ejemplo. l a creación de empresas rúblicas en los sectores de los servicios postales, las telecomunicaciones. el transporte y suministro de agua, la construcción y explotación de los puertos o el suministro ele electricidad. En Estados Unidos la regulación ha sido el modo tradicionalmente utilizado para solucionar el problema del monopolio natural (si bien existen empresas públicas en la producción de energía hidráulica y en el suministro de agua). En cambio. en Europa ha sido más común la producción pública.

\paragraph{Externalidades}

La intervención pública puede corregir las ineficiencias derivadas de la presencia de extcrnalidades o bienes públicos a través de la regulación o con la creación de impuestos o subvenciones pigouvianos. pero también mediante la creación dt: una empresa pública que opere en el sector y que se encargue de producir la cantidad óptima. Un ejemplo de este tipo de intervención sería el de las empresas públicas de transporte. En este sentido, puede entenderse que el transporte colectivo de pasajeros en las ciudades genera externalidades en la medida en que descongestiona el tráfico. lo que justifica la creación de una empresa pública que opere en el sector y ofrezca estos servicios de forma subsidiada.

Un caso particular este fallo de mercado, muy relevante en la actualidad, es la presencia de sectores estratégicos. Existen determinados bienes y servicios que se consideran fundamentales para la economía y que justifican la provisión pública. En estos casos. con el fin de evitar una dependencia exterior que pueda llevar en el futuro a estrangulamientos o cuellos de botella. el sector público interviene para garanti7ar una provisión suficiente, adecuada y segura de los mismos. E n este contexto. la empresa acomdería inversiones que, sin realizarse necesariamente en sectores con características de monopolio natural. tendrían la capacidad de aumentar la productividad a largo plazo (por ejemplo, de inversiones en ciertas infraestructuras o en sectores de mayor valor añadido).

Ejemplos de este tipo de intervención, que tuvo una gran relevancia en los países en desarrollo, serían la construcción de plantas de generación eléctrica o la realización de grandes proyectos de transporte en América Latina en los 70 (y posteriormente en África y Asia en los 80). En cualquier caso, muchas de estas infraestructuras te1minaron siendo infrautilizadas y no compensaron su coste de producción\footnote{Mohsin S. Khan: Manmohan S. Kumar (1997).
}

\subsubsection{Razones de equidad}


\paragraph{Servicios públicos esenciales}
A diferencia de las empresas privadas, las empresas públicas no están sometidas necesariamente a las dinámicas del mercado y pueden perseguir objetivos más allá de la maximización de los beneficios, siendo especialmente relevante en ese sentido la provisión universal de determinados servicios públicos a precios razonables. 

Un ejemplo típico de este tipo de intervención es la provisión de servicios ferroviarios en líneas que no son rentables, hasta el punto de que tan solo unos pocos países europeos son capaces de recuperar los costes de provisión de estos servicios mediante el cobro a los usuarios, tratándose en estos casos de países con una elevada densidad de población (FMI, 2020). En cambio, en muchos otros países las disparidades regionales siguen siendo muy relevantes (por ejemplo, en Francia el 15\% del coste de provisión de servicios ferroviarios se concentra en el 2\% de los usuarios) y son comunes los subsidios directos e indirectos a las empresas estatales en el sector.\footnote{%
    Por ejemplo, en los países de Europa del Este, que por lo general tienen una baja densidad de población y un bajo nivel de utilización de la infraestructura ferroviaria, solo se recupera en torno al 20\% de los costes de provisión mediante el cobro de tarifas a usuarios. En esta línea, también ha sido habitual la condonación de deuda. Un caso paradigmático fue la condondonación de deuda a la principal gestora de infraestructura ferroviaria en Grecia, que alcanzaron el 7\% del PfB en 2011.
}

Como alternativa, el objetivo de servicio universal puede alcanzarse mediante la concesión de subsidios a empresas privadas que operen en segmentos del mercado que no sean rentables, junto con otros instrumentos como las obligaciones de servicio público o la concesión de ayudas específicas a usuarios de baja renta. Sin embargo. la factibilidad de estas alternativas depende de la capacidad de implementarlas por parte del sector público. e implican también coses ligados a la definición de las zonas geográficas afectadas, la calibración de los pagos. la regulación y el control a los proveedores. Por ello, en función de las características del país, puede resultar más conveniente la provisión del servicio mediante empresas públicas.

\paragraph{Desigualdad horizontal: Corrección de disparidades regionales}
Los cambios tecnológicos han alterado la geografía de la producción. con un aumento de la importancia de las grandes ciudades. En cambio. otras ciudades pequeñas alejadas de las grandes urbes sufren procesos de despoblación, lo que aumenta las disparidades regionales. En este contexto. las empresas públicas pueden configurarse como una forma de apoyar económicamente a las regiones más desaventajadas, generando empleos que corten las dinámicas de despoblación, creando una mayor demanda de bienes y servicios y promoviendo con ello el afloramiento de más empresas privadas. también creadoras de empleo. Se entiende que este efecto va a ser mayor cuanto más cualificados sean los trabajadores trasladados.

A modo de ejemplo, puede destacarse la apertura de distintos establecimientos públicos en el este de Alemania tras la reunificación, o el más reciente plan de desarrollo de la región de Bavaria, que consistirá en el traslado de más de 50 entidades públicas a las zonas rurales de este estado.

Este tipo de medidas de política regional se han llevado a cabo en múltiples países.\footnote{%
    Como ejejemplos adicionales, destacan: en Reino Unido, el traslado de parte de la BBC a Salford (Manchester); en Corea del Sur, dos tercios de las agencias estatales se trasladaron de Seúl a otras ciudades menos pobladas; y en Noruega se han trasladado multiplicidad de organismos públicos, como la autoridad de la competencia, a distintas ciudades.
}

\subsubsection{Razones de estabilización: }
Tras la Segunda Guerra Mundial, el auge del intervencionismo estatal en Europa hizo que aumentara el rol de las empresas públicas como mecanismo de intervención del estado en la economía. En esta línea argumental, entendiendo que la inversión privada es volátil y se contrae en periodos de crisis, contar con empresas públicas permitiría acometer inversiones directamente por parte del Estado, lo que conferiría una mayor capacidad de impulsar el crecimiento en tiempos de crisis, efecto que se suma a la mayor estabilidad cíclica del empleo público.

En un marco más actual, múltiples estudios muestran que los niveles de inversión y de empleo de las empresas públicas tienden a ser menos sensibles al ciclo económico. De acuerdo con el \textit{Life in Transition Survey}, el porcentaje de trabajadores con contratos indefinidos es mayor en las empresas públicas que en las privadas, además de haber reportado una menor probabilidad de perder el empleo durante la crisis financiera de 2008. Estas diferencias son estadísticamente significativas incluso cuando se tienen en cuenta características individuales como la edad, el género o el tamaño de la empresa.

Por el lado de la oferta, las entidades financieras de titularidad pública son prestadores más estables en tiempos de crisis. Además, la evidencia sugiere que los trabajadores públicos ahorran más, lo que les puede permitir estabilizar en mayor medida su consumo y contribuir con ello a la suavización del ciclo económico..\footnote{%
    También de acuerdo con el Life in Transition Survery, el porcentaje de trabajadores que afirma tener ahorros para poder hacer frente a reducciones temporales de su renta es mayor entre empleados públicos que entre sus homólogos en el sector privado, para niveles de renta comparables. Por tanto, parece que la estabilidad en los flujos de renta de los trabajadores públicos (que les confiere una mayor capacidad de ahorro) compensa su menor necesidad de contar con ahorro precautorio (ligada a la menor volatilidad de sus ingresos), generando en conjunto un efecto positivo.
} 

Por tanto, las empresas públicas pueden funcionar como estabilizadores automáticos ante shocks económicos adversos. Este efecto positivo de cobertura frente al riesgo se contrapone a la ineficiencia relativa de las empresas públicas que, como se analizará en los siguientes apartados, pueden presentar menores niveles de productividad. Por tanto, existirá en este sentido un trade-off entre estabilización y motor de crecimiento.

\paragraph{Razones climáticas}
Algunas de las mayores empresas públicas a nivel global operan en el sector de la energía. Esta realidad está acrecentando el debate en torno al rol que pueden (o deben) jugar en la transición ecológica. 

Estudios recientes han mostrado que las empresas energéticas estatales prestan más atención a consideraciones medioambientales.\footnote{Por ejemplo, véase: BARNES (2019), PAN et al. (2020) y PRAG et al. (2018).
} Además, la evidencia sugiere que las empresas estatales presentan una mayor probabilidad de controlar sus emisiones y tener objetivos de reducción de las mismas; aunque esto contrasa, con que no se aprecien diferencias significativas en consumo energético y en su proceso de descarbonización si se comparan con sus homólogas en el sector privado (FMI, 2020).

Merece mención especial el rol de las empresas públicas en el sector del carbón que sigue representando en torno a un tercio de la generación eléctrica a nivel mundial a pesar de su elevado impacto ambiental. Se estima que el 60\% de las minas y centrales eléctricas de carbón son de titularidad pública. Por tanto, las empresas públicas ostentan un rol esencial en su necesaria reconversión. En este sentido, la mayoría de empresas públicas europeas ya han reducido el consumo de carbón como resultado de la introducción de estándares medioambientales más estrictos y de nuevos esquemas como el \textit{EU Emissions Trading System} (EU-ETS). Teniendo en cuenta la importancia social de la minería del carbón en algunas regiones, será de vital importancia la elaboración de planes de reconversión que reduzcan su impacto negativo sobre el empleo.\footnote{%
    Además. el empleo en la minería de carbón tiende a estar fuertemente concentrado, suponiendo. En algunos casos, se sitúa entre el 10\% y 20\% del empleo; por ejemplo, en algunas regiones en Bulgaria, Grecia o Polonia.
} 

Son buenos ejemplos de este tipo de políticas la conversión de minas en centrales petroquímicas en Países Bajos en los 70, así como la creación en Alemania de más empleos públicos en las regiones afectadas por la reconversión.

\subsection{Argumentos en contra:}
A pesar de que, como se ha visto en el apartado anterior, se puede defender la existencia de la empresa pública por motivos de eficiencia. existen argumentos teóricos que señalan su mayor ineficiencia con respecto a la empresa privada. Estos argumentos se apoyarán tanto en la teoría de la agencia y en la teoría de la elección pública como en, a efectos más prácticos, las m ayores restricciones legales y burocráticas a las que se enfrentan las entidades públicas empresariales a la hora de operar en el mercado.



\subsubsection{Eficiencia: }


\paragraph{Desalineamiento de incentivos: Diferencias en la organización interna}
El hecho de que una empresa sea pública o privada cambia de forma sustancial los distintos incentivos que se dan en la toma de decisiones, alterando su comportamiento. Existen motivos teóricos para pensar que la empresa pública va a ser relativamente más ineficiente en ciertos aspectos, si bien. como se verá posteriormente, se requerirá un análisis empírico que evalúe si los argumentos utilizados para justificar su existencia por razones de eficiencia pueden llegar a compensar o no estas ineficiencias. 

En este sentido, las principales críticas a la empresa pública en términos de eficiencia han venido de la mano de la teoría de la elección pública y de la teoría de la agencia.

En primer lugar, la teoría de la elección pública se ha encargado de analizar los incentivos relativos de burócratas y políticos (ambos encargados, de forma más o menos directa, de la toma de decisiones en la empresa públicos) para perseguir objetivos de interés público u obj etivos particulares en el contexto de la gestión empresarial. De esta forma: 

VICKERS y YARROW (1988). defienden que los políticos pueden perseguir la maximización de los votos, lo que les puede llevar a incurrir en ineficiencias como: (i) mantener puestos de trabajo por encima de los que llevarían a la maximización de los beneficios: (ii) pagar salarios superiores a la productividad marginal de los trabajadores; o (iii) favorecer con sus decisiones a grupos de presión o a ciertos colectivos específicos en contra del interés general

Por su parte, NISKANEN (1975) argumenta que los burócratas pueden perseguir la maximización de una función objetivo que incluya el bienestar social (el agregado de los beneficios de la empresa y el excedente del consumidor que genera su actuación) pero también sus preferencias personales. ligadas. por ejemplo, al prestigio o al reconocimiento público.

En cualquier caso, aunque los agentes privados tienen incentivos distintos a los de los políticos y los burócratas, tampoco tienen incentivos que se encuentren estrictamente alineados con los de los propietarios de la empresa. Por ello, se tiene que recurrir a la Teoría de la Agencia para analizar las diferencias, no sólo en el sistema de incentivos sino también en los sistemas de control que subyacen a cada relación contractual, para poder profundizar. en su caso. en las ineficiencias relativas de la empresa pública.

En la empresa privada, se entiende que existe una relación de agencia cuando un principal (propietario) le encarga al agente (directivo) que tome decisiones en su nombre para alcanzar determinados objetivos empresariales. Los problemas aparecen cuando: (i) los incentivos del agente no están completamente alineados con los del principal; y, (ii) existe asimetría informativa en torno a las actuaciones del agente, que no son perfectamente observables. En este sentido, el propietario (principal) va a intentar establecer sistemas de control sobre el agente y a alinear sus incentivos para minimizar las ineficiencias derivadas de su actuación (los costes de agencia).

En la empresa pública pueden acentuarse las ineficiencias derivadas del problema de agencia porque existe una doble relación de agencia: (i) los titulares de la empresa pública, que son en última instancia los ciudadanos, delegan en los políticos la gestión de la misma; y, (ii) los políticos delegan la gestión de la empresa en unos directivos (normalmente burócratas) con respecto a los que existe también una relación de agencia análoga a la de la empresa privada. Además, existen algunos factores adicionales a tener en cuenta que, en el marco de la teoría de la agencia aumentan la ineficiencia relativa de la empresa pública a nivel teórico:

- La empresa pública persigue, a menudo, objetivos distintos a la maximización del beneficio, por lo que su labor de evaluación y control externo resulta difusa y difícil. Dicho de otro modo, cuando una empresa pública presente una rentabilidad más baja a la de sus análogas privadas no implica que sea más ineficiente si en paralelo ha podido satisfacer otros objetivos sociales igualmente deseables.

- En las empresas privadas cotizadas, la exposición a los mercados de capitales aumenta la capacidad de control sobre los gestores en tanto que: (i) cualquier accionista puede vender sus acciones si está disconforme con la gestión realizada, de modo que cualquier indicio de ineficiencia que sea apreciable por parte de los accionistas se verá inmediatamente reflejado en el precio de las acciones; (ii) otros participantes en el mercado, aunque no sean propietarios de la empresa, van a ejercer un control indirecto sobre la gestión de la misma porque reducirán la demanda de participaciones en caso de percibir ineficiencias; y, (iii) existe la posibilidad de que los gestores sean apartados de la empresa mediante una oferta oferta de adquisición, lo que aumenta sus incentivos a ser más eficientes.

Estos sistemas de control no existen, o existen de forma tan solo muy parcial, en la empresa pública, puesto que: (i) los ciudadanos como titulares últimos de la empresa pública no pueden vender sus participaciones; (ii) no existe una señalización de mercado análofa al precio, donde pueda reflejarse el buen o mal comportamiento de los gestores; y, (ii) aunque algunas empresas públicas coticen en bolsa, el hecho de que el estado ostente su control (bien siendo accionista mayoritario o bien con una acción de oro) hace que sea imposible la destitución del equipo directivo mediante una OPA.

Además, es posible que la empresa pública se enfrente a una restricción presupuestaria blanda: en caso de quiebra, existe la posibilidad de que sea rescatada. Esta posibilidad genera un problema de riesgo moral en los gestores y en los políticos, que no internalizan las consecuencias de una posible quiebra del mismo modo en que lo hacen los gestores en la empresa privada y los accionistas, lo que puede llevar a un comportamiento ineficiente.


\paragraph{Rigideces: Contratación, procedimientos y financiación}
\textcolor{red}{
\textbf{OJO -> Esto es muy burdo, existe evidencia empírica que demuestra cómo la contraación pública resulta, en muchos casos, en procedimientos competitivos y muy ajustados al precio}. 

De manera similar. para asegurar que no existen abusos en la gestión de la hacienda pública, los gestores públicos se enfrentan a importantes restricciones a la hora de tomar ciertas decisiones, como pueda ser la contratación de determinados servicios o la compra de bienes. En muchas ocasiones se establece la obligación de someter estas compras a concurso público, algo que. si bien es necesario, puede introducir costes adicionales ligados a: (i) el propio coste burocrático del proceso, que deviene más complejo: (ii) el coste de oportunidad ligado a la duración del
procedimiento de contratación y (iii) el coste burocrático que soportan las empresas.

\textbf{OJO: Esto es interesante, pero hay que reestructuralo bien}

La empresa pública y la privada también difieren en la configuración de sus presupuestos. en especial en lo referido a las inversiones a largo plazo. En concreto. cuando existe una oportunidad de crecimiento e inversión, la empresa privada puede conseguir financiación de fonna relativamente rápida, mientras que es posible que la empresa pública requiera de una modificación presupuestaria que va a requerir un mayor plazo. teniéndose incluso que esperar a la aprobación de los presupuestos del año siguiente. 

En la empresa privada, los propietarios persiguen su interés panicular y para ello eligen gestores que actúen en consecuencia pudiendo prescindir de ellos cuando su gestión sea ineficiente; también pueden remunerarles con el salario que estimen conveniente de acuerdo con su productividad.

En cambio. en la empresa pública (y en el conjunto del sector público) el interés particular que se persigue no es más que el interés general, puesto que la buena gestión de la hacienda pública impacta de forma positiva sobre el conjunto de los ciudadanos. Para asegurar que los trabajadores públicos no abusan de su poder o se apartan del interés general, se establecen procedimientos estrictos para su contratación y se impone el principio de mérito y capacidad en el acceso a la función pública. 

Todas estas restricciones, si bien son necesarias, introducen rigideces tanto a la hora de despedir como a la hora de pagar salarios mayores para atraer talento, de modo que impactan negativamente en la productividad relativa de la empresa pública y en su capacidad de competir.
}


\subsubsection{Ineficiencias autoinfligidas: }



\textcolor{blue}{\textbf{OJO: Aquí lo que sería interesante es introducir las tésis de MAZZUCATO (2013)}}

A priori, identifico dos posibles ineficiencias autoinfligidas: la infrainversión y el sobredimensionamiento del riesgo (extremada aversión al riesgo, infligida). Aunque hace falta desarrollar, la idea central es la siguiente: excesivo castigo/escrutinio al fracaso -> desconexión total del riesgo (en el sentido de que la gente exige resultados, cuando los proyectos per se contienen un elemento de riesgo) -> infrainverisón y desilusión -> efectos sobre directivos, empleados, organismos y ciudadano. (Este esquema es provisional, pero lo importante es la idea de ineficiencias autoinfligidas)







































\clearpage
\section{Empresa Pública: Política de privatización}
\puntosclave{
    Cuando una empresa pública pierde su razón de ser, suele llevarse a cabo un proceso de privatización, a través del cual se transfiere al sector privad. Este también puede iniciarse por razones presupuestarias. 
    
    Una vez se toma la decisión de privatizar, es necesario asegurar que el sector público vende la empresa a un precio que refleje su valor de mercado, y que el sector en el que ésta opera va a seguir funcionando adecuadamente. Para ello, es fundamental contar con una buena regulación del mercado, y que el proceso de privatización se sLuete a los máximos estándares de transparencia.
}

\textcolor{red}{\textbf{OJO: Un buen nexo con el apartado anterior es presentar la victoria de los argumentos en contra, lo que ha conducido al estudio/análisis del proceso de privatización/desinversión: ¿cómo se puede hacer de la mejor forma?}}

Todos estos problemas en el funcionamiento de la empresa pública sirvieron para justificar procesos de privatización, especialmente prevalentes a partir de los años 80.\footnote{%
    Este proceso se produjo en muchos países desarrollados, pero tuvo su mayor alcance y extensión en Reino Unido, donde las empresas públicas pasaron de representar un 8% del empleo en 1979 a un 3'1/o en 1992.
} 

Se ha visto que existen argumentos tanto a favor como en contra de la privatización. En cualquier caso, una vez tomada la decisión de privatizar, es relevante entender los pasos a seguir para que se lleve a cabo de la mejor forma posible, entendiendo que los objetivos para el sector público a lo largo del proceso serían los siguientes:
• Maximizar los ingresos para el sector público: en principio. el estado debería intentar vender la empresa pública a un precio que refleje su valor fundamental. Sin embargo. una serie de factores que pueden hacer que la empresa se llegue a vender por un precio sustancialmente inferior a su valor teórico, entre los que destacamos (i) la asimetría informativa e incertidumbre en torno a .los rendimientos potenciales de la empresa y (ii) las condiciones de liquidez del mercado. Por ello. a lo largo del proceso de privatización se tendrán que tomar una serie de decisiones estratégicas con el fin ele minim izar su impacto.
• Asegurar el buen funcionamiento posterior del mercado en el que opera la empresa pública: cuando hablamos de privatización en sectores que no son plenamente competitivos. deviene esencial la regulación aplicable. como se ha explicado a lo largo del terna.

En general, el proceso de privatización consiste en tomar una serie de decisiones en las que podrán existir trade-offs que hay que considerar. En todo caso, la clave será siempre la transparencia y la rendición de cuentas.


\textcolor{blue}{Privatización: ¿por qué llevarla a cabo?}
Entendemos por privatización la venta de activos públicos a agentes privados que deberán pasar a ostentar la propiedad mayoritaria de la empresa. Este proceso se debe distinguir de la liberalización, que sería el proceso mediante el cual se abren a la competencia mercados que antes estaban regulados y dominados por un único actor. La privatización no implica necesariamente liberalización, del mismo modo que la privatización no es estrictamente necesaria para llevar a cabo un proceso de liberalización. Teniendo en cuenta que a lo largo del tema se han expuesto argumentos tanto a favor de la empresa pública como en su contra por razones de eficiencia, es necesario profundizar en el análisis teórico de la privatización en esos términos.

\textcolor{yellow}{La privatización por motivos de eficiencia}

La existencia de monopolios naturales o economías de escala que impiden la competencia en un determinado sector ha sido el principal argumento en favor de la existencia de la empresa pública en términos de eficiencia. 

En este sentido, cuando la empresa se encuentra en un sector que es competitivo o puede serlo potencialmente, no se puede justificar la existencia de la empresa pública por motivos de eficiencia (sin perjuicio de que se pueda mantener la titularidad pública de la empresa para perseguir otro tipo de objetivos, como la equidad). En cambio, en sectores donde la competencia no es posible o va a ser escasa, la alternativa a la que se enfrentan los poderes públicos no es un monopolio estatal en contraposición al libre mercado, sino en contraposición a la regulación (que sería necesaria en un contexto de escasa competencia).

En ese marco, va a ser difícil determinar desde un punto de vista teórico si la privatización implica ganancias netas de eficiencia. de modo que tendrá que acudirse a la evidencia empírica, como veremos posteriormente.


\textcolor{yellow}{Razones presupuestarias}
En la práctica, el cumplimiento de ciertos objetivos presupuestarios ha servido para justificar multiplicidad de procesos de privatización. 
- Las empresas públicas impactan de forma relevante en las finanzas del estado: en caso de acumular pérdidas, pueden suponer pasivos potenciales y mayor gasto público tanto en el presente o en el futuro, mientras que las empresas rentables son activos valiosos susceptibles de generar ingresos públicos.

Si una determinada empresa se clasifica como titularidad del estado en el marco del SEC-2010, sus ingresos, gastos y deuda se añaden a los del resto del sector público y, por tanto, afectan al cómputo del déficit y de la deuda del estado. En este sentido, debe tenerse en cuenta que: 
• En muchos casos, la participación del gobierno en empresas públicas es positiva para las finanzas públicas. Por ejemplo, en Finlandia. los ingresos públicos obtenidos a través de sus empresas alcanzaron en promedio el 1 ,5\% del PIB entre 2005 y 20 1 4. Otros estados como Suecia, Estonia o Malta, tuvieron ingresos de hasta el 1\% del PIB. En este marco, la posibilidad de privatizar activos del estado que son rentables puede impactar positivamente en la reducción de la deuda, disminuyendo también la necesidad de financiación del sector público. En cualquier caso, debe tenerse en cuenta que este tipo de operaciones, tanto si se refieren a la privatización de empresas públicas como a la venta de activos públicos en general. suponen una entrada de flujos de caja en el momento de la venta. y una reducción de los mismos en todos los periodos sucesivos. de modo que hablamos, en el fondo, de una fom1a de financiación, como pueda ser la emisión de deuda. Que esta forma de financiación sea deseable o no dependerá de si (i) con ella el estado consigue obtener liquidez en condiciones más ventajosas que mediante las emisiones de deuda. lo que será más probable en periodos de tensión financiera y (ii) de si existen ganancias netas de eficiencia derivadas de la privatización y, además. éstas se reflejan en el precio de venta y por tanto el estado tiene la capacidad de apropiárselas. Estas condiciones pueden cumplirse o no en función del momento y las circunstancias de la venta. de modo que no en todos los casos la privatización va a suponer una ganancia patrimonial neta para el estado, aunque suponga a corto plazo una inyección de liquidez.
• Por otro lado, en algunos casos la participación en el capital de determinadas empresas públicas supone un coste y un pasivo potencial para el gobierno. Cuando el gobierno asume la titularidad de una empresa, asume la titularidad de su deuda, y si la empresa llega a quebrar o a atravesar dificultades, puede que tenga que hacerse cargo de la misma. En este sentido, la necesidad de reducir los costes ligados a la financiación de estas empresas o como forma de reducir pasivos contingentes también ha supuesto tradicionalmente un motivo para la privatización.

Por tanto, la decisión de privatizar deberá ser tornada caso por caso, sin existir recomendaciones de política económica únicas y universalmente aplicables: en la privatización por motivos de eficiencia, se ha visto que las ganancias en productividad van a depender en gran medida del grado de competencia en el mercado; en la privatización por motivos presupuestarios, se ha visto que, para que su impacto sea positivo no solo a co110 sino también a largo plazo. se tienen que cumplir una serie de precondíciones que deberán evaluarse en cada caso particular.

\subsection{Pasos previos: ¿cómo preparar el proceso?}

1) Responsabilidad administrativa en el proceso de privatización

El gobierno debe garantizar que las unidades administrativas encargadas del proceso de privatización se sujeten a altos estándares de transparencia. Cuando el volumen de activos a privatizar sea elevado, puede ser recomendable crear una agencia de privat i::: ación especializada en el proceso y que ostente la responsabilidad derivada del mismo. Si el volumen de activos no es suficiente como para que la creación de este organismo resulte económicamente conveniente. se debe identificar claramente la unidad administrativa encargada de la privatización para sujetarla a los máximos estándares de transparencia.

2) El rol de los asesores externos

En ocasiones, los gobiernos pueden requerir la ayuda de asesores externos con mayor experi encia para diseñar el proceso de privatización. El mayor reto en este tipo de situaciones va a ser el garantizar que no incurran en conflicto de intereses y que tengan unos incentivos completamente alineados con los del sector público. En este sentido, entre los países OCDE, Australia, Polonia y Espafia se caracterizarían por encargar análisis externos de las em presas a privatizar. siendo en los dos últimos obligatorio por ley.

Por lo general. la decisión de contar con estos asesores debe tomarse caso por caso, sin evitarla por no incunir en costes a corto plazo, y sin utilizarla corno herramienta esencial y como fonna de reducir la responsabilidad del gobierno en el proceso. En todo caso, va a ser esencial evitar los conflictos de intereses, establecer procesos de selección claros y transparentes y disefiar sistemas de incentivos adecuados.

3) Regulación y competencia
Por lo general, no se recomienda acometer el proceso de privatización sin haber establecido antes un marco regulatorio idóneo. Cuando no sea posible la introducción de competencia en el mercado. el estado debe asegurar que existe una buena regulación previa. que deberá estar diseñada y ser gestionada por una unidad distinta a la encargada del procedimiento de privatización. Idealmente, debería encomendarse la regulación a un órgano dotado de autonomía con respecto al gobierno.

4) Condiciones laborales
Los cambios en las condiciones laborales de los trabajadores y las reducciones de plantilla al privatizarse las empresas públicas son una práctica habitual, y constituyen a menudo uno de los principales obstáculos al proceso de privatización. en la media en que pueden generar problemas jurídicos y problemas sociales. 

Cuando la empresa pública no tenga forma societaria y sus trabajadores sean funcionarios, se generará un problema jurídico para el sector público. que tendrá que recolocar a estos trabajadores o elaborar disposiciones legales específicas para amparar su traslado a la empresa privada. Si la empresa pública tiene forma societaria no deberían existir problemas jurídicos, pero sí existirán. probablemente, problemas ligados a protestas sindicales que pueden dificultar o incluso impedir el proceso de privatización.

En consecuencia, los po/icvmakers tendrán que encontrar un equilibrio dentro del siguiente tradeoff: si se alcanzan acuerdos por los cuales los trabajadores van a mantener, total o parcialmente. sus condiciones laborales una vez la empresa se privatice (así como acuerdos que garanticen el mantenimiento de la plantilla) se facilita el proceso de privatización desde un punto de vista social, pero a la vez se limitan las mejoras de eficiencia que se podrían alcanzar con el proceso. 

Tradicionalmente. una de las formas más comunes de gestionar las reducciones de personal en privatizaciones ha sido la concesión de prestaciones por jubilación o por despido. Un ejemplo paradigmático de la utilización de este tipo de técnicas fue la privatización de Bra:::i/ Railways en el afio 2000, en la que el gobierno inició la labor de reestructuración de la plantilla antes de la privatización. con el objetivo de atraer inversores y mitigar el impacto social de las reformas. En este caso. de los 42.000 trabaj adores que integraban la empresa. 1 2.000 optaron por la j ubilación anticipada y otros 6.000 abandonaron voluntariamente la empresa. El proceso fue llevado a cabo con varias refomrns de la seguridad social y su coste fue financiado por el gobierno y por diversas instituciones i nternacionales. El proceso se inició en 1 998 y en mayo del 2000 la plamilla de la empresa se había reducido en un 80\%, habiéndose considerado el programa. en términos generales, un éxito en la planificación del proceso de privatización desde la óptica de su impacto social.\footnote{Estacht\ A., de Schmitt. A. J. A., \& Sydenstricker. E. (2000).
}


5) Reestructuración de la empresa
Existen importantes consideraciones en torno a si el sector público debe intentar reestructurar la empresa pública antes de privatizarla. entendiendo esta reestructuración en sentido amplio (desinversión de activos poco rentables, reorientación de actividades, etc.). Tradicionalmente, se había entendido que estas reformas no debían acometerse antes de la privatización.\footnote{Véase. por ejemplo, Chong y Lopez-de-Silanes (200
} al considerarse que los futuros compradores estarían más capacitados para llevarlas a cabo (conociendo el m ercado en el que operan) y, por tanto, el coste de estas reformas. que se iba a descontar del precio, sería inferior al coste de llevar a cabo estos mismos ajustes en el seno de la Administración.

Sin embargo, estas consideraciones parten de la base de que el agente que va a adquirir la empresa pública está identificado y opera en el sector. de modo que conoce la rentabilidad potencial que la empresa puede llegar a tener tras la privatización. En cambio. cuando no exista un comprador identificado o la empresa vaya a privatizarse vía oferta pública de 1·enta ( ver posteriormente) puede ser conveniente realizar ajustes previos para sei'íali::ar los flujos de caja que la empresa a privatizar puede llegar. potencialmente. a generar. Además. la valoración de la empresa que hagan los agentes compradores estará afectada por la ince1tidumbre en torno a las posibilidades de reestructuración que lleguen a realizarse. existiendo distintos escenarios alternativos. Si el sector público inicia estas reformas. puede reducir la incertidumbre en torno al devenir de las m ismas, además de mitigar la asimetría informativa que existe con respecto al funcionamiento interno de la empresa. lo que en conjunto puede aumentar la disposición de los compradores a pagar por la empresa.

Por tanto, la elección deberá hacerse caso por caso, teniendo en cuenta las características de la empresa y también la forma de venta que se va a adoptar. Normalmente, en ventas directas la reestructuración no va a ser necesaria para señalizar el valor fundamental de la empresa, m ientras que en OPYs puede llegar a serlo. Por otro lado. existirán otras alternativas factibles como privatizar secuencialmcnte,\footnote{%
    Primero se privatizan una pane de las acciones de la empresa, de forma que entran accionistas privados que inician el proceso de reforma, y cuando éste ha avanzado (aumentando la rentabilidad o, al menos, señalizando una mayor rentabilidad potencial) se privatizan el resto de acciones.
} contratar a empresas dotadas de credibilidad para que valoren la empresa pública y lo anuncien al público.\footnote{%
    En realidad. la valoración del precio de la empresa antes de su privatización es una prácti.ca habitual en la gran mayoría de los países de la OCDE, siendo incluso obligatoria por ley en algunos estados como Francia, Italia. Polonia o España). Las diferencias principales se encuentran en las entidades encargadas de esta valoración, siendo que en la mayoría de países la llevan a cabo empresas del sector privado, y en unos pocos (como Corea del Sur o Turquía) la valoración es realizada por comisiones especializadas dentro del gobierno, con el objetivo de evitar los conflictos de intereses que podrían producirse al valorar la empresa en el seno de una entidad privada.
}

\subsection{Venta: ¿Qué forma elegir?}
Una vez expuestas las decisiones a tomar antes de iniciar la privatización, van a exponerse los pasos a seguir a lo largo del proceso:

Existirán distintas formas de realizar la venta, entre las que distinguiremos dos fundamentales: (1) la oferta pública de venta y (2) la venta directa. Cada una de ellas tendrá una ventajas e inconvenientes específicos. de modo que su conveniencia deberá evaluarse caso por caso.

\subsubsection{Oferta pública de venta (OPV}
Una OPV representa la incorporación de una empresa al mercado de valores al emitir acciones para atraer capital, siendo su compra accesible a cualquier inversor en la sociedad. Se habla de OPV inicial (Oferta pública inicial. o IPO. por sus siglas en inglés) cuando la empresa cuyas acciones se emiten no cotizaba antes en los mercados de capitales. La emisión de acciones a los mercados permite a la empresa pública captar financiación (que podrá ser utilizada para nuevas inversiones). y se podrá hablar de privatizac ión completa cuando los nuevos entrantes en la empresa pasen a suponer más del 50% de su capital social.

La JPO requiere que la empresa pública que busca pri vatizarse sea de suficiente tamafio y esté lo suficientemente saneada como para generar una demanda de acciones suficiente. También requiere que los mercados nacionales de capitales sean amplios y profundos, o en su defecto emitir las acciones también en mercados extranjeros. Normalmente las acciones se ofrecen tanto a inversores institucionales como minoristas y el precio de las acciones es fijo (decidido por el gobierno o por el comité en el que haya delegado el proceso).

La principal ventaja de este sistema es su transparencia, motivada por los fue11es requisitos re,gulatorios que se imponen a las empresas cotizadas y por el hecho de que la compra de acciones de la empresa vaya a ser accesible para cualquier inversor en la sociedad. Además. la disciplina que posterio1mente establecen los mercados de capitales sobre la empresa cotizada, como se ha explicado a lo largo del tema, va a imponer restri cciones más rígidas sobre su com portamiento.

El proc edimiento a seguir para la hora de lanzar una ofe11a pública de venta inicial sería el siguiente:

(i) Preparación del folleto: Se trataría de un folleto infonnativo que recoja de fo1ma veraz la principal información relativa al emisor y a las características de las acciones emitidas.

Las Administraciones Públicas deberán valerse de asesores financieros para su elaboración. suministrándoles toda la información que pueda ser necesaria.
(ii) Formación del sindic o/o Je colo cución.
(iii) Aprobación del follet o.
(iv) Lanzamiento de la campaña de venta. Consistirá en dar publicidad a la venta de las acciones, con la participación del gobierno, con el objetivo de maximizar su demanda.
(v) Aprobación del rango de precios. Teniendo información en torno a la demanda de acciones existente. el sector público va a fijar un precio mínimo por encima del cual está dispuesto a vender sus acciones, rechazándose las ofertas que se encuentren por debajo de dicho valor.
(vi) Apertura a la venta: Los distintos agentes interesados en comprar acciones van a hacer sus ofertas. que serán tomadas en consideración en la medida en la que se encuentren por encima del precio mínimo establecido. A partir de aquí, pueden ocurrir dos cosas: a. Si no se alcanza un importe mínimo de suscripción (es decir. si el capital social total suscrito en la OPV se encuentra por debajo de un mínimo aceptable) la Administración podrá decir extender la oferta en el tiempo o bien paralizar el proceso de privatización. b. Si se alcanza este importe mínimo. se establece el precio de adjudicación y se asignan las acciones a aquellos oferentes que se encontraran por encima de la misma.

En muchas ocasiones, si el capital social a asignar es especialmente elevado, el proceso de privatización se puede realizar en dos etapas: primero con una OPV inicia y posteriormente completándola con una OPV secundaria, cuando una parte de las acciones ya estén cotizando en bolsa. La OPV secundaria podrá realizarse también cuando la empresa pública ya cotizara en bolsa antes de iniciar el proceso de privatización. pero el sector público siguiera siendo el titular de la mayor parte del capital social de la misma.

En general la privatización secuencial se con:;idera adecuada cuando es poco probable que la OPV resulte en la colocación completa de las acciones ofe11adas. En este sentido, tiende a considerase que la realización de ofertas iniciales de venta reducidas, seguidas de posteriores ofertas secundarias y terciarias. puede aumentar los ingresos públicos derivados de la privatización, al i ncrementarse el precio de las acciones como consecuencia de la mejora en la gestión de la empresa parcialmente privatizada.\footnote{De acuerdo con la evidencia general recogida en OECD (2009).
} Un ejemplo histórico de esta práctica sería la privatización de ENEL ( Ente Nacional para la Energía Eléctrica) acometida en Italia durante los años 90, debido al elevado tarnai'io de l a empresa.

Como potenciales inconvenientes de las IPO como instrumento de privatización, debe mencionarse el hecho de que no permiten al sector público seleccionar a inversores específicos con know-how en la materia, sumado a la posibilidad de que las acciones de la empresa pasen a estar muy diseminadas en distintas manos (capitalismo pop;lar) lo que podría acentuar las ineficiencias generadas por el problema de agencia que surge entre accionistas y gestores.\footnote{%
    El capitalismo popular es una teoría económica que defiende la participación masiva de la población en la propiedad de acciones en empresas públicas o privadas. Este objetivo fue especialmente relevante en el diseño de las estrategias de privatización en algunos países, como Francia en los años 80-90. en los que se intentaba que una gran parte de la población tuviera acceso a la compra de acciones de las empresas privatizadas. Como inconvenientes, debe destacarse que este método genera una fuerte displ,rsión de l a propiedad de acciones que pued!: plantear problemas en l a gestión d e l a empresa, ya que puede d-ificultar la toma de decisiones al no existir un núcleo determinado de accionistas que ostente el control de l a empresa.
}

\subsubsection{Venta directa}

Implica la venta de las acciones propiedad del gobierno a un tercero, que normalmente será un inversor considerado estratégico. En muchos casos, el inversor estratégico es una cm presa que ( I ) tiene un tamaño elevado y una m ayor disponibilidad de recursos para invertir y (2) opera en el sector de la empresa pública que se está privatizando. En este sentido, en la medida en que la venta directa pennite al sector público tener en cuenta otros criterios (más allá del precio) a la hora de elegir al nuevo i nversor. va a facilitar:
• Asegurar que la gestión sea realizada por equipos que tengan experiencia en la provisión del bien o servicio objeto de la empresa pública privatizada.
• Incorporar las mejoras tecnológicas y el know how que haya podido desaiTOllar con anterioridad la empresa compradora. permitiendo una mayor eficiencia.
• Comprometer una serie de inversiones para garantizar la continuidad en el bien o servicio producido, o incluso diseñar una estrategia para su mejora (no pudiendo imponerse estas condiciones en una OPV).


Existen, a grandes rasgos, dos formas de realizar la venta directa:

\paragraph{Subasta} 
Se toman en consideración distintas ofc11as realizadas por inversores que se consideren elegibles para la adjudicación. fa un método más transparente y competitivo que. no obstante. sólo será realizable cuando exista más de un inversor con las características idóneas para asumir la gestión de la empresa privatizada. algo que no siempre tiene porqué ocurrir. No es poco habitual que los estados cuemen con disposiciones legales que recojan la existencia de otros objetivos a tener en cuenta en la privatización. más allá de la maximización del precio. como el mantenimiento de la plantilla de empleados. evitar la deslocalizac ión de la empresa, o impulsar la inversión en l+D. Sin embargo, en la práctica han aparecido problemas a la hora de disei'íar procesos de subasta que tengan en cuenta distintos objetivos a la vez. Por lo general, las principales opciones para los gobiernos descansan en procesos competitivos en los que se tengan en cuenta características y compromisos de la empresa compradora más allá del precio. o en los que se exijan ciertos requisitos previos a la propuesta de precio (aunque posteriormente el precio sea la variable de decisión).

\paragraph{Venta negociada}
Se elige al comprador y se negocian las condiciones de la compraventa. Es un método poco transparente que tiene que utilizarse cuando solamente existe un inversor idóneo para asumir la compra. En estos casos, al no existir un procedimiento de subasta que pem1 ita fijar el precio, la entidad privatizadora deberá ser especialmente cautelosa a la hora de garantizar que la empresa no se vende a un precio excesivamente conservador.\footnote{%
    Un precio de venta excesivamente bajo podría. además. entrar en conflicto con el Derecho Comunitario. si pudiera considerarse que el estado está indirectamente subsidiando a los compradores de la empresa privatizada (OCDE (2009). Privatisation in the 2 /si century • Recen/ expl!riencies of OL'('IJ coumr ies).
}

\subsubsection{Venta mixta}
En muchas ocasiones, la venta directa y la OPV se combinan, de forma que se vende parte de la empresa a un primer inversor estratégico, que va a realizar los ajustes necesarios para aumentar su rentabilidad y hacerla más atractiva, para posteriormente realizar una venta pública del resto de acciones que posea el sector público. Esta forma de privatización combinaría parte de las ventajas de los dos métodos explicados, en la medida en que permite elegir al primer inversor sin renunciar a la transparencia y accesibilidad de la OPV.


\subsection{Auditoría y control: ¿Cómo se verifica el resultado dela operación?}

En la mayoría de los países. el control ex post del proceso de privatización se realiza mediante auditoria y presentación ante el parlamento. Por lo general, la presentación ante el parlamento (y, por ende. ante el público en general) se considera una buena práctica porque aumenta la responsabilidad de los agentes implicados en el proceso. La institución encargada de esta auditoria debería. idealmente, estar presente en todo el proceso de privatización: reestructuración, valoración, elección de métodos, elección del vendedor. etc. y monitorear todas las transacciones.\footnote{OCDE (2009).
}



\clearpage
\section{Comparaciones internacionales}
\puntosclave{
    Las empresas públicas han aumentado su peso en la economía mundial en los últimos 20 años como consecuencia del auge de las economías emergentes. Llegan a suponer un 5\%, del empleo a nivel global y son más relevantes en los países en desarrollo.

    Las empresas públicas operan principalmente en sectores caracterizados por la presencia de economías de red, como las telecomunicaciones, la energía o el transporte. Destaca también su presencia en el sector financiero.
    
    En los últimos años ha aumentado el número de empresas públicas multinacionales, lo que ha acentuado la necesidad de diseñar un marco regulatorio que garantice un terreno de juego nivelado a nivel global.
}

Las empresas públicas juegan un papel clave en la economía mundial y en su crecimiento y desarrollo. En el año 2018, sus activos a nivel mundial tenían un valor de unos 45 trillones de dólares, en torno al 50\% del PIB global.\footnote{
\textcolor{blue}{\textbf{Pregunta:} ¿cuánto si excluimos las empresas chinas?}
}

Se ha estimado también que suponen un 20\% de la inversión y un 5\% del empleo en el mundo. Por países, se observa una gran disparidad entre economías en las que las empresas públicas tienen un mayor peso. frente a otras en las que su rol es más limitado: por lo general. su peso va a ser más importante en países en desa1Tollo. En cuanto a su descomposición sectorial, va a verse que son más prevalentes en sectores caracterizados por la existencia de externalidades en red y monopolios naturales, como la electricidad, las telecomunicaciones, el transporte o la provisión de agua.\footnote{lnternational Finance Corporation, Sta/e owned e11/e1prises (2018).
} así como en el sector financiero.

En la actualidad, el mayor po1tfolio de empresas públicas en términos de valor de activos lo ostenta China (29,2 trillones de dólares en activos distribuidos en más de 50.000 empresas públicas) seguida de lejos por la India (339,5 billones de dólares), Corea del Sur ( 2 1 7,8 billones) e Italia (207,5 billones).\footnote{OCDE (2017).
}

Dentro de la OCDE. los países en los que el peso de las empresas públicas en el total del empleo no agrícola es mayor son Noruega (9.6%). Letonia (6,7%), Estonia (4,8%), Hungría (4,2%), Francia (3.5%) y Finlandia (3.5%). Teniendo en cuenta no sólo las empresas de titularidad mayoritariamente pública. sino incorporando también aquellas en las que el sector público ostente una proporción m inoritaria de las acciones, la ordenación permanece inalterada, salvo por la entrada de Grecia y Alemania en el top- 1 5 de la OCDE. por la mayor prevalencia en estos países de este tipo de intervención.

\textbf{GRÁFICOS: OCDE 2017 (3)}

En la Unión Europea. la propiedad pública en varios sectores de la economía es particulannente extensa en algunos de los estados miembros que se adhi1ieron en 2004, como Polonia. Croacia. Rumania y Eslovenia. Sin embargo, las empresas públicas también tienen un lugar destacado en algunos Estados miembros de la U E 1 5, como Francia, Italia y Suecia. En l a mayoría de los Estados m iembros. la tendencia en el tiempo apunta a una reducción gradual en la propiedad pública. En términos de valor de mercado y empleo, las empresas públicas son particularmente relevantes en Finlandia, Eslovenia y Francia. y en menor medida en Bélgica y Letonia.\footnote{Comisión Europea (20 1 6).
}

España es el país de la Unión Europea con menor peso de las empresas públicas tanto en términos de valor de mercado sobre PIB como en términos de empleo. Destaca también a la baja Portugal. cuyo peso de las empresas de titularidad pública es similar al de España (aunque es sustancialmente más elevado en términos de empleo) y comparable al de Reino Unido. Ligeramente por encima. aunque en niveles reducidos, destaca también el bajo peso de las empresas públicas en Alemania y Países Bajos. lo que contrasta con su mayor relevancia en países comparables como Francia o Italia.

\textbf{GRÁFICO (1)}






\subsection{Composición sectorial}
De acuerdo con el estudio de la OCDE The Size and Sectoral Distribution uf State-Owned E11te1prises (201 7) las empresas públicas se encuentran altamente concentradas en sectores con presencia de economías de red. Conjuntamente, telecomunicaciones. electricidad y gas, transporte y otro tjpo de servicios básicos ( como los servicios postales) suponen la mitad del total de las empresas púhlicas en términos de calor de mercado, y hasta a un 70% en términos de empleo.

Entre estos sectores, la electricidad y el gas ostentarían un mayor peso, suponiendo el 2 1 o/o de las empresas púhlicas en términos de valor y un 1 0% en términos de empleo. A nivel nacional, las mayores empresas públicas de electricidad y gas (excluyendo China) se encuentran en India. Italia, Corea, Noruega y Francia, suponiendo en total un 66% del valor de mercado de la muestra analizada.

Más allá de los sectores de red, el sector financiero supone un 8% del empleo total y un 26% del total de las empresas públicas en términos de valor de mercado. de modo que. considerado individualmente, sería el que tendría mayor peso. En términos de empleo, los sectores con mayor peso serían el del transporte y otros servicios básicos, que conjuntamente suponen un 57% del empleo total.

\textbf{GRÁFICO (2): COMPOSICIÓN SECTORIAL EMPRESAS PÚBLICAS}

Considerando China, cabe destacar que las empresas públicas chinas difieren de forma relevante en cuanto a composición sectorial con respecto al resto del mundo. La mayor proporción de empresas públicas en China en términos de valor se encuentra en el sector financiero (que supone en torno a un 60% del valor de mercado total). seguido de otros sectores primarios (9%) y del sector del transporte, las manufacturas y la electricidad y el gas, que individualmente suponen en torno a un 5% del valor de mercado total. En cambio. los sectores primarios serían el mayor empleador. suponiendo un 29% del empleo total de las empresas públicas, seguido de las manufacturas, el transporte y. finalmente, el sector financiero.

\subsubsection{Empresas públicas en el sector energético}
Como ya se ha destacado. las empresas públicas juegan un rol fundamental en los sectores de red. Dentro de estos, merece mención especial el sector energético. que. especialmente en África y Asia. se encuentra dominado por empresas públicas. 

En concrew, las empresas públicas tienden a dominar los segmentos de la transmisión y distribución de electricidad. que tendrían cmacteristicas de monopolio natural. En cambio. el sector de la generación eléctrica permitiría una mayor participación de agentes privados, si bien las empresas públicas son un actor fundamental a nivel global incluso en ese segmento.

En los países desarrollados, la evidencia empírica es poco concluyente en cuanto a la capacidad de ciertas reformas en estos sectores. incluida la privatización. de generar ganancias de eficiencia. En los países menos avanzados y en desarrollo. los esfuerzos liberalizadorcs de los gobiernos han dado lugar también a resultados ambiguos. En este sentido. el acceso a la electricidad sigue siendo un reto en muchos países menos avanzados ( 840 millones de personas viven sin electricidad en África) y, si bien los actores privados podrían aumentar la capacidad de generación. por el momento la expansión de redes y del acceso ha ido de la mano de proyectos de infraestructuras promovidos por empresas públicas. que juegan. por tanto. un rol fundamental. 

Un problema común en estos casos es la dificultad de conseguir recuperar los costes a través de las tarifas. En los países en desarrollo y en los menos avanzados, es común que las empresas públicas del sector eléctrico realicen grandes inversiones en infraestructuras que posteriormente no se pueden recuperar a través de las tarifas, lo que reduce su capacidad de inversión y deteriora su situación financiera, además de generar distorsiones en la asignación de recursos. En todo caso, la necesidad de garantizar el acceso universal a la electricidad, que juega también un papel clave en el proceso de desarrollo de los países. ha hecho que estas prácticas sigan estando comúnmente extendidas.

\textbf{GRÁFICO (3): EMPRESAS PÚBLICAS SECTOR ELÉCTRICO}

Por lo que respecta a la producción de gas y de petróleo, muchos países productores crearon compañías estatales para explotar estos recursos y obtener graneles ganancias para el estado. Sin embargo, estas empresas públicas han tenido tradicionalmente menores beneficios que sus análogas en el sector privado. en parte como consecuenci a ele la voluntad del gobierno de ofrecer gas y petróleo a precios subsidiados como un tipo de gasto social, y han tenido también a generar unos mayores niveles de empleo y a tener una menor productividad del trabajo.\footnote{%
    De acuerdo con el FMJ: Fiscal Alonilor. />oficies fo Support People l)uri11g the Covid-19 Pa11de111ic (2020). la productividad aparente del trabajo (medida como millones de dólares de producción por trabajador) asciende a u n promedio de 0.09 en el promedio de las empresas públicas que operan en el sector de la producción de petróleo, mientras que sería. en promedio, de 0,2 en sus análogas en el sector privado.
}

\subsubsection{Sector financiero}
Como se ha explicado. la presencia de empresas pC1blicas en el sector financiero y, en concreto. en el sector bancario. es especialmente relevante, sobre todo en las economías emergentes. En concreto. las empresas públicas controlan un 40% de los activos bancarios en las economí<1s BRIC (Brasil, Rusia, India y China) y en algunos países menos avanzados. En países desarrollados. este porcentaje tiende a ser m enor, si bien llega a ser del 30% en Alenrnnia y Portugal.

\textbf{GRÁFICO (4): SECTOR BANCARIO}

Si bien la presencia de bancos comerciales públicos se ha reducido significativamente desde los años 90 (como consecuencia de la liberalización financiera llevada a cabo a nivel global) sigue siendo relevante en algunas grandes economías. En los países desaiTOllados, las recapitalizaciones y nacionalizaciones en el sector bancario tras la crisis de 2008 aum·entaron esta presencia.

Tradicionalmente, la existencia de empresas públicas en el sector financiero se ha justificado como forma de garantizar el crédito en recesiones. Sin embargo, la efectividad de los bancos públicos a la hora de estabil izar la economía tiene límites. En los últimos 20 años, los préstamos de los bancos públicos han tendido a ser menos procíclicos que los de sus análogos en el sector privado. si bien no ha sido así en economías en desarrollo con elevados niveles de deuda pública.

Por ejemplo. el banco brasileiio de desarrollo, BNDES. redujo repentinamente sus niveles de crédito durante la crisis ele 20 1 4-20 1 6 al costarse los préstamos que recibía del sector público.\footnote{FMI (2020).
}

Finalmente. los bancos públicos tienden a tener en sus activos una mayor cantidad de deuda pública que sus homólogos en el sector privado. especialmente en economías emergentes con mayores vulnerabilidades ligadas a la deuda pública. Esta situación puede generar el llamado circulo ,·icioso de deuda soberano y hancaria, suponiendo con ello una fuente de vulnerabilidad financiera.\footnote{%
    Esta exposición a l a deuda pública va l igada, normalmente, al otorgamiento de privilegios a estos bancos por parte del estado. Por ��jemplo, en la India, las garantías del gobierno a los bancos públicos permitieron expandir el crédito al estado durante la crisis financiera, al producirse un traslado de depósitos de bancos privados a públicos, ahora garantizados por el gobierno. Sin embargo, la calidad de los activos de estos bancos se deterioró. i ncrementando l a fragilidad financiera del sector y los riesgos asociados a los pasivoscontingentes del gobierno (FMI (2020)).
}

\subsection{Empresas públicas: actores globales}
Durante la última década. el porcentaje de empresas públicas sobre las 2.000 mayores empresas del mundo se ha duplicado, hasta alcanzar el 20\%. En este proceso. ha siclo clave el elevado crecimiento de las economías emergentes. especialmente de China. donde las empresas públicas siguen jugando un importante rol en la economía nacional. Esta expansión también se ha reflejado en transacciones intemacionales, destacando su rol en los mercados de deuda.\footnote{%
    La deuda de las mayores empresas públicas del mundo es en la actualidad de 7,4 billones de dólares. frente a 1 ,4 billones en el año 2000. Así, estas empresas suponen un tercio de la deuda en divisa de los países emergentes. jugando un rol esencial en los mercados de deuda corporativa.
} Destaca también que la propiedad de estas empresas es mixta en un 60\% de los casos, es decir. son empresas de titularidad pública y privada.\footnote{%
    Esta prominencia de la titularidad mixta tiene su origen en las estrategias europeas de privatización de los alios 80, en los que los gobiernos. en muchos casos. decidieron reservarse parte del control ele la empresa, especialmente en sectores estratégicos. Esta estrategia ha sido aplicada también en los mercados emergentes.
}

\textbf{GRÁFICO (5): MAYORES EMPRESAS DEL MUNDO}
Además, muchas de estas grandes empresas públicas son también multinacionales, es decir, controlan activos y realizan actividades en otros países. Las multinacionales públicas se han expandido en los últimos años. teniendo la mayoría de ellas su origen en China, en miembros de la Unión Europea, y en otros países como India o Malasia. Este proceso se explica por la necesidad de aumentar su rentabilidad, acceder a recursos naturales, obtener conocimiento tecnológico o incluso. de acuerdo con algunos autores, por razones políticas.\footnote{Cuervo-Cazurra et al. (2014).
}

El aumento del rol internacional de las empresas públicas genera importantes retos de gobernanza, en la medida en que pueden introducir distorsiones competitivas en los mercados nacionales y extranjeros. En este sentido. es común que los gobiernos den apoyo económico a las empresas públicas como compensación por la persecución de determinados objetivos de política económica, pudiendo este apoyo tomar la forma de subsidios, transferencias de capital, financiación en términos favorables, rescates en caso de quiebra o disposiciones regulatorias especiales. En caso de que estas ventajas superen el coste de cumplir con los objetivos a nivel público que tenga la empresa pública, ésta va a disponer de una mejor posición competitiva en los mercados internacionales.

Por todo ello, deviene especialmente importante el disef10 de un leve/ playing/leld para que las empresas públicas que operan internacionalmente compitan en igualdad, es decir, en neutralidad competitiva.\footnote{%
    La neutralidad competitiva puede definirse como la situación en la que ninguna entidad que opera en el mercado esta sujeta a ventajas o desventajas indebidas. Estas consideraciones no se aplican tan solo a las empresas públicas. sino que pueden aplicarse a entidades sin animo de lucro que operen en el mercado, o a empresas privadas que por alguna razón reciban apoyo público.
} En este sentido, la Unión Europea y Australia habrían presentado los mayores avances en términos de neutralidad competitiva; por ejemplo. Australia obliga a sus empresas públicas a compensar, mediante pagos al sector público, las ventajas regulatorias o financieras que puedan tener. Además, distintas instituciones multilaterales han establecido m arcos regulatorios a adoptar para crear un leve/ playing jield a nivel internacional, destacando, por parte de la OCDE, las Guidelines on Corporate Govemance ofSta/e-Owned emerprises (2015).\footnote{%
    El documento supone una reedición de una primera versión publicada en 201 5. e incorpora una serie de buenas prácticas legales y regulatorias. que incluyen la profesionalización de la gestión de las empresas públicas, una mayor transparencia y buenas prácticas en cuanto a la financiación de las empresas públicas. entre otros.
}

En los próximos años va a ser esencial el desarrollo de un m arco multilateral que garantice la competencia efectiva de las empresas públicas. Este proceso va a requerir de una voluntad coordinada de los estados y de las instituciones multilaterales, así como de un importante trabajo técnico que permita resolver muchos de los problemas de información que aparecen a la hora de regular de forma efectiva las empresas públicas.



\clearpage
\section{Evidencia empírica}
\puntosclave{
    La existencia de argumentos teóricos tanto a favor como en contra de la ineficiencia relativa de la empresa pública con respecto a la privada hace necesaria la realización de análisis empiricos.
    
    Las diferencias metodológicas entre los distintos estudios empíricos dificultan la existencia de un consenso en la materia. En cualquier caso, la mayoría de estudios apuntan a una mayor eficiencia relativa de la empresa privada cuando opera en sectores altamente competitivos, lo que obliga a poner el foco en asegurar la existencia de un marco regulatorio que promueva la competencia efectiva.
}

Como se ha visto a lo largo del tema, existen argumentos teóricos para defender la mayor eficiencia de la empresa pública en determinados contextos. a la vez que muchos otros vendrían a explicar su potencial ineficiencia relativa. Esta multiplicidad de argumentos teóricos ha hecho necesario realizar análisis empíricos con el fin de responder a la siguiente pregunta: ¿es la empresa pública. por lo general, menos eficiente que la privada?

En este apartado se va a hacer. en primer lugar. una revisión general de la evidencia empírica disponible para ver qué problemas presenta el análisis, así como el consenso en la materia. A continuación. se van a analizar de forma más pormenorizada algunos estudios específicos que se han considerado de mayor interés.

\subsection{Visión general}
La eficiencia de la empresa pública en comparación con la privada ha sido analizada por múltiples estudios que han llegado, en ocasiones. a conclusiones dispares. siendo dificil extraer un consenso general. En este sentido, algunos estudios sostienen la mayor ineficiencia de la empresa pública, mientras que otros llegarían a la conclusión de que no existe suficiente evidencia empírica en favor de esta tesis. Esta diversidad de los resultados se explica por las diferencias metodológicas entre los distintos estudios. que se proyectan, entre otros. sobre:
• Las definiciones de eficiencia utilizadas, que pueden ir desde la rentabilidad sobre el capital o sobre los activos (ROE y ROA) hasta medidas de eficiencia estrictamente productiva o eficiencia en costes.\footnote{%
    Como la productividad aparente del trabajo de cada empresa o la productividad global como relación entre el valor añadido y los inputs (capital y trabajo) utilizados en el proceso productivo.
}
• Los factores por los que se controla al realizar las comparaciones, que puedyn ser muy dispares y de distinta naturaleza (ej: sector regulado o no regulado, sector productivo específico. control por país. control de si estamos ante un país en desarrollo o no. etc.). Además. el planteamiento de los modelos permite distinguir entre dos grandes bloques de estudios: ( 1 ) los que comparan la eficiencia de empresas públicas y privadas que operan a la vez en un mismo sector y (2) los que comparan la eficiencia relativa de una misma empresa antes y después de ser privatizada. A pesar de la d i sparidad en las conclusiones, se pueden extraer unas cuantas conclusiones generales:

• Si l a empresa pública opera en sectores sometidos a un .-1lto grado de 1.:ompetencia (como pueda ser el sector manufacturero) la gran mayoría de estudios son favorables a l a mayor eficiencia de l.-1 empresa privada, con escasas excepciones. Además. se destaca que la creación de condiciones com petitivas es esencial para aumentar l a eficiencia de empresas públicas y privadas, siendo, en consecuencia, más relevante favorecer la competencia que establecer una determinada titularidad pública o privada.
• En sectores sometidos a escasa competencia o en los que la empresa privada opera con grandes restricciones regulatorias (ej: electricidad, fenocarril, telecomunicaciones) l a evidencia empírica no encuentra una mayor eficiencia d e l a empresa privada.\footnote{%
    Siguiendo a Hernández de Cos (2004) destacarían en estos sectores los estudios realizados por Shepherd (1966); Mann (1970); Yunker (1975): Spann (1977): Di Lorenzo y Robínson (1982): Edison Electric I.nstitute (1985): Atkinson y Halvorsen (1986); I-lolmes (1990): Hjalmarsson y Veiderpass (1991).
}
• En los países en desarrollo. caracterizados por tener un menor grado de competencia y de desarrollo de los mercados de capitales, no puede concluirse que las empresas públicas sean menos eficientes que las privadas, tanto desde un punto de vista productivo como asignativo. En cualquier caso, los datos en esos casos son escasos y de dudosa fiabilidad.
• Los estudios basados en el método (2). es decir, comparando empresas antes y después de su privatización. suelen haIJai, que el cambio en la titularidad genera mejoras sustanciales en la eficiencia. En cualquier caso, en muchas ocasiones las privatizaciones en m a��a van acompaliadas de reformas regulatorias que hacen d i fic i l atribuir causalmente los aumentos de la eficiencia al cambio de titularidad y no al aumento de la competencia derivado de las refonnas.\footnote{%
    Estos estudios han sido, tradicionalmente, menos numerosos. Muchos de ellos se centran en los países de Europa del Este, destacando el meta-análisis realizado por Djankov y Murrell (2002), que encuentra que la privatización suele conllevar un mayor nivel de reestructuración dentro de la empresa, con efectos positivos sobre la eficiencia técnica y asignativa. Sin embargo. existen pocos estudios de este tipo que se centren en procesos privatizadorcs ocurridos dentro de países desarrollados, y la mayoría de ellos se centran en Reino Unido. En este último caso. los resultados son más controve11idos y no apuntan necesariamente a una mayor eficiencia de la empresa privada (veáse Martin y Parker ( 1 995)).
}

\subsection{Análisis para el caso de la Un ión E u ropea: Com isión E u ropea (2016)}


\footnote{%
    Análisis contenido en el Informe de la Comisión Europea de julio de 20 1 6: ,'>'tate-nwned e11/e1prises inthe El/: Lessons learnt w1d waysforward in a pos1-crisis contexl.
}

En el estudio Stale-Owned Enlerprises in the EU: Lessons Leaml and Ways Forward in a PostCrisis Context (julio de 2 0 1 6 ) el Directorio General de Asuntos Económicos y Financieros de l a Comisión Europea realiza un análisis del rol que juegan l a s empresas públicas dentro d e l a Unión Europea. intentando estimar, entre otros. el impacto de la titularidad pública o privada en la eficiencia de las empresas que operan en los distintos subsectores del ferrocarri l y de las telecomunicaciones.


E l estudio analiza la eficiencia relati\'a de las empresas públicas )' pri\'adas que operan en e l sector de la energía ) del ferrocarril. distinguiendo 9 subsectores dentro del segmento energético y 3 subscctorcs dentro del segmento del fcrrocruTil, partiendo de la base de que algunos son monopolios naturales (y por tamo, sujetos a regulación) ) otros pueden ser relativamente competitivos. Se admite también la posibilidad de que ha) a empresas que operen a la vez en segmentos competitivos y no competitivos (conglomerados). Todo esto da lugar a la existencia de 1 4 subsectores en la est imación.\footnote{%
    9 en el segmento energético: generación de electricidad, transmisión de electricidad. distribución de electricidad. venta minorista )· mayorista de electricidad, transmisión de ga��. distribución de gas. venta minorista y mayorista de gas: 3 en el sector ferroviario: gestión de infraestructuras, transporte de mercancías y transporte de pasajeros: 1 conglomerados: uno en el sector energético y uno en el sector ferroviario.
}

Con datos de 2013. las empresas públicas en el sector de la energía en la UE supusieron hasta el 60% del total de las wntas del sector. teniendo, por tanto. una presencia muy destacada. aunque dispar entre estados miembros (por ejemplo. en Espai'la y Portugal su rol sería más reducido). En el sector ferroviario. las empresas públicas suponen la gran mayoría del volumen de ventas en ca�� todos lo�� estados miembros. a excepción de Estonia y República Checa. En los segmentos no competitivos. las empresas públicas suelen ser dominantes. m ientras que en los segmentos competitivos su peso tiende a ser menor y suelen operar las llamadas minority-SOEs (esto es, empresils que cuentan con capitHI público en participación m inoritaria, es decir. que no se encuentran completamente sujetas al control estatal).

\subsubsection{Consideraciones previas}
\textcolor{red}{\textbf{OJO -> Lo anterior era una subsubsección: NO TIENE SENTIDO} (meterla, o bien en el primer bloque, o bien como nota a pie de página resulta más coherente). 

Históricamente. la mayoría de los países organizaron la provisión de energía y servicios feiToviarios con monopolios estatales integrados verticalmente. Desde los 90. se impulsan una serie de reformas para pcnnitir la entrada de competencia que consistieron. en el sector energético, en separar las actividades de transmisión (monopolio natural) de las de generación. introduciendo requisitos en términos de provisión. E n el sector ferroviario, la transposición de los distintos paquetes de l iberalización permitió introducir competencia en distimos segmentos que se separaron de la gestión de la infraestructura. En consecuencia. coexisten en ambos casos un segmento competitivo y otro no competitivo.

En los segmentos regulados (1 igados a las infraestructuras) las tarifas deben pennitir el acceso a la red a un precio justo, si bien cada estado miembro tiene cierta autonomía a la hora de diseñar la tarifa. En consecuencia. la configuración de la tarifa puede hacer que la Hctividad en cada subsector resulte más o menos rentable dependiendo del país. lo que obliga a tener que controlar por esa variable si se quiere llevar a cabo un análisis riguroso.

En los segmentos no regulados. si bien se supone que existe competencia, la intensidad de la misma puede variar también entre estados miembros. en función de la profundidad de las reformas llevadas a cabo. lo que va a impactru· en la rentabilidad de las empresas públicas y privadas que operen en esos segmentos (uno esperaría que. por lo general. cuanto menor sea el grado de competencia mayor sea la rentabilidad). Además, el establecimiento de ohligaciones ele servicio público, que difieren también entre estados miembros. hace que no se pueda comparar la rentabilidad de las empresas que operan en un mismo segmento pero en distintos países.}

\subsubsection{Metodología}
El estudio consistiría en la metodología ( 1) explicada en el apaitado V.I: se compara la eficiencia relativa de empresas públicas y privadas que coexisten a la vez en un mismo sector. Así. se plantea la siguiente '.egresión. que se estima para cada uno de los 14 subsectores del ferrocarril y la energía:

$$Y_n=\beta_0+\beta_1SOE\_MAJOR_{it}+\beta_2SOE\_MINOR_{it}+\beta_3SIZE_{it}+\alpha_c+\alpha_t-\varepsilon_{it}$$

Donde:
• La variable dependiente (Y) busca ser una medida de eficiencia de las empresas. Corno ya se ha explicado, comparar la eficiencia de empresas públicas y privadas en términos de rentabilidad puede ser problemático porq ue las empresas públicas pueden renunciar en cie11os casos a paite de la rentabilidad para conseguir otros objetivos de interés público. Por ello. la regresión se estima con cada una de las siguientes variables dependientes:
o Rentabilidad del capital (ROA)
o Ratio de gastos operativos. que se calculará como el peso de los gastos operativos sobre el volumen total de ventas. Este indicador podrá interpretarse como proxy de la eficiencia productiva.
o Ratio de costes laborales, que se calculará como el peso de los costes laborales sobre el volumen total de ventas. Servirá para determinar si las empresas públicas tienden a ser más inten sivas en trabaj o (lo que implicaría algún grado de ineficiencia productiva).
o Rati o de capital sobre activos. que se calcularía como los recurs_os propios de la empresa divididos entre el total del activo: busca determinar si la titularidad de la empresa tiende a afectar a su estructura del capital.
o Ratio de inversiones, que se calculará como la inversión en activo fijo realizada en cada periodo entre el activo fijo total. Buscará determinar si las empresas públicas tienen más o menos propensión a invertir que las privadas.

• Como variables de control van a introducirse el país (ac ), el periodo (at ) y el tamaño de la empresa (SIZEit ), siendo todos estos factores que. como ya se ha explicado, influyen en la rentabilidad de las empresas y podrían sesgar los estimadores.

• Las variables independientes de interés serían dos variables dummy: SOE_MAJORit tomará el valor de 1 cuando la empresa sea de titularidad mayoritariamente pública y O en el caso contrario; SOE_MJNORít tomará el valor de l cuando la empresa esté partici pada por el sector público pero de forma minoritaria. y O en el caso contrario. De este modo, los coeficientes /31 y /32 van a indicarnos si las empresas públicas o participadas por el sector público de forma minoritaria tienden a obtener mejores o peores resultados medidos en términos de las variables dependientes.

De este modo, con datos de sección cruzada se estima una regresión para cada uno de los 14 subsectores de la energía y del ferrocarril y para cada una de las 5 variables dependientes explicadas.

\subsubsection{Resultados}
• En los segmentos regulados del sector energético se encuentra que: en la transmisión de electricidad la eficiencia de las empresas públicas es menor. mientras que ocurre lo contrario en la distribución de gas. En los demás casos. las diferencias no son significativas.
• En los segmentos regulados del sector fetTo viario no se encuentran diferencias signifi cativas entre la eficiencia de las empresas públicas y privadas. si bien las últimas tienen unos mayores costes laborales.
• En los segmentos no regulados tanto del sector energético como del ferroviario, se encuentra que por lo general las empresa�� públicas son menos rentables y menos eficientes en términos ele costes. presentando también unos mayores costes laborales.
• En términos de deuda, no se aprecian difere ncias significativas en el nivel de endeudamiento de empresas públicas y privadas en la mayoría de los subsectores: en alguno�� subsectores específicos. el endeudamiento de las empresas públicas es menor.
• En términos de inversión. no se encuentran diferencias significativas.
• En cuanto a las empresas con una participación minoritaria de capital público, se encuentra que su comportamiento tiende a semejarse más al de las empresas privadas en la gran mayoría de los casos. lo que confirmaría la tesis de que, a menor control público, menor impacto de los factores que lle van a una mayor ineficiencia de las empresas públicas.

En conclusión. los resultados están en línea con la literatura disponible. en el sentido de que. en los segmentos regulados (es decir, en los subsectores en los que 110 hay competencia efectiva) 110 hay diferencias significativas entre las empresas públicas y privadas. En los segmentos sometidos a competencia, las empresas privadas tienden a ser más eficientes, si bien no son resultados que se cumplan en todos y cada uno de los subsectores analizados. Tocio esto confirmaría la tesis de que la eficiencia puede ir más ligada al contexto competitivo y a la calidad ele la regulación que a la titularidad de la empresa.

Finalmente, la mayor ratio de costes laborales de las empresas públicas es un resultado robusto, que se cumple en todos los subsec tores y que sería consecuencia del mayor poder de los sindicatos en las empresas públicas y de la mayor generosidad de los salarios. así como de los menores ajustes de plantilla.

\subsection{España (HERNÁNDEZ DE COS, 2004)}
En este estudio de Banco de Espaíla se pretende analizar los cambios en la eficiencia de las empresas al pasar de ser públicas a privadas, para ofrecer evidencia sobre la hipótesis de que la privatización de las empresas públicas mejora la eficiencia empresarial. El trabajo estudia una muestra de empresas españolas privatizadas en el periodo 1983-1 996 e integradas en el sector manufacturero. para las que se dispone de información en los años anteriores y posteriores a la privatización mediante la Central de Balance�� del Banco de Espaíla (CBBE). La muestra total consta de 33 empresas privatizadas.

Por tanto, se trata de un estudio que se centra en la primera oleada de privatizaciones en España. y que elige el sector de las manufacturas por considerarlo especialmente competitivo, tanto antes como después de las privatizaciones; de este modo, se pretenden evitar los problemas derivados de no saber identificar si los aumentos de la eficiencia se deben a la privatización o a la creación de un contexto más competitivo

\subsubsection{Metodología}
El estudio se basa en la metodología (2) descrita en el apa1tado V.I: se compara el mismo bloque
de empresas antes y después de ser privatizada. Para solucionar los problemas ligados al hecho
de que las empresas públicas pueden ser menos rentables por el hecho de perseguir otros objetivos
(y no por st::r menos productivas) se intenta estimar el impacto de la privatización sobre la
eficiencia productiva, entendida como la productividad total de los factores. Así. se asume que la
tecnología de producción toma la forma de la siguiente función Cobb-Douglas: 

$$Y_{it}=AK_{it}^\alpha L_{it}^\alpha$$

De esta forma, se admite cualquier tipo de rendimientos a escala (y por ende, cualquier tipo de estructura competitiva). A partir de aquí, se loglineariza la expresión y se estima la siguiente regresión, con el objetivo de ver como la privatización afecta a la productividad (medida como su impacto en la producción controlando por la cantidad de factores utilizada):

$$ll Yít = Yo + (]kit + a lit + V¡ + Vt + Y1PRIVKPUit + L ({im Xmit + E¡t m=l$$

Donde:
• La variable dependiente (Y¡t ) es la producción en términos reales loglinearizada.
• kit y la representan la cantidad de capital y de trabajo incorporados a la producción, también en logaritmos. Como se observa, los coeficientes a y p se consideran comunes a toda la población de empresas, de forma que no se admite que puedan variar a nivel individual. Esta restricción parte de asumir que los rendimientos a escala van a ser los mismos en las distintas empresas (supuesto restrictivo que tiene que tomarse por el reducido tamaño de la muestra).
• V¡ y vt son coeficientes fijos para cada empresa y para cada momento en el tiempo, respectivamente. Con ellos, se busca capturar la mayor o menor productividad estructural de cada empresa (antes y después de la privatización) y los efectos comunes en la productividad derivados ele factores ambientales que cambien entre periodos.
• PRIVKPUit va a ser una variable dummy que tome el valor de I cuando la empresa ha sido privatizada y O cuando aún sea pública. Por tanto. el coeficiente y1 va a ser el relevante a efectos de determinar si la privatización impacta positivamente en la eficiencia productiva de la empresa. 
• ¿��= 1 cpm Xm i t son las variables de control introducidas en la regresión, que en concreto van a ser: ( 1 ) grado de concentración sectorial; (2) penetración sectorial de las importaciones; (3) nivel de endeudamiento y (4) proporción de empleo temporal sobre el total de empleo de la empresa. Nótese que las variables ( 1 ) y (2) pretenden capturar elgrado de competencia del sector en el que opera cada empresa. algo especialmente relevante teniendo en cuenta que en los aíios analizados ( 1983-1 996) se entra en la CEE, que aumentó la competencia exterior de los productos con el consecuente impacto en la product ividad, que podría confundirse con efectos derivados de la privatización. Las variables (3) y (4) permiten tener en cuenta que la estructura financiera de la empresa. así como la temporalidad en t'I empleo, puedan afectar a la productividad.

\subsubsection{Resultados}
Se comprueba que la competencia tiene un impacto relevante en la productividad. siendo que el coeficiente de la variable de control ( 1) es negativo y significativo, y el de la variable de control (2) es positivo y signi ficativo. La variable de endeudamiento (3) muestra también un coeficiente negativo y significativo.
• Como es lógico. los factores prodqctivos (capital y trabajo) tienen un impacto positivo en la producción.
• La variable de privatización (PRIVKPU) presenta un coeficiente positivo y significativo. lo que indica que la privatización tuvo. en promedio. un efecto positivo sobre la productividad de las empresas. El signo de este coeficiente. además, es robusto ante otras especificaciones del modelo (que. entre otros, introducen retardos o miden la productividad en términos de rentab ilidad).
• Además. con otros contrastes se encuentra un efecto positivo de la privatización so bre el stock de capital real. la ratio capital/trabajo y la remuneración real por trabajador. Se observa también un efecto negativo sobre le empleo. el endeudamiento y la cuota de mercado.

En conclusión .. el estudio viene a demostrar que. para el periodo y la muestra tomada en consideración. la privatización aumenta la eficiencia productiva incluso controlando por el contexto competitivo (que también confirma tener un impacto significativo en la producción). En todo caso, no deben olvidarse las limitaciones del trabajo realizado, relacionados especialmente con el tamaño de la muestra disponible, así como las circunstancias particulares en las que se realizó la privatización de las empresas analizadas (en el sentido de que. otro proceso de privatización con otras normas y en otros sectores podría dar lugar a resultados distintos.





\clearpage
\section{Empresa Pública: ¿Por qué existe?}
\puntosclave{
    
}

Este tema responde, en última instancia, a una sola pregunta —¿cuándo debe el Estado producir directamente bienes y servicios, en lugar de limitarse a regular o subvencionar a quien los produce?—, pero la respondemos en tres tiempos, no de una vez: primero reconstruyendo por qué, históricamente y en teoría, se decidió que sí debía hacerlo (bloque 1); después exponiendo por qué esa decisión fue objeto de una crítica teórica e histórica de gran calado que llevó a revertirla en buena parte del mundo (bloque 2); y cerrando con la respuesta más matizada y actual, que ni confirma la tesis inicial ni la antítesis que la sustituyó, sino que identifica las condiciones concretas —de competencia, de regulación, de tipo de fallo de mercado— bajo las cuales una u otra opción resulta preferible (bloque 3). Antes de empezar, conviene fijar una advertencia de enfoque: este no es un tema sobre el sector público empresarial español. Es un tema de Economía Pública general, y España —el INI, la SEPI, las privatizaciones de los noventa— aparecerá en él exactamente con el mismo estatus que el Reino Unido, Italia, Francia o México: como una ilustración más de un fenómeno mundial, nunca como su centro de gravedad.

\subsection{Desarrollo histórico: Construcción del tejido público empresarial}
Antes de narrar esta historia, conviene una precisión conceptual que evita una confusión habitual: este relato no es el de la planificación centralizada de las economías del bloque soviético, donde la propiedad estatal de los medios de producción sustituía por completo al mercado como mecanismo de asignación. El objeto de este tema —y de la propia definición del SEC 2010 que veremos en el subapartado 1.3— es la empresa pública que opera dentro de una economía de mercado, compitiendo o coexistiendo con agentes privados, sujeta a un balance y a una cuenta de resultados, aunque bajo control estatal. Es un fenómeno propio de las economías mixtas, no de la planificación integral, y esa distinción es la que explica por qué el relato que sigue transcurre, casi en su totalidad, en Europa Occidental y en el mundo en desarrollo no alineado, y no en la URSS o sus satélites.

[Probable] La empresa de titularidad estatal tiene antecedentes muy anteriores al siglo XX —arsenales reales, casas de moneda, monopolios postales, compañías comerciales con privilegio real como la Compañía Holandesa de las Indias Orientales—, pero el fenómeno que interesa a este tema, el del tejido público empresarial moderno operando en sectores industriales y de servicios en competencia potencial con el sector privado, tiene dos genealogías distintas que conviene narrar por separado antes de que confluyan, porque responden a lógicas históricas diferentes y porque comprenderlas por separado es lo que permite, más adelante, entender por qué la crítica neoliberal de los años setenta y ochenta las trató como si fueran un único fenómeno cuando en realidad no lo eran.

La primera genealogía es la europea de posguerra, y dentro de ella conviven, a su vez, dos motivaciones distintas que rara vez se distinguen con la precisión que merecen. En el Reino Unido, la ola de nacionalizaciones del primer Gobierno laborista de Clement Attlee (1945-1951) —Banco de Inglaterra en 1946, la industria del carbón mediante la creación del National Coal Board en 1947, los ferrocarriles con la creación de British Railways en 1948, el gas, la electricidad, y más tarde el acero— fue, ante todo, un proyecto ideológico: respondía al compromiso programático del Partido Laborista, formalizado en la Cláusula IV de sus estatutos de 1918, de asegurar "la propiedad común de los medios de producción, distribución e intercambio". [Seguro] El fundamento doctrinal más elaborado de esta postura llegaría algo después, con Anthony Crosland (The Future of Socialism, 1956), quien reformuló el argumento en términos que resultan centrales para todo lo que sigue en este tema: la propiedad pública no era, para Crosland, un fin en sí mismo, sino uno más de los instrumentos disponibles para alcanzar fines socialistas —igualdad, pleno empleo, control democrático de la economía— dentro de una economía todavía fundamentalmente de mercado, lo que anticipaba ya, en 1956, la idea de que la titularidad y el objetivo perseguido son cuestiones lógicamente separables, precisamente la distinción que va a vertebrar el bloque 3 de este tema setenta años después.

[Probable] En Francia, la primera ola de nacionalizaciones de 1945-1946 —Renault, la energía con la creación de Électricité de France y Gaz de France, buena parte de la banca y los seguros— combinó ese mismo impulso ideológico con un motivo adicional que el caso británico no tuvo: el punitivo. La nacionalización de Renault en 1945 se ejecutó, formalmente, como sanción contra su propietario, Louis Renault, acusado de colaboración con la ocupación alemana —un origen que conviene conocer porque desmiente la idea de que toda nacionalización responde a un cálculo puramente económico de fallo de mercado; a veces responde a un cálculo de legitimidad política en un contexto de posguerra muy concreto—. Francia repetiría el patrón, ya sin componente punitivo, cuatro décadas después: la segunda ola francesa, bajo François Mitterrand en 1981-1982, nacionalizó los principales grupos bancarios e industriales del país, un episodio tardío que demuestra que la genealogía europea de posguerra no se cerró, como suele simplificarse, en los años cuarenta, sino que tuvo una réplica tardía justo en el umbral cronológico de la contrarreforma neoliberal que estudiaremos en el bloque 2.

[Seguro] Italia ofrece un origen genealógico distinto, y más interesante analíticamente porque no nace de una ideología socialista sino de una urgencia financiera: el IRI (Istituto per la Ricostruzione Industriale), creado en 1933 bajo el régimen fascista de Mussolini, a instancias del banquero Alberto Beneduce, no como proyecto doctrinal sino como rescate financiero de emergencia: el Estado italiano se vio obligado a asumir la propiedad de la banca mixta italiana —y, con ella, de las participaciones industriales que esa banca tenía en cartera— tras la crisis bancaria derivada de la Gran Depresión, cuando esos bancos resultaron insolventes y arrastraban consigo a buena parte de la industria pesada del país. [Probable] Lo notable del caso italiano es que esa titularidad, nacida como solución de urgencia y pensada en su momento como transitoria, se convirtió en permanente y se transformó, tras la caída del fascismo, en el conglomerado público más extenso de Europa Occidental durante décadas —una circunstancia que conviene tener presente porque revela una vía de creación de empresa pública distinta tanto de la ideológica como de la estrictamente teórica: el rescate financiero que nunca se deshace—. España, con la creación del INI (Instituto Nacional de Industria) en 1941 bajo el franquismo, siguió explícitamente el modelo institucional italiano —el propio INI se diseñó a imagen del IRI—, aunque con una motivación adicional propia del régimen autárquico español de posguerra: la sustitución de importaciones y la industrialización forzada de un país aislado internacionalmente, más que el rescate de un sistema bancario en quiebra.

La segunda genealogía es la de la descolonización y la soberanía sobre los recursos naturales, y tiene una lógica completamente distinta de la europea: no es la construcción de un tejido industrial nuevo, sino la nacionalización de activos extractivos ya existentes y previamente en manos de empresas multinacionales de países desarrollados. [Seguro] El episodio pionero, incluso anterior a la Segunda Guerra Mundial, es la expropiación petrolera mexicana de 1938, decretada por el presidente Lázaro Cárdenas frente a compañías petroleras británicas y estadounidenses, que dio origen a Petróleos Mexicanos (PEMEX) y que sigue conmemorándose en México como fecha de soberanía económica nacional. [Seguro] Un episodio posterior, políticamente mucho más dramático, es la nacionalización iraní del petróleo de 1951, decretada por el primer ministro Mohammad Mossadegh frente a la Anglo-Iranian Oil Company —episodio que desencadenaría, dos años después, el golpe de Estado de 1953 respaldado por los servicios de inteligencia británico y estadounidense, ilustrando hasta qué punto la nacionalización de activos extractivos podía convertirse en un asunto de geopolítica de primer orden y no solo de política económica doméstica—.

[Seguro] La teoría que mejor explica la lógica económica subyacente a esta segunda genealogía es el modelo del "pacto obsolescente" (obsolescing bargain) de Raymond Vernon (Sovereignty at Bay, 1971): el poder de negociación entre una multinacional extractiva y el Estado que alberga el yacimiento no es constante, sino que se desplaza sistemáticamente con el tiempo. En el momento inicial de la inversión, la multinacional tiene todo el poder de negociación —posee el capital, la tecnología y el know-how que el país anfitrión no tiene—, lo que le permite negociar condiciones muy favorables (concesiones a largo plazo, tipos impositivos reducidos). Pero, una vez realizada la inversión —una vez el pozo está perforado, la mina excavada, la infraestructura construida—, esos costes son ya hundidos (sunk costs): la multinacional no puede retirar la inversión sin perderla por completo, mientras que el país anfitrión ya ha aprendido a operar la infraestructura y ha dejado de depender tan estrechamente de la tecnología extranjera. El poder de negociación se desplaza, pues, hacia el Estado, que puede entonces renegociar condiciones más favorables —o, en el extremo, nacionalizar directamente— sin temer que el inversor abandone el país, precisamente porque ya no puede hacerlo sin pérdidas irrecuperables. [Probable] Este modelo se ha aplicado empíricamente a casos tan diversos como el petróleo saudí (con el acuerdo "50-50" entre Aramco y la monarquía saudí en 1950, antesala de la nacionalización plena posterior), la propia Anglo-Iranian de 1951, o la bauxita de Guyana nacionalizada a Alcan en 1971 —y culmina institucionalmente con la creación, en 1960, de la OPEP, que coordinó precisamente a los países productores que habían recorrido, cada uno a su ritmo, este mismo proceso de desplazamiento del poder de negociación frente a las grandes petroleras occidentales—.

[Probable] Esta doble genealogía —europea de posguerra por un lado, extractiva-descolonizadora por otro— confluye, hacia mediados de siglo, en un consenso doctrinal común que conviene nombrar antes de cerrar este recorrido histórico: el llamado consenso keynesiano de la economía mixta, bajo el cual la propiedad pública de sectores estratégicos se consideraba, en prácticamente todo el espectro político democrático occidental —no solo en la izquierda socialista—, un instrumento legítimo y hasta deseable de gestión macroeconómica, coherente con el papel activo que el Estado había asumido en la estabilización del ciclo tras la Gran Depresión y la Segunda Guerra Mundial. [Seguro] Este consenso, sin embargo, empieza a resquebrajarse a partir de la crisis del petróleo de 1973 y la subsiguiente estanflación, que erosionaron simultáneamente la capacidad fiscal de los Estados para sostener el tejido público empresarial y la confianza teórica en que la intervención directa fuera la respuesta correcta a los problemas económicos del momento —un diagnóstico de la crisis fiscal del Estado que James O'Connor formalizó, ya en pleno vórtice de la crisis, en su influyente The Fiscal Crisis of the State (1973)—. Es exactamente en ese punto de ruptura donde termina este subapartado histórico y donde arranca, de forma natural, el bloque 2 de este tema: la crítica teórica que explica por qué ese modelo, dominante durante tres décadas, entró en una crisis de legitimidad de la que no se ha recuperado del todo.

\subsection{Argumentos clásicos: ¿por qué debería existir?}

Cerrado el recorrido histórico, toca ahora el aparato teórico que legitimó doctrinalmente esa historia. Conviene distinguir, desde el principio, tres familias de argumentos —de eficiencia, de equidad y de estabilización—, porque cada una responde a un tipo de fallo o de objetivo distinto, y porque esa distinción es la que después nos va a permitir, en el bloque 3, separar con precisión qué críticas afectan a cuáles.

\subsubsection{Razones de eficiencia}

La familia de argumentos más antigua y mejor formalizada es la que justifica la empresa pública como respuesta a un **fallo de mercado que impide la existencia de una estructura plenamente competitiva**. Dentro de ella conviene distinguir dos vertientes con una lógica interna distinta.

\paragraph{Monopolio natural}
[Seguro] Un mercado presenta monopolio natural cuando la estructura de costes de la actividad —típicamente, costes fijos muy elevados y costes marginales decrecientes o muy reducidos, como ocurre en redes de transporte de electricidad, gas o ferrocarriles— hace que una única empresa pueda satisfacer toda la demanda del mercado a un coste total inferior al que resultaría de fragmentar esa misma producción entre varias empresas competidoras. La formalización moderna de esta condición —la llamada **subaditividad de costes**— se debe a Baumol, Panzar y Willig (*Contestable Markets and the Theory of Industry Structure*, 1982), que además introdujeron un matiz importante para el debate de este tema: un monopolio natural no siempre requiere regulación de precios si el mercado es "disputable" —es decir, si la amenaza creíble de entrada de nuevos competidores basta para disciplinar al monopolista incumbente, incluso sin que la entrada llegue a producirse—, lo que anticipa ya la idea, central en el bloque 3, de que la variable relevante puede no ser la titularidad de la empresa sino la contestabilidad del mercado en que opera.

[Probable] Frente a esta misma condición de monopolio natural, la teoría económica identifica tres respuestas de política pública posibles, no solo la de crear una empresa pública: la **regulación de precios** de un proveedor privado (la vía tradicionalmente dominante en Estados Unidos, con la Comisión Federal de Energía y comisiones estatales de servicios públicos regulando tarifas de empresas privadas, tradición cuya sistematización clásica se debe a Alfred Kahn, *The Economics of Regulation*, 1970-71), el **establecimiento de subsidios o impuestos** que aseguren la provisión también en los segmentos no rentables del mercado (el caso típico de subsidiar el servicio postal en zonas rurales), o directamente la **provisión pública** —la vía tradicionalmente dominante en Europa continental—. La elección entre estas tres vías no es indiferente: mientras la regulación mantiene la propiedad privada y sus incentivos asociados pero exige una capacidad técnica regulatoria considerable para no ser capturada por el propio regulado, la provisión pública resuelve el problema de la captura regulatoria a costa de introducir los problemas de agencia que analizaremos en detalle en el bloque 2.

\paragraph{Externalidades y sectores estratégicos}

[Probable] La segunda vertiente de eficiencia responde no a un problema de estructura de costes sino a la existencia de **externalidades** —positivas o negativas— que el mercado no internaliza correctamente por sí solo. El ejemplo canónico es el transporte colectivo urbano: un operador privado, al fijar su oferta según la demanda que capta directamente, no tiene en cuenta el beneficio social adicional que su servicio genera al descongestionar el tráfico general de la ciudad, un beneficio del que se apropian terceros que ni siquiera usan el servicio. La respuesta clásica de la teoría pigouviana —impuestos o subvenciones que corrijan el precio privado hasta igualarlo al coste o beneficio social— es una alternativa a la empresa pública, no su única solución posible; pero cuando el cálculo y la gestión de ese subsidio resultan administrativamente más costosos que la provisión pública directa, esta última puede resultar preferible en términos puramente de eficiencia.

Dentro de esta misma vertiente, conviene distinguir un caso particular que adquiere renovada relevancia en la actualidad: la existencia de **sectores estratégicos**, cuya provisión pública se justifica no tanto por una externalidad convencional cuanto por el riesgo de **dependencia exterior** y de estrangulamientos de suministro que podrían comprometer la productividad de largo plazo de toda la economía. [Suposición] Este argumento —que sustentó buena parte de las grandes inversiones públicas en infraestructura energética y de transporte en América Latina durante los años setenta, y en África y Asia durante los ochenta— tuvo, sin embargo, un balance histórico ambivalente: muchas de esas infraestructuras terminaron infrautilizadas y no llegaron a recuperar su coste de producción (Khan y Kumar, 1997), un dato que conviene retener porque anticipa ya, dentro del propio bloque "a favor", una de las críticas de eficiencia que se sistematizará con mucho mayor rigor en el bloque 2.

\subsubsection{Razones de equidad}

La segunda gran familia de argumentos no busca corregir una ineficiencia de mercado, sino perseguir objetivos distributivos que el mercado, aunque funcione perfectamente en términos de eficiencia asignativa, no garantiza por sí solo.

\paragraph{Servicios públicos esenciales y bienes de mérito}

[Seguro] A diferencia de la empresa privada, que maximiza beneficios sujeta a las condiciones de mercado, la empresa pública puede perseguir la **provisión universal** de determinados bienes y servicios a precios asequibles, con independencia de que esa provisión sea rentable en cada segmento concreto del mercado. El fundamento doctrinal más elaborado de esta idea procede del concepto de **bienes de mérito** (*merit goods*) de Richard Musgrave (*The Theory of Public Finance*, 1959): bienes cuyo consumo la sociedad considera deseable con independencia de las preferencias individuales reveladas por la demanda de mercado, lo que justifica que el Estado module su provisión por encima del nivel que el mercado generaría espontáneamente. El ejemplo empírico más citado es el de los servicios ferroviarios de cercanías en líneas no rentables: solo un puñado de países europeos, todos ellos con densidad de población muy elevada, consiguen recuperar el coste de provisión de estos servicios mediante el cobro directo a los usuarios (FMI, 2020), mientras que en la mayoría de los países las disparidades territoriales de rentabilidad son tan severas —en Francia, el 15% del coste total del servicio ferroviario se concentra en apenas el 2% de los usuarios— que exigen, o bien subsidios sostenidos y condonaciones de deuda recurrentes a la operadora pública, o bien mecanismos alternativos como obligaciones de servicio público impuestas a operadores privados.

\paragraph{Desigualdad horizontal: corrección de disparidades regionales}

[Probable] Una segunda vertiente del argumento de equidad, de naturaleza más territorial que sectorial, sostiene que la empresa pública puede utilizarse deliberadamente como instrumento de **cohesión regional**: en un contexto de cambio tecnológico que ha reforzado sistemáticamente las economías de aglomeración de las grandes ciudades a costa de las regiones periféricas y de menor densidad —acentuando procesos de despoblación—, la localización deliberada de empleo público cualificado en regiones desfavorecidas puede generar efectos de arrastre sobre la demanda local de bienes y servicios y favorecer, indirectamente, el surgimiento de iniciativa privada adicional. Los ejemplos comparados son numerosos y de perfil político muy distinto entre sí: la apertura de establecimientos públicos en el este de Alemania tras la reunificación; el traslado de parte de la BBC a Salford, cerca de Mánchester, en el Reino Unido; el traslado de dos tercios de las agencias estatales surcoreanas desde Seúl a ciudades de menor población; o el reciente plan de desarrollo de Baviera, que contempla el traslado de más de cincuenta entidades públicas a zonas rurales del estado. [Suposición] Conviene señalar, aunque el propio argumento rara vez lo hace explícito, que su eficacia se sostiene sobre un supuesto empírico no siempre verificado: que el efecto de arrastre sobre la economía local sea proporcional a la cualificación de los trabajadores trasladados, lo que convierte a esta política en potencialmente mucho más eficaz cuando se traslada un organismo regulador con personal técnico de alta cualificación que cuando se traslada una función meramente administrativa.

\subsubsection{Razones de estabilización}

[Seguro] La tercera familia de argumentos, de raíz explícitamente keynesiana, cobra fuerza tras la Segunda Guerra Mundial: si la inversión privada es intrínsecamente volátil y tiende a contraerse precisamente en los periodos de crisis —cuando más se necesitaría sostenerla—, disponer de un tejido empresarial público permite al Estado acometer inversión directa de forma anticíclica, sin depender de la disposición a invertir de agentes privados desanimados por el propio ciclo. [Probable] La evidencia contemporánea reciente ofrece cierto respaldo empírico a esta lógica: los datos del *Life in Transition Survey* muestran que el porcentaje de trabajadores con contrato indefinido es sistemáticamente mayor en las empresas públicas que en las privadas, y que estos trabajadores reportaron una probabilidad significativamente menor de perder su empleo durante la crisis financiera de 2008 —diferencias que persisten siendo estadísticamente significativas incluso controlando por edad, género y tamaño de la empresa—. Por el lado de la oferta de crédito, la evidencia sugiere que las entidades financieras de titularidad pública tienden a ser prestamistas más estables durante las crisis que sus homólogas privadas, mientras que, por el lado de los hogares, los trabajadores del sector público —beneficiados por esa mayor estabilidad de renta— muestran una mayor tasa de ahorro que sus homólogos privados de renta comparable, lo que en conjunto refuerza su capacidad de suavizar su propio consumo ante shocks adversos.

[Probable] Conviene cerrar este subapartado señalando, sin embargo, que este mismo efecto estabilizador tiene un reverso que la propia literatura reconoce: la estabilidad del empleo público, precisamente por no estar sujeta a los mismos mecanismos de ajuste que la empresa privada, es también uno de los canales por los que se cuela la ineficiencia relativa que analizaremos con rigor en el bloque 2 —existe, por tanto, un **trade-off estructural entre la función estabilizadora y la función de motor de eficiencia y crecimiento** de la empresa pública, que ningún diseño institucional ha conseguido eliminar por completo, y que conviene que traslades al tribunal como conclusión de cierre de todo este primer bloque: los mismos rasgos que hacen a la empresa pública un buen amortiguador del ciclo son, con frecuencia, los que la literatura de agencia señala como fuente de ineficiencia relativa en tiempos normales.


\subsection{Marco conceptual: qué es, formalmente, una empresa pública}

Cierro el bloque 1 con la pregunta que, en rigor, debería haber abierto el tema entero, y que sin embargo llega aquí a propósito: solo después de saber por qué y cuándo surgió este fenómeno tiene sentido preguntarse con precisión qué es, porque la propia historia que acabamos de recorrer —rescates bancarios italianos, nacionalizaciones punitivas francesas, expropiaciones petroleras mexicanas— muestra que "empresa pública" no ha significado siempre lo mismo, y que cualquier definición demasiado estrecha dejaría fuera a la mitad de los casos que hemos estudiado.

\subsubsection{Aproximación económica: un híbrido de dos lógicas}

[Probable] Antes de acudir a cualquier definición jurídico-estadística, conviene fijar la aproximación puramente económica al concepto, que la literatura clásica sobre gestión de empresas públicas —con la obra de referencia de Yair Aharoni, *The Evolution and Management of State Owned Enterprises* (1986), como punto de partida obligado— construye como un **híbrido de dos conjuntos de rasgos** que rara vez concurren por completo. Del lado del sector público, la empresa pública se caracteriza porque sus decisiones las toman agentes vinculados a la esfera política, cuyo criterio no está estrictamente ligado a la rentabilidad financiera; porque sus beneficios pertenecen, al menos parcialmente, al erario público; y porque se espera de ella una rendición de cuentas ante la sociedad en su conjunto, no solo ante sus accionistas. Del lado de la empresa, en cambio, se espera que sea financieramente viable a largo plazo, y que sus precios guarden al menos alguna relación con sus costes de producción —dos rasgos que la distinguen de un organismo administrativo ordinario, que no necesita cubrir sus costes con ingresos propios—. [Suposición] La utilidad de esta doble caracterización no es meramente descriptiva: explica por qué, en la práctica, resulta tan difícil trazar una frontera nítida entre empresa pública y organismo administrativo —muchas entidades cumplen los criterios del primer grupo con intensidad variable y los del segundo de forma solo parcial—, lo que hace inevitable acudir a un criterio más operativo y menos ambiguo para fines de clasificación estadística. Ese criterio es el del **control efectivo**, que desarrollo a continuación.

\subsubsection{Aproximación jurídico-estadística: el criterio de control del SEC 2010}

[Seguro] El Sistema Europeo de Cuentas 2010 (SEC 2010), marco estadístico de referencia de la Unión Europea, no define la empresa pública por la propiedad formal de su capital, sino por el **control efectivo** que el sector público ejerce sobre ella, y establece un catálogo de indicadores que deben valorarse en su conjunto —sin que sea necesario cumplir todos, ni suficiente con cumplir solo uno—: la titularidad pública de la mayoría de los derechos de voto; el control público del consejo de administración o del órgano ejecutivo; el control público del nombramiento o la revocación del personal directivo; el control público de los principales comités de la entidad; la posesión de una **acción de oro** que otorgue derecho de veto en operaciones clave, aunque la participación de capital sea minoritaria; la existencia de una **normativa especial** que condicione sustancialmente las decisiones empresariales de una entidad formalmente privada; el hecho de que la Administración represente el grueso de la demanda de la empresa; o que la Administración le preste recursos financieros en condiciones que exceden las de cualquier entidad crediticia ordinaria.

[Probable] Lo que el documento original no explicaba —y que conviene incorporar aquí, porque es la razón por la que este ejercicio de clasificación importa tanto en la práctica— es **por qué** esta frontera tiene consecuencias que van mucho más allá de lo terminológico. El propio SEC 2010 introduce, junto al criterio de control, un segundo filtro decisivo para la contabilidad nacional: el **test de productor de mercado**, que compara los ingresos por ventas de la entidad con sus costes de producción. Si esos ingresos cubren menos del 50% de los costes —el llamado *criterio del 50%*—, la entidad deja de clasificarse como "empresa pública" a efectos estadísticos y se reclasifica como parte del **sector Administraciones Públicas**, con una consecuencia nada trivial: su déficit y su deuda pasan a computar, íntegramente, en el déficit y la deuda pública del Estado a efectos de los criterios de convergencia europeos —el mismo umbral del 3% y el 60% del PIB que rige el Pacto de Estabilidad y Crecimiento—. [Suposición] Esta es, precisamente, la razón técnica por la que la clasificación sectorial de una empresa pública deficitaria puede convertirse en objeto de disputa política e incluso de ingeniería contable: mantener una entidad por encima del umbral del 50% de cobertura de costes —aunque sea mediante subidas de tarifa o reducciones artificiales del gasto declarado— permite mantenerla fuera del perímetro de las cuentas públicas a efectos de Maastricht, mientras que caer por debajo de ese umbral la convierte, de la noche a la mañana, en una fuente directa de déficit público computable. Este dato conecta directamente, y lo retomaremos así en el bloque 2, con una de las razones presupuestarias que con más frecuencia impulsan procesos de privatización: no siempre es el rendimiento económico real de la empresa lo que está en juego, sino su tratamiento estadístico a efectos de las reglas fiscales europeas.

\paragraph{Singularidad de la acción de oro y su choque con el Derecho europeo}

[Probable] De entre todos los criterios de control que acabamos de enumerar, la **acción de oro** merece un desarrollo propio, porque ilustra mejor que ningún otro la tensión entre la lógica de control estatal que legitimó históricamente la empresa pública y la lógica de integración de mercados que rige la Unión Europea desde los tratados fundacionales. [Seguro] Se trata de un instrumento surgido, precisamente, durante las privatizaciones europeas de los años ochenta —el momento histórico que abrirá el bloque 2—: al ceder el control accionarial de sus antiguas empresas públicas a inversores privados, varios gobiernos europeos se reservaron una participación simbólica, a menudo una sola acción, que sin embargo llevaba aparejado un derecho de veto sobre determinadas operaciones societarias —fusiones, cambios de objeto social, adquisiciones de participaciones significativas por terceros—, típicamente en sectores considerados estratégicos: energía, telecomunicaciones, defensa, transporte.

[Seguro] Este instrumento fue sometido a un intenso escrutinio por parte del Tribunal de Justicia de las Comunidades Europeas a comienzos de la década de 2000, en una serie de sentencias que conviene conocer en su conjunto porque revelan que el Tribunal no adoptó una postura uniforme, sino un test de proporcionalidad caso por caso. El **4 de junio de 2002**, el Tribunal resolvió simultáneamente tres asuntos —Comisión contra Portugal (C-367/98), Comisión contra Francia (C-483/99) y Comisión contra Bélgica (C-503/99)—, declarando **contrarios al Derecho comunitario** los regímenes de acción de oro portugués y francés, por constituir restricciones injustificadas a la libre circulación de capitales y a la libertad de establecimiento, mientras que **validó el régimen belga**, precisamente porque estaba diseñado con criterios más precisos, conocidos de antemano y sujetos a control judicial, y limitado a decisiones muy concretas relacionadas con la seguridad del abastecimiento energético. [Seguro] El **13 de mayo de 2003**, el Tribunal resolvió los asuntos paralelos contra España (C-463/00, referido a las acciones de oro del Estado en Repsol, Telefónica, Endesa, Argentaria y Tabacalera) y contra el Reino Unido (C-98/01, referido a BAA, el operador de los principales aeropuertos británicos), condenando a **ambos** Estados —un desenlace que conviene resaltar porque el Abogado General Ruiz-Jarabo Colomer, en sus conclusiones de febrero de ese mismo año, había propuesto al Tribunal desestimar el recurso contra España y solo estimarlo contra el Reino Unido; el Tribunal no siguió esa recomendación en lo relativo a España, un ejemplo poco frecuente pero real de divergencia entre la propuesta del Abogado General y el fallo final—. [Probable] Esta saga judicial tendría todavía una réplica posterior en Alemania, con el caso de la llamada "Ley Volkswagen" (Comisión contra Alemania, C-112/05, 2007), que protegía a ese fabricante de una toma de control hostil mediante un límite legal —no una acción de oro en sentido estricto, pero con un efecto funcional equivalente— a los derechos de voto de cualquier accionista individual por encima del 20%, también declarado contrario a la libre circulación de capitales.

[Suposición] El patrón que emerge de esta línea jurisprudencial, y que conviene que traslades como conclusión de cierre de todo el bloque 1, es que el Tribunal europeo nunca prohibió la acción de oro como instrumento en abstracto —el caso belga de 2002 lo demuestra—, sino que exigió que cualquier restricción a la libre circulación de capitales se justificara en un interés general o estratégico **preciso, conocido de antemano y proporcionado**, y no en una reserva de control genérica y discrecional del Estado sobre lo que fueron sus propias empresas. Es, en cierto modo, el primer indicio —todavía dentro del bloque que defiende la existencia histórica de la empresa pública— de la tensión que va a dominar por completo el bloque 2: la desconfianza creciente, tanto política como jurídica, hacia cualquier forma de control estatal sobre la empresa que no pueda justificarse con un argumento económico preciso y no meramente heredado de una época distinta.


\clearpage
\section{Empresa pública: ¿Por qué debería no existir?}
\puntosclave{
    
}

El bloque 1 ha mostrado un tejido público empresarial construido sobre argumentos sólidos y sostenido, durante tres décadas de posguerra, por un consenso doctrinal amplio que trascendía las fronteras ideológicas convencionales. Este segundo bloque explica cómo y por qué ese consenso se rompió. No fue una ruptura repentina ni exclusivamente política: fue la confluencia de una crisis macroeconómica que deslegitimó el paradigma keynesiano en su conjunto y de un aparato teórico —construido durante años, en buena medida al margen del debate público— que estaba ya preparado para ofrecer una explicación alternativa en el momento en que esa crisis estalló. Sigo, por tanto, el orden que hemos acordado: primero el contexto que explica por qué unos economistas concretos empezaron a investigar en esta dirección (2.1); después los argumentos teóricos que construyeron, con ese contexto como telón de fondo (2.2); después la ola histórica de privatización que esos argumentos legitimaron (2.3); y cierro con el propio proceso técnico de privatización, que solo tiene sentido una vez se ha decidido, por las razones de los subapartados anteriores, que la desinversión es preferible al mantenimiento de la titularidad pública (2.4).

Ya cerramos el bloque 1 señalando el punto de ruptura: la crisis del petróleo de 1973 y la estanflación que la siguió. [Seguro] Conviene ahora detenerse en por qué ese episodio fue, específicamente, una crisis **teórica** y no solo económica. La ortodoxia keynesiana de posguerra descansaba sobre la curva de Phillips —la relación empírica inversa entre inflación y desempleo que permitía a los gobiernos, en teoría, "comprar" menos paro a cambio de algo más de inflación mediante política de demanda—. La estanflación de los años setenta, con inflación y desempleo elevados **simultáneamente**, contradecía esa relación de forma tan visible que ningún ajuste menor del modelo bastaba para salvarlo: hacía falta, o bien reformular por completo la macroeconomía keynesiana —lo que efectivamente ocurrió, con la incorporación de expectativas racionales que asociamos a Robert Lucas—, o bien abandonar el paradigma completo en favor de una alternativa. [Probable] Fue precisamente en ese vacío de legitimidad teórica donde el monetarismo de Milton Friedman —que llevaba defendiendo desde los años cincuenta y sesenta, con escaso eco político, la primacía de la política monetaria sobre la fiscal y un papel mínimo del Estado en la economía (*Capitalism and Freedom*, 1962)— encontró, por fin, una audiencia política dispuesta a escucharlo. El propio O'Connor, que ya citamos al cerrar el bloque 1, diagnosticaba el problema desde el lado opuesto del espectro ideológico: el Estado de posguerra había asumido compromisos de gasto —incluida la financiación de un tejido público empresarial extenso— que su propia base fiscal, erosionada por la crisis, ya no podía sostener sin generar déficits crecientes. Es decir: **izquierda y derecha diagnosticaban, por razones distintas, el mismo síntoma de agotamiento del modelo**, lo que explica por qué la crítica que sigue no fue nunca un fenómeno exclusivamente de un color político, aunque su instrumentalización más visible —Thatcher, Reagan— sí lo fuera.

\begin{specialblock}{Public Choice School: Precursor intelectual directo}
    [Seguro] Pero el aparato teórico específico que se aplicaría a la empresa pública no nació de la crisis de 1973: llevaba gestándose desde más de una década antes, en un rincón entonces marginal de la profesión económica conocido como la **escuela de la elección pública** (*public choice*), fundada por James Buchanan y Gordon Tullock con su obra fundacional *The Calculus of Consent* (1962). [Probable] Su innovación metodológica, que puede parecer obvia hoy precisamente porque tuvo tanto éxito, era en su momento radical: aplicar el mismo supuesto de comportamiento racional y maximizador que la economía usa para modelar a consumidores y empresas a los **políticos y burócratas**, en lugar de asumir, como hacía implícitamente buena parte de la economía del bienestar de la época, que el Estado actúa como un planificador benevolente cuyo único objetivo es maximizar el bienestar social. [Seguro] Buchanan recibiría el Premio Nobel de Economía en 1986 "por su desarrollo de las bases contractuales y constitucionales de la teoría de la toma de decisiones económicas y políticas" —un reconocimiento tardío a una escuela que, en el momento de su fundación en 1962, apenas tenía predicamento fuera de un pequeño círculo académico en Virginia—.
    
    [Probable] Dentro de este mismo programa de investigación, la contribución específica de **Gordon Tullock** sobre la **búsqueda de rentas** (*rent-seeking*, formulada en su artículo de 1967, "The Welfare Costs of Tariffs, Monopolies, and Theft", *Western Economic Journal*) resulta especialmente relevante para lo que viene después en este bloque: Tullock mostró que, cuando el Estado tiene capacidad de conceder privilegios económicos —una licencia, un monopolio protegido, o, en lo que aquí nos interesa, la gestión de una empresa pública—, los agentes privados no se limitan a competir produciendo con mayor eficiencia, sino que invierten recursos reales en **influir políticamente** sobre quién obtiene ese privilegio, un gasto que no genera ningún valor social nuevo y que constituye, por tanto, una pérdida de eficiencia adicional a la del propio monopolio. Este concepto —que un privilegio estatal genera, a su alrededor, un ecosistema de actores dispuestos a gastar recursos reales para capturarlo o influir en su gestión— es el eslabón que conecta directamente la teoría de la elección pública con la crítica específica a la empresa pública que desarrollaremos en el subapartado 2.2: cuando allí veamos a **Vickers y Yarrow (1988)** describir a los políticos persiguiendo la maximización de votos mediante el mantenimiento de plantillas sobredimensionadas o salarios superiores a la productividad marginal, o a **Niskanen (1975)** describir al burócrata persiguiendo su propio prestigio antes que el bienestar social agregado, no estaremos ante ideas surgidas de la nada, sino ante la **aplicación directa, veinte años después, del mismo programa metodológico que Buchanan y Tullock habían inaugurado en 1962**: sustituir el supuesto del planificador benevolente por el del agente político racional que persigue su propio interés dentro de las restricciones institucionales que enfrenta.
    
    [Suposición, dato de interés expositivo] Un precursor adicional, menos citado en el debate específico de la empresa pública pero igualmente influyente en la genealogía de esta escuela, es Anthony Downs (*An Economic Theory of Democracy*, 1957), quien había aplicado ya, antes que Buchanan y Tullock, el supuesto de racionalidad económica al comportamiento electoral —el votante que se abstiene de informarse plenamente porque el coste de hacerlo supera el beneficio esperado de su voto individual (la "ignorancia racional")—, sentando así una de las piezas complementarias del rompecabezas: si ni el votante tiene incentivos para vigilar con atención lo que hace el político, ni el político tiene incentivos para actuar exclusivamente en interés general, la cadena de rendición de cuentas democrática sobre la gestión de una empresa pública —de la ciudadanía al político, del político al gestor— queda, en el mejor de los casos, debilitada en cada uno de sus eslabones.
    
    #### 2.1.3. Un matiz que conviene anticipar
    
    [Suposición] Antes de desarrollar con todo detalle, en el subapartado siguiente, el aparato crítico que esta escuela construyó específicamente contra la empresa pública, vale la pena dejar apuntada —sin desarrollarla todavía, porque pertenece por derecho propio al bloque 3— una objeción que la propia disciplina ha planteado contra el supuesto de partida de la escuela de la elección pública: que reducir sistemáticamente la motivación de políticos y funcionarios al interés propio es, en sí mismo, un supuesto empírico discutible y no una verdad axiomática. Julian Le Grand, en su distinción entre "caballeros y bribones" (*knights and knaves*, 1997), sostiene que buena parte del diseño de política pública de los ochenta y noventa asumió, sin suficiente evidencia, que los profesionales del sector público eran "bribones" egoístas cuando la evidencia sobre **motivación de servicio público** (*public service motivation*) sugiere que muchos actúan, al menos parcialmente, como "caballeros" genuinamente orientados al interés general. Dejo esta tensión apuntada porque es, precisamente, el tipo de matiz "granular" que hemos reservado para el bloque 3 de este tema, y que sería prematuro resolver aquí, en un bloque cuya función es presentar la crítica en su versión más fuerte y mejor construida, no en su versión ya atenuada por décadas de contraargumento posterior.
\end{specialblock}

\subsection{Argumentos teóricos en contra}

Con el contexto ya fijado, desarrollo ahora el aparato crítico en sí, recuperando y ampliando sustancialmente el contenido que ya tenía el documento original —Vickers-Yarrow, Niskanen, la doble relación de agencia— e incorporando el marco teórico que le faltaba para alcanzar el nivel de profundidad que corresponde a este apartado.

### 2.2.1. La elección pública aplicada: los incentivos de políticos y burócratas

[Seguro] La aplicación específica del programa de Buchanan y Tullock a la empresa pública tiene su formulación de referencia en **Vickers y Yarrow (1988, *Privatization: An Economic Analysis*, MIT Press)** —la obra que, más que ninguna otra, sistematizó el argumento económico de la privatización británica de los ochenta—. Su tesis central, ya presente en el documento original pero que conviene ahora desarrollar en sus mecanismos concretos, es que los políticos, si persiguen la maximización de votos, tienen incentivos sistemáticos para desviar la gestión de la empresa pública de la maximización de beneficios hacia comportamientos electoralmente rentables: mantener plantillas por encima del nivel que maximizaría beneficios (porque el despido de trabajadores es visible y electoralmente costoso, mientras que la ineficiencia derivada del sobredimensionamiento es difusa y se reparte entre todos los contribuyentes); pagar salarios superiores a la productividad marginal de los trabajadores (porque comprar la paz social con los sindicatos del sector, especialmente en vísperas electorales, tiene un rédito político inmediato); o favorecer con sus decisiones de localización e inversión a colectivos o territorios concretos de interés electoral estratégico, en lugar de a los criterios de eficiencia asignativa que seguiría un gestor privado.

[Probable] La pieza complementaria es la de **William Niskanen** (*Bureaucracy and Representative Government*, 1971, y su desarrollo posterior "Bureaucrats and Politicians", *Journal of Law and Economics*, 1975), quien modela al burócrata —el gestor de la empresa pública, en última instancia— no como un maximizador puro de bienestar social ni tampoco como un maximizador puro de interés personal, sino como alguien que persigue maximizar el **presupuesto** bajo su control, en la medida en que un presupuesto mayor suele traer aparejados mayor prestigio, mayor poder discrecional y mejores perspectivas de carrera. [Probable] Conviene señalar, porque añade matiz y rigor al argumento, que el propio modelo de Niskanen fue objeto de una crítica influyente por parte de **Migué y Bélanger** (1974, "Toward a General Theory of Managerial Discretion", *Public Choice*), quienes argumentaron que el burócrata no maximiza el presupuesto total —parte del cual se destina, inevitablemente, a cubrir costes de producción que no le reportan ningún beneficio personal—, sino específicamente el **presupuesto discrecional**, el margen de holgura que queda una vez cubiertos los costes estrictamente necesarios: una versión más afinada del mismo argumento, que explica mejor por qué la ineficiencia burocrática tiende a manifestarse en gastos suntuarios o en personal poco productivo antes que en un simple sobredimensionamiento uniforme de toda la plantilla.

### 2.2.2. La teoría de la agencia: la doble relación y la ausencia de mecanismos de mercado

[Seguro] La formalización moderna del problema de agencia se debe a **Jensen y Meckling** (1976, "Theory of the Firm: Managerial Behavior, Agency Costs, and Ownership Structure", *Journal of Financial Economics*), que define la relación de agencia como aquella en la que un principal encarga a un agente la toma de decisiones en su nombre, y que surge un coste —el **coste de agencia**— cuando los incentivos de ambos no están perfectamente alineados y existe, además, asimetría informativa sobre las actuaciones del agente que el principal no puede observar directamente. En la empresa privada, este problema ya existe entre accionistas y directivos; en la empresa pública, como recogía correctamente el documento original, se **duplica**: los ciudadanos, titulares últimos de la empresa, delegan en los políticos, que a su vez delegan en los directivos, generando una **cadena de agencia con dos eslabones** en lugar de uno, cada uno con su propio problema de alineación de incentivos e información asimétrica.

[Probable] Lo que el documento original no explicaba con suficiente profundidad es **por qué**, más allá de esta duplicación estructural, la empresa pública carece de los mecanismos correctores que sí operan sobre la empresa privada cotizada. El primero es el **precio de las acciones** como señal continua de desempeño: cualquier indicio de mala gestión percibido por el mercado se traduce de inmediato en una caída del precio, una señal pública, cuantificada y difícil de ocultar que no tiene equivalente en una empresa cuyo "accionista" es el conjunto de la ciudadanía y no puede, individualmente, vender su participación. Esta imposibilidad de "salida" —**exit**, en la terminología clásica de **Albert Hirschman** (*Exit, Voice, and Loyalty*, 1970)— deja al ciudadano-propietario únicamente con el mecanismo de la **voz** —el voto, la protesta, la presión política— como instrumento de disciplina sobre la gestión de "su" empresa, un mecanismo mucho más lento, más costoso de ejercer individualmente, y sujeto a los mismos problemas de acción colectiva e ignorancia racional que ya identificamos con Downs en el subapartado anterior. El segundo mecanismo ausente es el de la **oferta pública de adquisición hostil**, cuya función disciplinadora sobre los gestores de empresas cotizadas formalizó Henry Manne (1965, "Mergers and the Market for Corporate Control", *Journal of Political Economy*): la amenaza creíble de que un equipo gestor ineficiente sea destituido mediante la compra de la empresa por un tercero que aprecia el potencial de mejora de gestión no operativa en la empresa pública, precisamente porque, como ya vimos al hablar de la acción de oro en el subapartado 1.3, el control estatal —accionarial o mediante privilegios especiales— hace inviable ese tipo de operación incluso cuando la empresa cotiza formalmente en bolsa.

### 2.2.3. La restricción presupuestaria blanda

[Seguro] El tercer mecanismo, ya apuntado en el documento original sin el respaldo bibliográfico que merece, es la **restricción presupuestaria blanda** (*soft budget constraint*), concepto acuñado por el economista húngaro **János Kornai** (*Economics of Shortage*, 1980) para explicar las ineficiencias sistemáticas de las economías centralmente planificadas: una empresa opera bajo restricción presupuestaria blanda cuando sabe, de antemano, que en caso de pérdidas o incluso de quiebra técnica será rescatada por una autoridad superior —el Estado—, lo que elimina el incentivo a operar con la misma disciplina de costes que una empresa sujeta a la amenaza real de bancarrota. [Probable] El concepto, nacido para explicar el socialismo de Estado, resultó tener un poder explicativo mucho más amplio del que Kornai anticipó originalmente: **Kornai, Maskin y Roland**, en su artículo de síntesis (2003, "Understanding the Soft Budget Constraint", *Journal of Economic Literature*), extienden el marco a cualquier organización —bancos "demasiado grandes para caer", entidades de crédito al desarrollo, y, por supuesto, empresas públicas en economías de mercado— cuyo comportamiento venga determinado, no ya por su rentabilidad esperada, sino por la expectativa racional de rescate. Este problema genera un **riesgo moral** doble, que afecta tanto al gestor —que no internaliza las consecuencias de una gestión temeraria del mismo modo que lo haría un directivo cuyo empleo depende de la supervivencia financiera de la empresa— como al propio político que la controla, que puede preferir posponer decisiones dolorosas de reestructuración hasta después del siguiente ciclo electoral, sabiendo que el coste de esa demora recaerá, en última instancia, sobre el conjunto de los contribuyentes.

### 2.2.4. La frontera moderna: contratos incompletos y la propiedad como derecho de control residual

[Seguro] El desarrollo teórico más reciente y más riguroso de todo este bloque —ausente por completo del documento original, y que conviene incorporar porque es hoy el punto de partida obligado de cualquier análisis académico serio sobre titularidad pública frente a privada— procede de la teoría de los **contratos incompletos**, cuyo fundamento se remonta a Grossman y Hart (1986, "The Costs and Benefits of Ownership: A Theory of Vertical and Lateral Integration", *Journal of Political Economy*): dado que ningún contrato puede especificar, de antemano, cómo debe actuar cada parte ante cualquier contingencia futura imaginable, la propiedad de un activo importa porque otorga a su titular los **derechos de control residual** —el derecho a decidir sobre todo aquello que el contrato no ha previsto explícitamente—.

[Seguro] **Hart, Shleifer y Vishny** (1997, "The Proper Scope of Government: Theory and an Application to Prisons", *Quarterly Journal of Economics*) aplican directamente este marco a la disyuntiva entre provisión pública y privada. Su modelo distingue dos dimensiones del desempeño de cualquier proveedor de un servicio: la **reducción de costes**, que es relativamente fácil de especificar y verificar en un contrato (y por tanto "contratable"), y la **mejora de calidad**, que a menudo no lo es (resulta difícil, por ejemplo, especificar contractualmente con precisión el nivel exacto de trato humano que debe recibir un preso en una prisión gestionada por un tercero, el ejemplo concreto que los autores desarrollan). Su resultado central es que, bajo contratos incompletos, un proveedor privado tiene **siempre más incentivos** que un empleado público tanto para reducir costes como para mejorar la calidad —porque, a diferencia del empleado público, se queda con parte del excedente que genera cualquier mejora—, pero ese incentivo a reducir costes es, con frecuencia, **excesivo**: el proveedor privado tiende a recortar exactamente en aquello que no está contractualmente especificado —la calidad no verificable—, porque hacerlo reduce sus costes sin que el contrato lo penalice. [Probable] La conclusión normativa del modelo es, por tanto, mucho más matizada que un simple "lo privado es más eficiente": la privatización resulta preferible cuando la calidad es relativamente fácil de contratar y verificar, y la provisión pública resulta preferible cuando el margen para un recorte de calidad no contratable —y potencialmente grave— es amplio, precisamente el tipo de análisis "granular" que anticipábamos como propio del bloque 3, pero cuyo fundamento teórico pertenece, con toda propiedad, a este bloque crítico. [Seguro] **Andrei Shleifer**, en su artículo de síntesis posterior (1998, "State versus Private Ownership", *Journal of Economic Perspectives*), generaliza esta misma lógica de derechos de control residual a la elección entre titularidad pública y privada en general, convirtiéndose en la referencia de cabecera de toda la literatura moderna sobre esta materia y en el puente natural hacia el debate empírico que desarrollaremos en el bloque 3.

### 2.2.5. Rigideces prácticas: contratación, procedimientos y financiación

[Seguro] Cierro este subapartado recuperando, con la extensión que merece, el contenido más aplicado que ya recogía correctamente el documento original. Más allá de los problemas de incentivos que acabamos de desarrollar en términos puramente teóricos, la empresa pública se enfrenta a restricciones **operativas** de las que la empresa privada está exenta, y que tienen su propia justificación —evitar el abuso en la gestión de fondos públicos— pero también su propio coste de eficiencia. La obligación de someter la contratación de bienes y servicios a **concurso público** introduce, junto a la garantía de transparencia que persigue, tres costes añadidos: el coste burocrático directo de un procedimiento más complejo, el coste de oportunidad derivado de su mayor duración, y el propio coste burocrático que deben asumir las empresas licitadoras para poder concurrir. En el terreno presupuestario, mientras la empresa privada que detecta una oportunidad de inversión puede movilizar financiación con relativa rapidez, la empresa pública puede necesitar una modificación presupuestaria sujeta a plazos administrativos y, en el peor de los casos, esperar a la aprobación de los presupuestos del ejercicio siguiente —una asimetría temporal que puede resultar decisiva en sectores donde la ventana de oportunidad de una inversión es estrecha—. Y en el terreno laboral, el principio de mérito y capacidad que rige el acceso a la función pública —diseñado, precisamente, para impedir que el interés particular sustituya al interés general en la selección de personal— introduce rigideces tanto en la contratación como en el despido, y limita la capacidad de la empresa pública para competir por talento mediante una política salarial flexible, en un contraste directo con la libertad casi total de la empresa privada para remunerar, contratar y prescindir de sus gestores según su propia valoración de productividad.

## 2.3. La ola histórica de privatización

### 2.3.1. El Reino Unido de Thatcher: el caso paradigmático

[Seguro] Si el subapartado anterior explicaba de dónde procedía el aparato teórico de la crítica, este explica dónde y cómo se llevó, por primera vez y a mayor escala, a la práctica de gobierno. El **Reino Unido bajo Margaret Thatcher** (1979-1990) es el caso paradigmático no porque fuera el primer país en privatizar —Alemania Occidental había vendido ya, en los años sesenta, su participación mayoritaria en Volkswagen—, sino por la **escala, la sistematicidad y la ambición ideológica explícita** del programa, que llegó a convertirse en el modelo de referencia exportado, con variantes, a prácticamente todo el mundo desarrollado y en desarrollo durante las dos décadas siguientes.

[Probable] Conviene narrar el proceso en sus dos fases, porque revela algo interesante sobre la propia lógica de la reforma: en su primer mandato (1979-1983), el programa fue, de hecho, **modesto** y se limitó a empresas que ya operaban en mercados razonablemente competitivos —British Aerospace, Cable & Wireless, Amersham International—, sin tocar todavía los grandes monopolios de red. Fue en el **segundo mandato** cuando el programa se radicalizó y alcanzó a los sectores regulados propiamente dichos: **British Telecom en 1984**, la primera gran privatización de una utility, seguida de **British Gas en 1986** —cuya campaña publicitaria "Tell Sid", que invitaba a cualquier ciudadano corriente ("Sid") a comprar acciones asequibles de la empresa, se convirtió en el símbolo cultural de toda la era—, la British Airports Authority en 1987, las compañías de agua en 1989, y la electricidad, ya bajo el sucesor de Thatcher, John Major, en 1990-91. [Seguro] El resultado agregado, medido en tres décadas, fue una transformación estructural de primer orden: más de 60.000 millones de libras en activos estatales vendidos, y la proporción del empleo total correspondiente a las industrias nacionalizadas cayendo del **9% al menos del 2%**.

[Seguro] Pero el aspecto más interesante de la experiencia británica para este tema —el que conecta directamente con la nota sobre "capitalismo popular" que ya avanzamos al hablar de las formas de venta en el bloque 3— no es solo el volumen de la desinversión, sino su **objetivo político explícito**: Thatcher no perseguía únicamente eficiencia económica, sino construir deliberadamente una "democracia de accionistas" (*shareholding democracy*), en palabras del propio Gobierno de la época, que redujera el poder sindical y ampliara la base social del conservadurismo mediante la difusión masiva de la propiedad accionarial entre trabajadores que tradicionalmente habían votado laborista. [Seguro] El indicador más citado de este objetivo es que la proporción de ciudadanos británicos que poseían acciones pasó de apenas el **7% a más de una cuarta parte de la población adulta** en poco más de una década, con algunos recuentos situando en su momento álgido hasta 32,6 millones de personas —tres quintas partes de la población británica— como titulares de acciones de alguna empresa recién privatizada. [Probable] Este es, en sí mismo, un argumento **de economía política y no de eficiencia**, que conviene distinguir con precisión de los argumentos de agencia y contratos incompletos del subapartado 2.2: Thatcher no estaba corrigiendo un problema de incentivos gerenciales, sino usando la privatización como instrumento para **reconfigurar la coalición electoral** de la que dependía la propia continuidad del programa reformista —un uso instrumental de la política económica que Boycko, Shleifer y Vishny (1996) formalizarían después, en clave más general, como el argumento de que la privatización sirve para comprometerse creíblemente frente a la interferencia política futura, precisamente porque una base amplia de pequeños accionistas-votantes hace políticamente muy costoso revertir la venta—.

[Suposición, dato de actualidad relevante] Conviene cerrar este caso con una nota que desmiente cualquier lectura de éxito indiscutido: buena parte de la literatura crítica reciente señala que numerosas privatizaciones británicas, especialmente las de agua y electricidad, se vendieron a precios sustancialmente **inferiores** a su valor real —en parte por el interés del Gobierno en garantizar el éxito de la colocación—, lo que permitió a los primeros inversores obtener ganancias extraordinarias a costa del erario público; y las encuestas de opinión británicas más recientes muestran mayorías muy amplias —en torno al 75-83% según el sector— a favor de **renacionalizar** el agua, la electricidad, el gas y los ferrocarriles, en un giro de opinión pública que sugiere que el legado de la "democracia de accionistas" thatcherista está hoy lejos de ser un consenso estable, cuarenta años después de su puesta en marcha.

### 2.3.2. Otras trayectorias: Francia, la transición postcomunista, el Consenso de Washington y España como una ilustración más

[Probable] El modelo británico no se replicó de forma idéntica en ningún otro país, y las variantes son, en sí mismas, instructivas. En **Francia**, el proceso siguió un patrón de vaivén político que no tiene equivalente en el caso británico: tras la segunda ola de nacionalizaciones de Mitterrand de 1981-1982 que ya vimos en el bloque 1, el Gobierno de cohabitación de Jacques Chirac (1986-1988) revirtió buena parte de esas nacionalizaciones mediante un programa de privatizaciones, que a su vez sería parcialmente frenado tras el retorno de la izquierda al poder en 1988 —una alternancia que muestra hasta qué punto la titularidad empresarial siguió siendo, en Francia, un terreno de disputa política explícita mucho más tiempo que en el Reino Unido—.

[Seguro] En **Italia**, el desmantelamiento del IRI —el conglomerado que, como vimos en el bloque 1, había nacido en 1933 como rescate financiero y se había convertido en el mayor holding público de Europa Occidental— se aceleró en los años noventa por una razón que conecta directamente con el criterio de clasificación estadística que desarrollamos en el subapartado 1.3.2: la necesidad de **reducir la deuda pública italiana** para cumplir los criterios de convergencia del Tratado de Maastricht y poder acceder a la moneda única europea convirtió la venta de activos del IRI en un instrumento de consolidación fiscal tan importante como cualquier argumento de eficiencia productiva —el mismo mecanismo, "razones presupuestarias" antes que "razones de eficiencia", que ya identificamos como una de las dos grandes vías de justificación de cualquier privatización—.

[Seguro] La **transición postcomunista** de Europa del Este tras 1989 constituye un tercer patrón, radicalmente distinto de los dos anteriores por su velocidad y su escala: no se trataba de privatizar un número limitado de grandes empresas estratégicas en el contexto de una economía de mercado ya funcionante, sino de transformar de la noche a la mañana la totalidad del tejido productivo de economías enteramente planificadas. Los métodos variaron drásticamente de un país a otro —desde la **privatización mediante cupones o vales** (*voucher privatization*) checoslovaca, que distribuyó certificados de propiedad gratuitos entre la población para canjear por acciones de las antiguas empresas estatales, hasta el infame programa ruso de **"préstamos por acciones"** (*loans-for-shares*) de 1995-1996, mediante el cual un pequeño grupo de bancos privados concedió préstamos al Estado ruso a cambio de participaciones en las mayores empresas energéticas y de recursos naturales del país como garantía, préstamos que el Estado —previsiblemente, y según numerosos análisis posteriores, de forma deliberadamente pactada— no pudo devolver, lo que permitió a ese reducido grupo de banqueros hacerse con activos de un valor muy superior al de los préstamos concedidos, dando origen a la clase de los llamados **oligarcas** rusos—. Ya citamos, en la auditoría inicial de este tema, el meta-análisis de **Djankov y Murrell (2002)** sobre los efectos de la privatización en estas economías de transición, que documenta mejoras significativas de eficiencia técnica y asignativa derivadas de la reestructuración empresarial asociada al cambio de titularidad, pero sin que ese resultado agregado oculte la enorme disparidad de legitimidad y de resultados distributivos entre el modelo checo y el modelo ruso.

[Seguro] En **América Latina**, el proceso de privatización de los años ochenta y noventa se integró dentro de un paquete de política económica más amplio, sistematizado por el economista John Williamson bajo la etiqueta que se haría célebre —el **Consenso de Washington** ("What Washington Means by Policy Reform", 1989)—: disciplina fiscal, liberalización comercial y financiera, y **privatización de empresas estatales**, formaban parte de un mismo recetario de reforma estructural promovido por las instituciones financieras internacionales con sede en Washington (el Fondo Monetario Internacional y el Banco Mundial) como condición, con frecuencia, para el acceso a financiación externa —lo que explica por qué, a diferencia del caso británico, la privatización latinoamericana estuvo, en muchos países, tan condicionada por la presión de acreedores externos como por una convicción ideológica doméstica plenamente asumida—.

[Probable] **España**, en este mapa comparado, es una ilustración más entre varias, no un caso excepcional: el INI que conocimos en el bloque 1 se transformó, a través de sucesivas reformas, en la actual SEPI (Sociedad Estatal de Participaciones Industriales), y el país vivió su propia ola de privatizaciones en dos tiempos —una primera fase más gradual bajo los gobiernos de Felipe González a finales de los ochenta y principios de los noventa, y una segunda fase de privatización plena y acelerada bajo los gobiernos de José María Aznar entre 1996 y 2004, que completó la venta de Telefónica, Repsol, Endesa, Argentaria (posteriormente fusionada con BBV para formar BBVA), Iberia y Tabacalera— un proceso cronológicamente posterior al británico y, en su diseño técnico, deudor de la experiencia acumulada por el Reino Unido y por los organismos internacionales que estudiaremos en el subapartado siguiente.


\subsection{Proceso de privatización}
Una vez que el Estado ha decidido, por las razones teóricas e históricas de los subapartados anteriores, que la desinversión es preferible al mantenimiento de la titularidad pública, se abre una pregunta distinta y de naturaleza mucho más práctica, que es la que da contenido a este subapartado: **¿cómo debe ejecutarse esa venta para que cumpla su propósito?** [Seguro] La literatura identifica dos objetivos que el sector público persigue simultáneamente a lo largo de todo el proceso, y que no siempre son compatibles entre sí. El primero es **maximizar los ingresos** obtenidos por la venta, idealmente vendiendo la empresa a un precio que refleje su valor fundamental; el segundo es **asegurar el buen funcionamiento posterior del mercado** en el que la empresa privatizada va a operar, lo que exige atención a la regulación aplicable, especialmente cuando el sector no es plenamente competitivo —recordando aquí toda la discusión sobre monopolio natural del subapartado 1.2—. La tensión entre ambos objetivos no es anecdótica: un gobierno que priorice exclusivamente maximizar el precio de venta puede sentirse tentado a vender la empresa **sin** desmantelar previamente su posición de monopolio —un monopolio privado se vende más caro que una empresa que va a enfrentarse a competencia inmediata—, sacrificando así el segundo objetivo en beneficio del primero, un dilema al que volveremos en el subapartado 3.1 cuando revisemos la evidencia empírica sobre si la privatización, por sí sola, genera ganancias de eficiencia en sectores poco competitivos.

\subsubsection{Preparación del proceso: el problema de la información asimétrica}

[Probable] Antes de exponer, como hacía el documento original, los cinco pasos previos de cualquier proceso de privatización bien diseñado, conviene dotarlos de un marco teórico común que el documento no tenía: todos ellos son, en el fondo, **respuestas distintas al mismo problema de información asimétrica** que George Akerlof formalizó en su artículo fundacional sobre el mercado de "limones" (1970, "The Market for Mechanisms", *Quarterly Journal of Economics* — título real: "The Market for 'Lemons': Quality Uncertainty and the Market Mechanism"). Akerlof mostró que, cuando el vendedor de un bien conoce su calidad real mejor que el comprador potencial —el caso típico del mercado de coches usados, donde el propietario sabe si su vehículo es fiable o un "limón" y el comprador no—, los compradores, racionalmente desconfiados, solo están dispuestos a pagar un precio que refleje la **calidad media esperada** del conjunto de bienes en venta, no la calidad real del bien concreto que se les ofrece. Esto empuja a los vendedores de bienes de buena calidad a abandonar el mercado —porque no pueden obtener un precio que refleje su verdadera calidad—, dejando progresivamente solo "limones" en circulación, en una espiral de selección adversa que puede llegar a colapsar el mercado por completo. [Probable] Trasladado a la privatización, el problema es formalmente idéntico: el Estado conoce la verdadera rentabilidad potencial de la empresa que va a vender mejor que cualquier comprador externo, y los compradores, conscientes de esa asimetría, descontarán sistemáticamente el precio que están dispuestos a pagar para protegerse frente al riesgo de estar comprando un "limón" —una empresa pública con pasivos ocultos, activos sobrevalorados o problemas operativos que el Estado prefiere no revelar—. Cada uno de los cinco pasos que desarrollo a continuación puede leerse, bajo esta luz, como un mecanismo para **reducir esa asimetría informativa** y evitar así el descuento que los compradores aplicarían por pura precaución.

[Seguro] El primero es la **responsabilidad administrativa**: cuando el volumen de activos a privatizar es elevado, resulta recomendable crear una **agencia de privatización especializada** que concentre la responsabilidad del proceso y se someta a estándares de transparencia máximos; si el volumen no justifica ese coste fijo, debe al menos identificarse con claridad la unidad administrativa responsable. El segundo son los **asesores externos**: los gobiernos recurren con frecuencia a asesores de mayor experiencia técnica para diseñar el proceso, con el riesgo principal de que incurran en conflictos de interés —entre los países de la OCDE, Australia, Polonia y España destacan por encargar sistemáticamente valoraciones externas de las empresas a privatizar, siendo obligatorio por ley en los dos últimos—. El tercero es la **regulación y competencia previas**: no se recomienda acometer la privatización sin haber establecido de antemano un marco regulatorio idóneo cuando la introducción de competencia efectiva no sea posible, encomendando idealmente esa regulación a un órgano dotado de autonomía respecto del propio gobierno vendedor —un punto que conecta directamente con el objetivo de "buen funcionamiento posterior del mercado" del subapartado anterior—.

[Seguro] El cuarto —las **condiciones laborales**— es, con frecuencia, el obstáculo político más determinante de cualquier proceso: cuando la empresa pública no tiene forma societaria y sus trabajadores son funcionarios, su traslado a manos privadas genera un problema jurídico propio que exige recolocación o disposiciones legales específicas; cuando sí tiene forma societaria, el problema es más bien de protesta sindical y conflicto social. El *policymaker* se enfrenta aquí a un trade-off que el propio documento original ya identificaba con precisión: garantizar el mantenimiento de las condiciones laborales facilita la aceptación social del proceso, pero limita las ganancias de eficiencia que la reestructuración laboral podría generar. El caso paradigmático de gestión de este trade-off sigue siendo la privatización de **Brazil Railways en el año 2000**, donde el gobierno inició la reestructuración de plantilla —de 42.000 trabajadores, 12.000 optaron por jubilación anticipada y otros 6.000 por abandono voluntario, reduciendo la plantilla en un 80% entre 1998 y mayo de 2000— **antes** de la propia privatización, financiando el coste social del ajuste con fondos gubernamentales e internacionales, precisamente para no trasladar ese coste de golpe al futuro comprador y facilitar así una valoración más limpia de la empresa.

[Probable] El quinto —la **reestructuración previa de la empresa**— es el que mejor ilustra la lógica de Akerlof que hemos usado para unificar todo este subapartado. La postura tradicional sostenía que estas reformas no debían acometerse antes de la privatización, porque el futuro comprador, conocedor del sector, estaría mejor capacitado para llevarlas a cabo, y el coste de esas reformas —descontado del precio de venta— sería inferior al de intentarlas dentro de la propia Administración (Chong y López-de-Silanes, 2004). Pero esa lógica **presupone que existe ya un comprador identificado**, capaz de valorar con precisión el potencial de la empresa pese a la asimetría informativa. Cuando no hay comprador identificado —singularmente en una Oferta Pública de Venta dirigida a miles de inversores anónimos, que carecen individualmente de la capacidad de investigación que sí tendría un único comprador estratégico—, realizar ajustes previos cumple exactamente la función que la teoría de la señalización, desde Akerlof, atribuye a cualquier mecanismo costoso de revelación de información: reduce la incertidumbre sobre los flujos de caja futuros y mitiga la asimetría informativa, lo que en conjunto **aumenta la disposición a pagar** de los compradores potenciales, precisamente porque ya no necesitan aplicar el descuento defensivo que aplicarían frente a un "limón" no verificado. Es coherente, por tanto, que en la práctica la reestructuración previa resulte más necesaria en las OPV que en las ventas directas a un único inversor estratégico —donde ese inversor puede investigar y valorar la empresa por sí mismo con mucho mayor detalle—, y que la valoración externa previa a la privatización sea, como ya señalamos, obligatoria por ley en varios países de la OCDE precisamente por esta misma lógica de reducción de asimetría informativa.

\subsubsection{Formas de venta}

Una vez completada la preparación, el proceso se articula en torno a dos grandes formas de venta —oferta pública y venta directa—, cada una con ventajas e inconvenientes que conviene contrastar antes de entrar en su mecánica concreta.

#### 2.4.3.1. La Oferta Pública de Venta

[Seguro] La OPV incorpora la empresa al mercado de valores mediante la emisión de acciones accesibles a cualquier inversor, hablándose de Oferta Pública Inicial (OPV/IPO) cuando la empresa no cotizaba previamente, y de privatización completa cuando los nuevos entrantes superan el 50% del capital. Requiere un tamaño suficiente de la empresa y mercados de capitales domésticos amplios y profundos —o, en su defecto, recurrir también a mercados extranjeros—, y su principal ventaja es la **transparencia** que imponen los estrictos requisitos regulatorios de las empresas cotizadas, junto con la disciplina que el propio mercado de capitales ejerce después sobre la gestión de la empresa ya privatizada, en línea con lo que ya vimos, en el subapartado 2.2.2, sobre el papel disciplinador del precio de las acciones y de la amenaza de OPA hostil.

[Seguro] El procedimiento sigue seis fases: (i) preparación de un folleto informativo veraz sobre el emisor y las acciones, con asesores financieros; (ii) formación del sindicato de colocación; (iii) aprobación del folleto; (iv) campaña de venta para maximizar la demanda; (v) fijación del rango de precios mínimo aceptable, rechazándose ofertas inferiores; (vi) apertura a la venta, con dos desenlaces posibles según se alcance o no el importe mínimo de suscripción. [Probable] Cuando el capital social a colocar es especialmente elevado, resulta habitual dividir el proceso en una **OPV inicial y una OPV secundaria** posterior —el ejemplo histórico citado en el documento original es la privatización de ENEL (Ente Nazionale per l'Energia Elettrica) en la Italia de los años noventa, debido a su enorme tamaño—. [Seguro] Esta práctica de privatización secuencial tiene, más allá de la mera conveniencia práctica de colocar un volumen manejable de acciones cada vez, un fundamento teórico riguroso que el documento original no explicaba: **Enrico Perotti** (1995, "Credible Privatization", *American Economic Review*) muestra formalmente que, cuando los inversores temen que el gobierno pueda interferir políticamente en el futuro de la empresa —revirtiendo parcialmente la privatización, imponiendo obligaciones de servicio público no compensadas, o simplemente renacionalizándola, como ya vimos que reclama hoy una mayoría de la opinión pública británica—, un gobierno que **retiene deliberadamente una participación pasiva** en la empresa parcialmente privatizada está, precisamente por asumir ese mismo riesgo residual que asumen los inversores privados, **señalizando de forma creíble su compromiso** de no interferir en el futuro: si el gobierno fuera a expropiar o interferir, también perdería valor en su propia participación retenida, lo que hace la amenaza de interferencia menos creíble a ojos de los inversores. [Probable] El propio Perotti extiende el argumento a la cuestión de la **infravaloración sistemática** que documentamos empíricamente al hablar del caso británico: cuando la participación estatal retenida entra en conflicto con la necesidad de transferir el control efectivo a los nuevos accionistas, puede hacer falta ofrecer las acciones con un descuento —infravalorarlas deliberadamente— para lograr esa separación de señales, lo que explicaría, con un argumento teórico y no solo como anomalía atribuible a un mal cálculo del gobierno vendedor, por qué tantas privatizaciones europeas de los ochenta y noventa se colocaron sistemáticamente por debajo de lo que resultó ser su valor de mercado posterior.

#### 2.4.3.2. La venta directa

[Seguro] La venta directa transfiere las acciones a un tercero, habitualmente un **inversor estratégico** de tamaño considerable que ya opera en el sector, lo que permite al vendedor público valorar criterios distintos del precio —experiencia previa en la gestión del bien o servicio, incorporación de mejoras tecnológicas y *know-how*, compromiso de inversiones futuras para garantizar la continuidad o la mejora del servicio, condiciones que no pueden imponerse en una OPV anónima—. Dentro de la venta directa se distinguen la **subasta**, más transparente y competitiva pero solo viable cuando existe más de un inversor idóneo, y con frecuencia complicada por la existencia de objetivos adicionales a la maximización del precio —mantenimiento de plantilla, no deslocalización, compromisos de inversión en I+D—, cuyo diseño simultáneo resulta técnicamente complejo; y la **venta negociada**, poco transparente pero inevitable cuando solo existe un comprador idóneo, y que exige de la entidad privatizadora especial cautela para no vender a un precio excesivamente conservador —una cautela con fundamento jurídico y no solo prudencial, porque un precio de venta anormalmente bajo puede llegar a calificarse, bajo el Derecho de la competencia europeo, como una ayuda de Estado indirecta al comprador (OCDE, 2009)—.

#### 2.4.3.3. La venta mixta

[Seguro] En muchas ocasiones ambas formas se combinan: se vende primero una parte a un inversor estratégico que acomete los ajustes necesarios para hacer la empresa más rentable y atractiva, y se coloca después el resto del capital mediante OPV pública. Esta combinación reúne parte de las ventajas de ambos métodos —la elección deliberada del primer inversor sin renunciar a la transparencia y accesibilidad de la colocación pública posterior—, y puede leerse, a la luz del modelo de Perotti que acabamos de exponer, como una forma particularmente sofisticada de secuenciación: el propio proceso de mejora de gestión bajo el inversor estratégico inicial funciona como una segunda forma de señalización de calidad frente a los inversores públicos que participarán en la fase posterior.

\subsubsection{Auditoría y control}

[Seguro] Cierro este subapartado con el control ex post del proceso, que en la mayoría de los países se articula mediante **auditoría y presentación ante el parlamento**. Esta rendición de cuentas pública se considera buena práctica precisamente porque aumenta la responsabilidad de los agentes implicados en un proceso que, como hemos visto a lo largo de todo este subapartado, está plagado de asimetrías de información entre el Estado vendedor, los compradores, y —no debe olvidarse— la propia ciudadanía, que es la titular última del activo que se vende. La institución encargada de esta auditoría debería, idealmente, estar presente a lo largo de todo el proceso —reestructuración, valoración, elección del método de venta, elección del comprador— y no limitarse a una revisión final desconectada de las decisiones que la precedieron (OCDE, 2009).



\clearpage
\section{Empresa Pública: ¿Cuándo debería existir?}
\puntosclave{
    
}   
    
Los dos bloques anteriores han presentado, cada uno con su propio rigor, dos respuestas completas y opuestas a la misma pregunta: primero, un cuerpo de argumentos que justificaba la empresa pública allí donde el mercado fallaba; después, un cuerpo de argumentos, igualmente sólido, que explicaba por qué esa misma solución generaba sus propios fallos, a menudo peores que el problema que pretendía resolver. Si el tema terminara aquí, el opositor se vería obligado a elegir un bando, y la propia historia reciente demuestra que ese no es el estado actual del debate. [Probable] El énfasis se ha desplazado, en las últimas dos décadas, desde la pregunta binaria "¿pública o privada?" hacia la pregunta condicional "¿bajo qué condiciones?", por razones que conviene enumerar porque no son solo académicas. La crisis financiera de 2008 obligó a los gobiernos de medio mundo a nacionalizar o recapitalizar entidades financieras que el propio consenso neoliberal había dado por definitivamente privatizadas, sin que nadie —ni siquiera sus críticos más acérrimos— cuestionara la necesidad de esa intervención en aquel momento concreto. La pandemia de COVID-19 reactivó, en 2020, un debate sobre la resiliencia de las cadenas de suministro estratégicas que recuerda, con sorprendente fidelidad, al argumento de los "sectores estratégicos" que ya desarrollamos en el subapartado 1.2.1.2. El ascenso de China como potencia económica de primer orden, construida en buena medida sobre un modelo de capitalismo de Estado que Occidente daba por superado, ha obligado a revisar la premisa implícita de que la propiedad privada es una condición necesaria del éxito económico a gran escala. Y, sobre todo, cuatro décadas de privatizaciones acumuladas han generado un volumen de evidencia empírica que permite, por fin, sustituir la predicción ideológica por la respuesta condicionada a los datos. Dividimos, en consecuencia, este último bloque en dos aproximaciones metodológicamente distintas: primero la cuantitativa, que contrasta hipótesis mediante análisis estadístico de muestras amplias de empresas; después la cualitativa —en el sentido, ya precisado, de argumentación teórica sustentada en estudios de caso extensos, no de metodología de entrevistas o encuesta—, que representa el estado más reciente del debate y que cierra el tema con las "nuevas razones" para la existencia de la empresa pública.
    
\subsection{Aproximación cuantitativa}

3.1.1. El problema metodológico de fondo

[Seguro] Antes de repasar los estudios concretos, hace falta entender por qué esta literatura, pese a llevar décadas acumulándose, no ha producido un consenso cerrado. La dificultad no es de escasez de datos, sino de diseño de la comparación. La literatura se articula en dos grandes bloques metodológicos, cada uno con sus propios sesgos potenciales. El primero compara la eficiencia de empresas públicas y privadas que coexisten simultáneamente en un mismo sector; el segundo compara la eficiencia de la misma empresa antes y después de ser privatizada. [Probable] El primer bloque se enfrenta a un problema de variables omitidas: si las empresas que permanecen en manos públicas no se reparten aleatoriamente entre los distintos segmentos de un sector, sino que tienden a concentrarse, precisamente, en aquellos considerados menos atractivos para el capital privado —por su menor rentabilidad esperada, por su función de servicio universal, por su relevancia estratégica—, cualquier diferencia de rentabilidad observada podría reflejar esa selección previa y no un efecto causal de la titularidad en sí misma. El segundo bloque se enfrenta a un problema simétrico de selección: los gobiernos no privatizan sus empresas al azar, sino que tienden a elegir, en primer lugar, aquellas que ya han sido objeto de reformas de gestión previas, o cuyas perspectivas de mejora son más visibles —precisamente el tipo de "reestructuración previa" que analizamos en el subapartado 2.4.2—, lo que puede inflar artificialmente la mejora de eficiencia atribuida a la privatización cuando en realidad una parte de esa mejora ya se había producido, o era ya previsible, antes del cambio de titularidad. A esto se añade una tercera fuente de confusión, que el propio documento original ya apuntaba correctamente: muchas privatizaciones van acompañadas de reformas regulatorias simultáneas —introducción de competencia, liberalización de precios—, lo que dificulta atribuir cualquier mejora observada específicamente al cambio de propietario y no al cambio de marco regulatorio que lo acompañó.

\subsubsection{Estudios de gran muestra: titularidad pública, privada y mixta en coexistencia}
El estudio fundacional de este primer bloque metodológico es el de BOARDMAN y VINING (1989), que analiza una muestra de las 500 mayores empresas industriales no financieras del mundo fuera de Estados Unidos, comparando simultáneamente empresas de titularidad privada, pública y mixta. Su resultado más citado, y el que aporta un matiz que buena parte de la literatura posterior no siempre recoge con la misma precisión, es que tanto las empresas públicas puras como las de titularidad mixta presentan una rentabilidad y una eficiencia sistemáticamente inferiores a las de titularidad privada plena, siendo el resultado de las mixtas, en algunos indicadores, incluso peor que el de las públicas puras: los autores interpretan este hallazgo como evidencia de que la titularidad mixta puede combinar lo peor de ambos mundos: conserva los problemas de agencia política de la titularidad pública sin beneficiarse plenamente de la disciplina de mercado que ejercería una titularidad privada completa. Vining y Boardman retomarían y afinarían esta tesis pocos años después (1992), formulando de manera todavía más explícita la conclusión que atraviesa toda la literatura posterior: es el grado de competencia del entorno en que opera la empresa, mucho más que su titularidad formal, la variable que mejor explica las diferencias de eficiencia observadas -el mismo hallazgo, formulado quince años antes, que la Comisión Europea confirmaría en 2016 con metodología más sofisticada—.

[Seguro] Recupero ahora, con el detalle que merece, el estudio de la Comisión Europea que ya presentaba el documento original —State-Owned Enterprises in the EU: Lessons Learnt and Ways Forward in a Post-Crisis Context (julio de 2016)—, elaborado por su Dirección General de Asuntos Económicos y Financieros sobre los sectores energético y ferroviario, con datos de 2013 en los que las empresas públicas suponían hasta el 60\% del sector energético comunitario. El estudio distingue 14 subsectores —nueve dentro de la energía (generación, transmisión y distribución de electricidad; venta minorista y mayorista de electricidad; transmisión y distribución de gas; venta minorista y mayorista de gas) y tres dentro del ferrocarril (gestión de infraestructuras, transporte de mercancías, transporte de pasajeros), más dos categorías de conglomerados que operan simultáneamente en ambos—, reconociendo explícitamente que en los segmentos regulados —típicamente ligados a la infraestructura de red, con tarifas de acceso reguladas por cada Estado miembro— la rentabilidad depende en parte del diseño tarifario nacional, lo que obliga a controlar por esa variable, mientras que en los segmentos no regulados la intensidad de la competencia efectiva puede también variar de un Estado miembro a otro según la profundidad de sus reformas liberalizadoras. [Seguro] La estimación se articula mediante la regresión Y_it = β0 + β1·SOE_MAJOR_it + β2·SOE_MINOR_it + β3·SIZE_it + α_c + α_t − ε_it, estimada separadamente para cada uno de los catorce subsectores y para cinco variables dependientes distintas —rentabilidad sobre activos (ROA), ratio de gastos operativos sobre ventas (proxy de eficiencia productiva), ratio de costes laborales sobre ventas, ratio de recursos propios sobre activo total, y ratio de inversión—, controlando por país, periodo y tamaño de la empresa, con dos variables dummy de interés que distinguen titularidad pública mayoritaria de titularidad pública minoritaria. [Seguro] Los resultados, ya recogidos en el documento original pero que ahora podemos leer con el marco teórico completo: en los segmentos regulados del sector energético, la eficiencia pública es menor en transmisión de electricidad pero mayor en distribución de gas, sin diferencias significativas en el resto; en los segmentos regulados del ferrocarril, no hay diferencias significativas salvo unos mayores costes laborales privados; en los segmentos no regulados, tanto energéticos como ferroviarios, las empresas públicas son, por lo general, menos rentables y menos eficientes en costes, con mayores costes laborales de forma sistemática y robusta en todos los subsectores analizados —un resultado que los autores atribuyen al mayor poder sindical y a la mayor generosidad salarial de la empresa pública, y a sus menores ajustes de plantilla—; y, dato especialmente relevante para la tesis vertebradora de todo este bloque 3, las empresas de participación pública minoritaria se comportan de forma mucho más parecida a las privadas que las de participación mayoritaria, confirmando empíricamente, con datos europeos de 2013, la misma intuición que Boardman y Vining habían anticipado teóricamente en 1989 sobre el papel decisivo del grado de control público, no solo de su existencia.

\subsubsection{Estudios de antes y después: el efecto causal de la privatización}

[Seguro] El estudio fundacional del segundo bloque metodológico —comparar la misma empresa antes y después de privatizarse— es el de William Megginson, Robert Nash y Matthias van Randenborgh (1994, "The Financial and Operating Performance of Newly Privatized Firms: An International Empirical Analysis", The Journal of Finance), que analiza 61 empresas de 18 países y 32 industrias distintas privatizadas mediante oferta pública entre 1961 y 1990, comparando sus ratios financieros y operativos en los años previos y posteriores a la privatización. [Probable] Sus resultados, que se convertirían en el punto de referencia obligado de toda la literatura posterior, muestran mejoras estadísticamente significativas tras la privatización en rentabilidad, eficiencia operativa (medida como producción por empleado), inversión de capital y nivel de producción, junto con un descenso significativo del apalancamiento financiero —las empresas privatizadas reducen su dependencia de la deuda—, sin que se observe, a diferencia de lo que la crítica sindical anticipaba, un descenso sistemático y significativo del empleo total, aunque sí, en muchos casos, un descenso de la proporción de trabajadores sobre el total de activos.

[Seguro] Recupero ahora, íntegro y con la profundidad ya expuesta en el documento original, el estudio español de Pablo Hernández de Cos (2004, Banco de España), centrado en 33 empresas manufactureras españolas privatizadas entre 1983 y 1996, con datos de la Central de Balances del Banco de España antes y después de cada privatización, eligiendo deliberadamente el sector manufacturero por su carácter ya competitivo tanto antes como después de la privatización, precisamente para aislar el efecto de la titularidad del efecto de un contexto competitivo cambiante. La metodología asume una función de producción Cobb-Douglas, Y_it = A·K_it^α·L_it^β, log-linealizada y estimada mediante la regresión ΔY_t = γ0 + α·k_it + β·l_it + ν_i + ν_t + γ1·PRIVKPU_it + Σφ_m·X_mit + ε, donde PRIVKPU es la variable dummy de interés que distingue empresas ya privatizadas de empresas todavía públicas, y las variables de control —grado de concentración sectorial, penetración de las importaciones, endeudamiento, y proporción de empleo temporal— capturan tanto el grado de competencia del sector, especialmente relevante en un periodo, 1983-1996, marcado por la entrada de España en la CEE y el consiguiente aumento de la competencia exterior, como la estructura financiera y laboral de cada empresa. [Seguro] Los resultados confirman un coeficiente positivo y significativo de la variable de privatización sobre la productividad total de los factores, robusto ante distintas especificaciones alternativas, con efectos adicionales positivos sobre el stock de capital real, la ratio capital/trabajo y la remuneración real por trabajador, y efectos negativos sobre el empleo, el endeudamiento y la cuota de mercado —un patrón, dicho sea de paso, muy coherente con el hallazgo internacional de Megginson, Nash y van Randenborgh una década antes—, si bien el propio autor advierte de las limitaciones derivadas del tamaño reducido de la muestra y de la especificidad del periodo y del contexto regulatorio español analizados.

\subsubsection{La síntesis de toda la evidencia: Megginson y Netter (2001)}

[Seguro] Cierro el repaso empírico con la revisión de literatura de referencia de todo este campo: William Megginson y Jeffry Netter (2001, "From State to Market: A Survey of Empirical Studies on Privatization", Journal of Economic Literature), que sintetizan varias decenas de estudios empíricos publicados durante los años noventa, en economías tanto desarrolladas como en transición y en desarrollo. [Probable] Su conclusión general es la que ya anticipaba, de forma más impresionista, el documento original, pero ahora con el respaldo de una revisión sistemática de toda la literatura disponible hasta la fecha: la evidencia respalda de forma consistente que la privatización mejora la eficiencia operativa y financiera de las empresas privatizadas, especialmente cuando coincide con la introducción de competencia efectiva o con una regulación de calidad en el sector correspondiente, mientras que la evidencia es mucho más débil, e incluso ambigua, cuando la privatización se produce en sectores que permanecen altamente concentrados o débilmente regulados tras el cambio de titularidad —precisamente la misma variable de competencia que Vining y Boardman habían identificado ya en 1992—. Un hallazgo complementario, aportado por Dewenter y Malatesta (2001, "State-Owned and Privately Owned Firms: An Empirical Analysis of Profitability, Leverage, and Labor Intensity", American Economic Review), sobre una muestra de las 500 mayores empresas del mundo en cuatro años distintos entre 1975 y 1995, confirma además que la brecha de rentabilidad entre empresas públicas y privadas es especialmente pronunciada en los años previos inmediatos a la privatización —lo que podría interpretarse, y así lo hacen los propios autores, como evidencia adicional de que los gobiernos privatizan de forma no aleatoria, seleccionando precisamente los años en que la brecha resulta más visible y políticamente urgente de corregir—.

\subsubsection{Conclusiones generales de la aproximación cuantitativa}

[Seguro] Recupero, ya con todo el aparato desarrollado, las conclusiones generales que el documento original enunciaba de forma correcta pero sin el respaldo bibliográfico suficiente. Primera: en sectores sometidos a alto grado de competencia, la gran mayoría de estudios —desde Boardman-Vining hasta la síntesis de Megginson-Netter— son favorables a la mayor eficiencia de la empresa privada, con escasas excepciones, y la propia literatura subraya que crear condiciones competitivas es más determinante para la eficiencia agregada que decidir la titularidad de cada empresa individual. Segunda: en sectores de escasa competencia o sometidos a fuerte regulación —electricidad, ferrocarril, telecomunicaciones en su vertiente de red—, la evidencia empírica, desde el estudio de la Comisión Europea de 2016 hasta la literatura anterior que este mismo estudio cita, no encuentra sistemáticamente una mayor eficiencia de la empresa privada. Tercera: en los países en desarrollo, con mercados de capitales y de competencia menos desarrollados, no puede concluirse con la misma seguridad que las empresas públicas sean menos eficientes que las privadas, aunque la escasez y menor fiabilidad de los datos disponibles obliga a tratar esta conclusión con especial cautela. Y cuarta, la más metodológica de todas: los estudios de tipo "antes-después" tienden sistemáticamente a encontrar mejoras de eficiencia tras la privatización, pero esas mejoras suelen coincidir temporalmente con reformas regulatorias que dificultan aislar con precisión si el efecto causal procede del cambio de titularidad en sí mismo o del nuevo marco competitivo que casi siempre lo acompaña —la advertencia metodológica con la que abríamos este subapartado, y que conviene que mantengas presente como el límite epistemológico real de toda esta literatura, por mucho volumen de estudios que se haya acumulado.

\subsection{La aproximación cualitativa: "las nuevas razones"}
La obra de referencia de esta segunda aproximación es MAZZUCATO (*The Entrepreneurial State: Debunking Public vs. Private Sector Myths*, 2013), y conviene precisar, antes de exponer su tesis, por qué la trato aquí y no en el subapartado cuantitativo: Mazzucato no contrasta una hipótesis mediante regresión sobre una muestra amplia de empresas, sino que construye su argumento mediante **estudios de caso históricos extensos** sobre episodios concretos de innovación tecnológica, un método argumentativo-narrativo que exige un aparato distinto al de la econometría del subapartado 3.1, aunque persiga la misma finalidad última de informar la respuesta a la pregunta vertebradora del tema.

[Seguro] La tesis central de Mazzucato invierte el relato convencional sobre el origen de la innovación tecnológica: sostiene que, lejos de ser el sector privado el motor primario de las revoluciones tecnológicas más importantes de las últimas décadas, ha sido el **Estado**, actuando como inversor de capital riesgo paciente y dispuesto a asumir el riesgo más radical e incierto, quien ha financiado sistemáticamente las tecnologías fundacionales que el sector privado se ha limitado, después, a integrar y comercializar. [Seguro] Su ejemplo más célebre, y el que conviene dominar con precisión porque es, con diferencia, el gancho expositivo más citado de toda esta literatura, es la **deconstrucción del iPhone**: Mazzucato traza el origen de cada una de las tecnologías que lo hacen "inteligente" hasta programas de investigación financiados públicamente y durante décadas —internet, nacido como ARPANET, financiado por el Departamento de Defensa estadounidense a través de DARPA en los años sesenta para resolver problemas de comunicación militar; el GPS, desarrollado por el programa militar NAVSTAR para la localización de equipamiento militar; la pantalla táctil multitáctil, resultado de investigación de posgrado en la Universidad de Delaware financiada por la National Science Foundation y la CIA; Siri, desarrollado originalmente por el Stanford Research Institute bajo encargo de DARPA para crear un asistente virtual destinado a personal militar; el protocolo HTTP, creado por Tim Berners-Lee en el CERN, institución de investigación pública europea; y hasta la batería de ión-litio, desarrollada con financiación del Departamento de Energía estadounidense—. [Probable] La conclusión que Mazzucato extrae de este recorrido no es que Steve Jobs y Apple no aportaran nada de valor —reconoce explícitamente que la integración de todas estas tecnologías dispersas en un producto coherente, deseable y comercialmente exitoso fue, en sí misma, un logro empresarial genuino—, sino que el relato cultural dominante, que atribuye la innovación radical casi en exclusiva al genio emprendedor privado, **invisibiliza sistemáticamente** el papel del Estado como asumidor del riesgo más temprano y más incierto, generando una asimetría en la distribución de las recompensas: el riesgo se socializa —lo asume el conjunto de los contribuyentes, a través de la financiación pública de décadas de investigación sin garantía de retorno—, mientras que la recompensa final se privatiza casi por completo en manos de quienes llegan al final de la cadena de innovación, cuando el riesgo ya se ha disipado.

\subsubsection{El fundamento teórico: riesgo frente a incertidumbre radical}
[Seguro] Conviene dotar de fundamento teórico riguroso a la intuición de Mazzucato, porque de lo contrario corre el riesgo de leerse como una simple anécdota persuasiva y no como un argumento económico con genealogía propia. La distinción que sustenta todo su argumento —por qué el capital privado sistemáticamente infrainvierte en las fases más tempranas y radicales de la innovación— se remonta a **Frank Knight** (*Risk, Uncertainty, and Profit*, 1921), quien distinguió con precisión dos situaciones que el lenguaje cotidiano confunde: el **riesgo**, en el que se desconoce el resultado concreto pero se conoce la distribución de probabilidad de los resultados posibles —de modo que puede calcularse una prima de riesgo y, en principio, asegurarse—, y la **incertidumbre** en sentido estricto, en la que ni siquiera se conoce la distribución de probabilidad, porque el propio abanico de resultados posibles es desconocido de antemano. [Probable] Una tecnología verdaderamente novedosa —internet en los años sesenta, la secuenciación genómica en los años ochenta, la inteligencia artificial generativa hace apenas una década— no presenta riesgo en el sentido actuarial del término, sino **incertidumbre radical** en el sentido de Knight: nadie, ni siquiera el propio inversor, puede calcular una probabilidad de éxito con ningún fundamento estadístico sólido, porque no existe precedente alguno sobre el que construir esa distribución. [Suposición] Bajo esta lectura, la infrainversión privada en investigación verdaderamente radical no es un fallo de mercado en el sentido convencional —una externalidad no internalizada, un monopolio natural—, sino la consecuencia lógica y racional de que ningún capital privado, sujeto a horizontes de rentabilidad relativamente cortos y a inversores que exigen retornos calculables, está dispuesto a asumir una incertidumbre que ni siquiera puede cuantificar: el capital riesgo convencional espera, típicamente, retornos en horizontes de tres a cinco años y asume que solo una de cada diez inversiones tendrá éxito, un cálculo actuarial que presupone precisamente el tipo de conocimiento probabilístico que, por definición, no existe en la fase más temprana de una tecnología genuinamente nueva.

\subsubsection{De la tesis histórica a la propuesta de política: misiones y transición ecológica}
[Seguro] Mazzucato desarrolla, en una obra posterior (*Mission Economy: A Moonshot Guide to Changing Capitalism*, 2021), la implicación normativa de su tesis histórica: si el Estado ha demostrado capacidad histórica para asumir riesgo radical con éxito, el instrumento de política pública adecuado no es limitarse a corregir fallos de mercado puntuales una vez identificados, sino organizar la intervención pública en torno a **misiones** —objetivos ambiciosos, concretos y con plazo definido, siguiendo el modelo del **Programa Apolo** de la NASA en los años sesenta, que fijó el objetivo político de llevar un ser humano a la Luna antes de que terminara la década y organizó, en torno a ese objetivo, la coordinación de miles de empresas, universidades y agencias públicas—, aplicando esa misma lógica organizativa a los grandes retos contemporáneos, muy señaladamente la **transición ecológica**. Bajo esta óptica, la empresa pública —o, más ampliamente, el capital público desplegado a través de bancos de inversión estatales, agencias de innovación o participaciones accionariales directas— no es un instrumento residual reservado a los monopolios naturales que el mercado no puede proveer, sino un actor activo y deliberado en la **creación** de mercados enteramente nuevos, no solo en su corrección.

[Probable] Esta misma lógica de renovación del argumento industrial ha encontrado eco en economistas de corriente principal ajenos a la escuela específica de Mazzucato, lo que conviene señalar para no dar la impresión de que se trata de una posición aislada: **Dani Rodrik**, en su trabajo sobre política industrial verde ("Green Industrial Policy", *Oxford Review of Economic Policy*, 2014, y trabajos posteriores sobre la "normalización" de la política industrial), defiende, desde una tradición de economía del desarrollo mucho más convencional que la de Mazzucato, que la descarbonización de la economía exige una coordinación de inversión a gran escala que los mecanismos de mercado puro, incluidos los impuestos pigouvianos sobre el carbono, no proporcionan por sí solos con la velocidad que la urgencia climática exige, precisamente porque las tecnologías de descarbonización de próxima generación comparten el mismo problema de incertidumbre radical que ya identificamos con Knight.

\subsubsection{La respuesta granular: recuperando el rol de la empresa pública energética}

[Seguro] Recupero aquí, con la profundidad que merece, el contenido sobre el papel de la empresa pública en la transición ecológica que ya recogía el documento original, y que trasladamos deliberadamente desde el bloque 1 a este subapartado, porque encaja mejor como ejemplo concreto y "granular" de estas nuevas razones que como argumento clásico. Algunas de las mayores empresas públicas del mundo operan en el sector energético, y la evidencia disponible muestra que las empresas energéticas de titularidad estatal prestan, en conjunto, mayor atención a consideraciones medioambientales y presentan una mayor probabilidad de fijar objetivos explícitos de reducción de emisiones que sus homólogas privadas —aunque, matiza el propio FMI (2020), esa mayor atención declarada no se traduce en diferencias significativas de consumo energético real ni de ritmo efectivo de descarbonización cuando se compara con el sector privado—. El caso más relevante, por su magnitud e impacto, es el del **carbón**: en torno al 60% de las minas y centrales eléctricas de carbón del mundo son de titularidad pública, y ese dato convierte a la empresa pública en actor decisivo, casi por definición, de cualquier proceso serio de reconversión del sector. La propia experiencia europea ofrece ejemplos de gestión de esa reconversión que conviene retener como ganchos expositivos: la conversión de minas en centrales petroquímicas en los Países Bajos durante los años setenta, o la creación deliberada de empleo público alternativo en las regiones alemanas afectadas por el cierre de la minería del carbón, en ambos casos como respuesta a la relevancia social —muy concentrada territorialmente, hasta representar entre el 10% y el 20% del empleo local en regiones concretas de Bulgaria, Grecia o Polonia— que esta industria conserva en determinadas comunidades, y que cualquier estrategia de descarbonización debe internalizar si quiere evitar el mismo tipo de conflicto social que ya analizamos, en otro contexto, al hablar de las condiciones laborales de cualquier proceso de reestructuración en el subapartado 2.4.2.

\subsubsection{Una lectura crítica, no acrítica, de Mazzucato}

[Probable] Cierro este subapartado, y con él el aparato teórico central del tema, con el mismo ejercicio de equilibrio que ya aplicamos al presentar a LE GRAND: la tesis de MAZZUCATO, pese a su enorme influencia en el debate de política económica de la última década, no es unánimemente aceptada, y conviene presentar también su crítica. 

[Suposición] La objeción más citada, procedente de economistas de tradición más próxima al mercado, es de naturaleza contrafactual: que una tecnología haya sido financiada originalmente por el Estado no demuestra que el sector privado no hubiera llegado, eventualmente, al mismo resultado por una vía distinta (quizá más lenta, quizá con un diseño institucional diferente), y que el ejemplo de DARPA en particular puede no ser generalizable, porque DARPA opera con un grado de autonomía respecto de la burocracia estatal ordinaria, con presupuestos plurianuales protegidos y con gestores de programa dotados de una discrecionalidad que dista mucho del funcionario medio descrito por NISKANEN. 

[Probable] Esta crítica, bien entendida, no invalida el argumento de fondo de Mazzucato sobre la incertidumbre radical, pero sí matiza sustancialmente su implicación de política: sugiere que la clave no es la titularidad pública en abstracto, sino un diseño institucional específico —autonomía de gestión, protección frente a la interferencia política de corto plazo, tolerancia organizativa al fracaso— que la literatura de contratos incompletos ya había anticipado en términos más generales: la provisión pública resulta preferible precisamente cuando la calidad (o el objeto de la innovación) es difícil de especificar contractualmente de antemano, que es exactamente la situación de la investigación en condiciones de incertidumbre knightiana. 

**Las dos aproximaciones, la teórica de contratos incompletos del bloque 2 y la histórico-narrativa de Mazzucato en este bloque 3, convergen, por caminos completamente distintos, en la misma conclusión**: no es la titularidad por sí sola lo que determina el resultado, sino el ajuste entre el tipo de problema —contratable o no, cierto o radicalmente incierto— y el diseño institucional elegido para resolverlo.



### Cierre del tema

[Probable] Llegados aquí, podemos por fin responder con precisión a la pregunta que enunciamos, deliberadamente, antes incluso de empezar la exposición: **¿es la titularidad pública o privada de una empresa lo que determina su eficiencia, o lo es el grado de competencia y la calidad regulatoria del sector en que opera?** La respuesta que este recorrido dialéctico —tesis histórica y clásica, antítesis teórica y empírica, síntesis cuantitativa y cualitativa— permite dar no es "depende", en el sentido vago y evasivo del término, sino una respuesta **condicionada y verificable**: en sectores de competencia efectiva, la evidencia acumulada durante más de tres décadas respalda de forma consistente la mayor eficiencia de la titularidad privada; en sectores de monopolio natural o fuerte regulación, la titularidad deja de ser la variable decisiva, y lo que importa es la calidad del diseño institucional, sea cual sea el propietario formal; y en el terreno de la innovación bajo incertidumbre radical, la evidencia histórica sugiere que el Estado ha jugado, y probablemente deba seguir jugando, un papel de inversor de primer riesgo que ningún argumento de agencia o elección pública, por sólido que sea en su ámbito propio, ha conseguido sustituir satisfactoriamente por una alternativa puramente privada.



\clearpage
\section{Conclusión} 


\subsection{Recapitulación}


\subsection{Extensiones y relación con otras partes del Temario}



\subsection{Opinión}

\subsubsection{Idea final (Salida o cierra)}

\clearpage
\pagenumbering{gobble}
\MyBibliography


CROSLAND, C. A. R. (1956): The Future of Socialism, Jonathan Cape, Londres.
VERNON, R. (1971): Sovereignty at Bay: The Multinational Spread of U.S. Enterprises, Basic Books, Nueva York.
O'CONNOR, J. (1973): The Fiscal Crisis of the State, St. Martin's Press, Nueva York.
Historia institucional del IRI (Istituto per la Ricostruzione Industriale), Italia, 1933.
Historia institucional del INI (Instituto Nacional de Industria), España, Ley de 25 de septiembre de 1941.
Decreto de expropiación petrolera de México, 18 de marzo de 1938 (Lázaro Cárdenas).
- MUSGRAVE, R. A. (1959): *The Theory of Public Finance: A Study in Public Economy*, McGraw-Hill, Nueva York.
- BAUMOL, W. J., PANZAR, J. C. y WILLIG, R. D. (1982): *Contestable Markets and the Theory of Industry Structure*, Harcourt Brace Jovanovich, Nueva York.
- KAHN, A. E. (1970-1971): *The Economics of Regulation: Principles and Institutions*, John Wiley & Sons, Nueva York.
- KHAN, M. S. y KUMAR, M. S. (1997): estudio sobre infraestructuras públicas en países en desarrollo (ya referenciado en el documento original del tema).
- FMI (2020): *Fiscal Monitor: Policies to Support People During the COVID-19 Pandemic*.
- AHARONI, Y. (1986): *The Evolution and Management of State Owned Enterprises*, Ballinger Publishing, Cambridge (Massachusetts).
- Reglamento (UE) n.º 549/2013, relativo al Sistema Europeo de Cuentas Nacionales y Regionales de la Unión Europea (SEC 2010).
- Sentencias del Tribunal de Justicia de 4 de junio de 2002, asuntos C-367/98 (Comisión/Portugal), C-483/99 (Comisión/Francia) y C-503/99 (Comisión/Bélgica).
- Sentencias del Tribunal de Justicia de 13 de mayo de 2003, asuntos C-463/00 (Comisión/España) y C-98/01 (Comisión/Reino Unido).
- Sentencia del Tribunal de Justicia de 23 de octubre de 2007, asunto C-112/05 (Comisión/Alemania, "Ley Volkswagen").
- BUCHANAN, J. M. y TULLOCK, G. (1962): *The Calculus of Consent: Logical Foundations of Constitutional Democracy*, University of Michigan Press, Ann Arbor.
- FRIEDMAN, M. (1962): *Capitalism and Freedom*, University of Chicago Press, Chicago.
- TULLOCK, G. (1967): "The Welfare Costs of Tariffs, Monopolies, and Theft", *Western Economic Journal*, vol. 5, núm. 3.
- DOWNS, A. (1957): *An Economic Theory of Democracy*, Harper & Row, Nueva York.
- LE GRAND, J. (1997): "Knights, Knaves or Pawns? Human Behaviour and Social Policy", *Journal of Social Policy*, vol. 26, núm. 2.
- VICKERS, J. y YARROW, G. (1988): *Privatization: An Economic Analysis*, MIT Press, Cambridge (Massachusetts).
- NISKANEN, W. A. (1971): *Bureaucracy and Representative Government*, Aldine-Atherton, Chicago; y (1975): "Bureaucrats and Politicians", *Journal of Law and Economics*, vol. 18, núm. 3.
- MIGUÉ, J.-L. y BÉLANGER, G. (1974): "Toward a General Theory of Managerial Discretion", *Public Choice*, vol. 17.
- JENSEN, M. C. y MECKLING, W. H. (1976): "Theory of the Firm: Managerial Behavior, Agency Costs, and Ownership Structure", *Journal of Financial Economics*, vol. 3, núm. 4.
- HIRSCHMAN, A. O. (1970): *Exit, Voice, and Loyalty: Responses to Decline in Firms, Organizations, and States*, Harvard University Press, Cambridge (Massachusetts).
- MANNE, H. G. (1965): "Mergers and the Market for Corporate Control", *Journal of Political Economy*, vol. 73, núm. 2.
- KORNAI, J. (1980): *Economics of Shortage*, North-Holland, Ámsterdam.
- KORNAI, J., MASKIN, E. y ROLAND, G. (2003): "Understanding the Soft Budget Constraint", *Journal of Economic Literature*, vol. 41, núm. 4.
- GROSSMAN, S. J. y HART, O. D. (1986): "The Costs and Benefits of Ownership: A Theory of Vertical and Lateral Integration", *Journal of Political Economy*, vol. 94, núm. 4.
- HART, O., SHLEIFER, A. y VISHNY, R. W. (1997): "The Proper Scope of Government: Theory and an Application to Prisons", *Quarterly Journal of Economics*, vol. 112, núm. 4.
- SHLEIFER, A. (1998): "State versus Private Ownership", *Journal of Economic Perspectives*, vol. 12, núm. 4.
- OXERA (2013): "The Thatcher Privatisation Legacy", *Oxera Agenda*, mayo de 2013.
- BOYCKO, M., SHLEIFER, A. y VISHNY, R. W. (1996): "A Theory of Privatisation", *The Economic Journal*, vol. 106, núm. 435.
- WILLIAMSON, J. (1989): "What Washington Means by Policy Reform", en *Latin American Adjustment: How Much Has Happened?*, Institute for International Economics, Washington D.C.
- DJANKOV, S. y MURRELL, P. (2002): "Enterprise Restructuring in Transition: A Quantitative Survey", *Journal of Economic Literature*, vol. 40, núm. 3 (ya referenciado en el documento original).
- AKERLOF, G. A. (1970): "The Market for 'Lemons': Quality Uncertainty and the Market Mechanism", *Quarterly Journal of Economics*, vol. 84, núm. 3.
- CHONG, A. y LÓPEZ-DE-SILANES, F. (2004): *Privatization in Latin America: Myths and Reality*, Banco Mundial / Stanford University Press.
- PEROTTI, E. C. (1995): "Credible Privatization", *American Economic Review*, vol. 85, núm. 4.
- OCDE (2009): *Privatisation in the 21st Century: Recent Experiences of OECD Countries*.
- ESTACHE, A., DE SCHMITT, A. J. A. y SYDENSTRICKER, E. (2000): estudio sobre la privatización de Brazil Railways (ya referenciado en el documento original).
BOARDMAN, A. E. y VINING, A. R. (1989): "Ownership and Performance in Competitive Environments: A Comparison of the Performance of Private, Mixed, and State-Owned Enterprises", Journal of Law and Economics, vol. 32, núm. 1.
VINING, A. R. y BOARDMAN, A. E. (1992): "Ownership versus Competition: Efficiency in Public Enterprise", Public Choice, vol. 73, núm. 2.
Comisión Europea (2016): State-Owned Enterprises in the EU: Lessons Learnt and Ways Forward in a Post-Crisis Context, Directorate-General for Economic and Financial Affairs (ya referenciado en el documento original, ahora en su contenido metodológico completo).
MEGGINSON, W. L., NASH, R. C. y VAN RANDENBORGH, M. (1994): "The Financial and Operating Performance of Newly Privatized Firms: An International Empirical Analysis", The Journal of Finance, vol. 49, núm. 2.
HERNÁNDEZ DE COS, P. (2004): estudio del Banco de España sobre privatización y productividad en la industria manufacturera española (ya referenciado en el documento original, ahora en su contenido metodológico completo).
MEGGINSON, W. L. y NETTER, J. M. (2001): "From State to Market: A Survey of Empirical Studies on Privatization", Journal of Economic Literature, vol. 39, núm. 2.
DEWENTER, K. L. y MALATESTA, P. H. (2001): "State-Owned and Privately Owned Firms: An Empirical Analysis of Profitability, Leverage, and Labor Intensity", American Economic Review, vol. 91, núm. 1.
- MAZZUCATO, M. (2013): *The Entrepreneurial State: Debunking Public vs. Private Sector Myths*, Anthem Press, Londres.
- MAZZUCATO, M. (2021): *Mission Economy: A Moonshot Guide to Changing Capitalism*, Allen Lane, Londres.
- KNIGHT, F. H. (1921): *Risk, Uncertainty, and Profit*, Houghton Mifflin, Boston.
- RODRIK, D. (2014): "Green Industrial Policy", *Oxford Review of Economic Policy*, vol. 30, núm. 3.
- FMI (2020): *Fiscal Monitor: Policies to Support People During the COVID-19 Pandemic* (ya referenciado en el documento original, ahora en su contenido específico sobre empresas públicas energéticas y descarbonización).

% ANEXOS
\clearpage










\end{document}